% Acides et bases en solution aqueuse : eau, autoprotolyse et pH — Chimie Tle D — Cours signé OnBuch+
% Programme officiel MINESEC, Terminale C, D, E. Compiler avec : tectonic L06-eau-autoprotolyse-et-ph.tex
\newcommand{\DOCMATIERE}{Chimie}
\newcommand{\DOCNIVEAU}{Tle D}
\newcommand{\DOCMODULE}{Module 2 · Acides et bases}
\newcommand{\DOCLECON}{06}
\newcommand{\DOCTITRE}{Acides et bases en solution aqueuse : eau, autoprotolyse et pH}
% OnBuch+ — préambule « cours signés » ENRICHI (compilé avec tectonic / XeLaTeX)
% Système visuel « Pop » d'OnBuch+ : fond crème (jamais blanc), bordures encre
% épaisses, ombres « dures » décalées, titres Archivo Black en capitales.
% Version enrichie pour les cours de chimie : chimie (mhchem, chemfig),
% graphes (pgfplots), boîtes pédagogiques supplémentaires (définition,
% propriété, méthode, attention, le savais-tu, exercice résolu, exercice,
% TP, bilan), page de garde refaite et signature OnBuch+ ancrée en bas.
%
% Macros attendues dans main.tex AVANT \input{preamble} :
%   \DOCMATIERE  \DOCNIVEAU  \DOCMODULE  \DOCLECON (numéro)  \DOCTITRE
\documentclass[11pt]{article}

\usepackage{fontspec}
\usepackage[french]{babel}
\usepackage[a4paper,margin=2cm,top=3.4cm,bottom=2.8cm,headheight=1.4cm,footskip=1.3cm]{geometry}
\usepackage{amsmath,amssymb,amsfonts}
\usepackage[table]{xcolor} % [table] : fournit aussi \rowcolors (alternance de lignes)
\usepackage{pagecolor}
\usepackage{tikz}
\usetikzlibrary{arrows.meta,positioning,calc,shapes.geometric,shapes.misc,shapes.arrows,decorations.pathmorphing,decorations.markings,patterns,shadows,mindmap,trees,fit,backgrounds,matrix,intersections}
\usepackage{pgfplots}
\pgfplotsset{compat=1.18}
\usepackage[version=4]{mhchem}
\usepackage{chemfig}
\usepackage{enumitem}
\usepackage{fancyhdr}
% [most] charge toutes les bibliothèques tcolorbox (listings, minted, raster,
% documentation...) et gonfle inutilement la mémoire TeX sur un document long
% (~300 boîtes) : on ne charge que ce qui est réellement utilisé.
\usepackage[skins,breakable]{tcolorbox}
\usepackage{graphicx}
\usepackage{colortbl}
\usepackage{booktabs}
\usepackage{tabularx}
\usepackage{array}
\usepackage{multicol}
\usepackage{siunitx}
% parse-numbers=false : accepte aussi \SI{2,5 \times 10^{-3}}{...}
\sisetup{locale=FR,output-decimal-marker={,},group-minimum-digits=4,parse-numbers=false}
\usepackage{qrcode}
% \degree : commande fréquemment utilisée par les modèles pour un angle ou une
% température en dehors de \SI{}{} (non fournie par les paquets chargés).
\providecommand{\degree}{^{\circ}}
\usepackage{lastpage}
\usepackage{hyperref}
\hypersetup{hidelinks}

% ============================================================
% Palette « Pop » OnBuch+ — mêmes hex que la classe P (plus_ui.dart)
% ============================================================
\definecolor{poporange}{HTML}{F59321}
\definecolor{popdark}{HTML}{E07A0C}
\definecolor{popcream}{HTML}{FFF6E8}
\definecolor{popink}{HTML}{1C1714}
\definecolor{poppurple}{HTML}{7A5AE0}
\definecolor{popgreen}{HTML}{2E9E6A}
\definecolor{poppink}{HTML}{E0457B}
\definecolor{popblue}{HTML}{2F7DE1}
\definecolor{popgold}{HTML}{C98A12}
\definecolor{popmuted}{HTML}{8A7F76}
\definecolor{popline}{HTML}{EADFD2}
\definecolor{popgray}{HTML}{8A7F76}
\colorlet{popgrey}{popgray}
% Alias tolérants (le modèle nomme parfois les couleurs différemment)
\colorlet{popwhite}{white}
\colorlet{popblack}{popink}
\colorlet{popred}{poppink}
\colorlet{popyellow}{popgold}
% Teintes claires pour fonds de tableaux / zones
\colorlet{popblueL}{popblue!10!white}
\colorlet{poporangeL}{poporange!14!white}
\colorlet{popgreenL}{popgreen!12!white}
\colorlet{poppurpleL}{poppurple!12!white}
\colorlet{poppinkL}{poppink!12!white}
\colorlet{popgoldL}{popgold!14!white}

\pagecolor{popcream}

% ============================================================
% Typographie
% ============================================================
\defaultfontfeatures{Ligatures=TeX}
\setmainfont{PlusJakartaSans-Medium.ttf}[Path=../../../fonts/,BoldFont=PlusJakartaSans-Bold.ttf]
% Maths en sans-serif (Fira Math) avec chiffres et lettres droites de Plus
% Jakarta, pour que les formules se fondent dans le texte.
\usepackage[math-style=ISO,bold-style=ISO]{unicode-math}
\setmathfont{FiraMath-Regular.otf}[Path=../../../fonts/]
\setmathfont{PlusJakartaSans-Medium.ttf}[Path=../../../fonts/,range={up/num,up/{Latin,latin},\mathup}]
\usepackage{icomma} % virgule décimale sans espace en mode math
% Caractères absents de la police de texte (sinon carré vide) : rendus par
% leur équivalent mathématique, en texte comme en formule.
\usepackage{newunicodechar}
\newunicodechar{α}{\ensuremath{\alpha}}
\newunicodechar{Α}{\ensuremath{\alpha}}
\newunicodechar{β}{\ensuremath{\beta}}
\newunicodechar{δ}{\ensuremath{\delta}}
\newunicodechar{✓}{\popcheck{popgreen}}
\newfontfamily\popdisplay{ArchivoBlack-Regular.ttf}[Path=../../../fonts/]
\newfontfamily\poplogo{SpaceGrotesk-Bold.ttf}[Path=../../../fonts/]
\newfontfamily\popsemi{PlusJakartaSans-SemiBold.ttf}[Path=../../../fonts/]
\newfontfamily\popxbold{PlusJakartaSans-ExtraBold.ttf}[Path=../../../fonts/]
\newfontfamily\popxboldls{PlusJakartaSans-ExtraBold.ttf}[Path=../../../fonts/,LetterSpace=3.0]

\setlist[enumerate]{leftmargin=1.7em,itemsep=5pt,topsep=4pt}
\setlist[itemize]{leftmargin=1.7em,itemsep=4pt,topsep=4pt,label={\color{poporange}\textbullet}}
\setlength{\parindent}{0pt}
\setlength{\parskip}{5pt}
\renewcommand{\arraystretch}{1.25}

% chemfig : liaisons un peu plus épaisses, encre Pop
\setchemfig{atom sep=2em,bond style={line width=0.9pt,popink},cram width=3pt}

% pgfplots : style Pop par défaut
\pgfplotsset{
  every axis/.append style={axis line style={popink,line width=1pt},tick style={popink},
    grid style={popline,line width=0.5pt},label style={font=\small},tick label style={font=\footnotesize}},
  popcurve/.style={poporange,line width=1.8pt,no markers,smooth},
}

% ============================================================
% Pill / squiggle
% ============================================================
\newcommand{\poppill}[3]{%
  \tikz[baseline=-0.55ex]\node[draw=popink,line width=1.2pt,rounded corners=2.6pt,%
    fill=#2,inner xsep=6.5pt,inner ysep=3pt]{%
    \popxboldls\fontsize{7}{7}\selectfont\color{#3}\MakeUppercase{#1}};%
}
\newcommand{\popsquiggle}[1]{%
  \tikz[baseline=-0.4ex]{
    \draw[#1,line width=2.6pt,line cap=round]
      plot [smooth] coordinates {(0,0.09) (0.34,0.01) (0.68,0.21) (1.02,0.01) (1.36,0.10)};
  }%
}
% Mot-clé surligné (marqueur orange sous le texte)
\newcommand{\cle}[1]{{\popxbold\color{popdark}#1}}

% ============================================================
% En-tête / pied
% ============================================================
\pagestyle{fancy}
\fancyhf{}
\renewcommand{\headrulewidth}{0pt}
\renewcommand{\footrulewidth}{0pt}
\lhead{\raisebox{-0.15ex}{\poplogo\fontsize{13}{13}\selectfont\color{popink}OnBuch\color{poporange}+}}
\rhead{\poppill{\DOCMATIERE\ \textbullet\ \DOCNIVEAU\ \textbullet\ Leçon \DOCLECON}{white}{popink}}
\lfoot{\popxbold\fontsize{7.6}{7.6}\selectfont\color{popmuted}\MakeUppercase{Cours signé OnBuch+}\ \textbullet\ {\popsemi\DOCTITRE}}
\rfoot{\tikz[baseline=-0.6ex]\node[rounded corners=8.5pt,fill=popink,minimum height=17pt,minimum width=17pt,inner xsep=5pt,inner sep=0]{\popxbold\fontsize{8}{8}\selectfont\color{popcream}\thepage\,/\,\pageref*{LastPage}};}

% ============================================================
% Titres
% ============================================================
\newcommand{\chaptitre}[2]{%
  \begin{tcolorbox}[enhanced,colback=poporange,colframe=popink,boxrule=2.6pt,arc=11pt,
      drop shadow=popink,left=15pt,right=15pt,top=12pt,bottom=11pt,before skip=3pt,after skip=18pt]
    \poppill{#2}{popink}{popcream}\par\vspace{9pt}
    {\raggedright\hyphenpenalty=10000\exhyphenpenalty=10000\popdisplay\fontsize{22}{24}\selectfont\color{popcream}\MakeUppercase{#1}\par}
  \end{tcolorbox}
}
% Section numérotée : pastille orange avec le numéro + titre Archivo
\newcounter{popsec}
\newcommand{\coursec}[1]{%
  \stepcounter{popsec}\setcounter{popsub}{0}%
  \par\vspace{16pt}\noindent
  \begin{minipage}{\linewidth}
  \tikz[baseline=-0.7ex]\node[draw=popink,line width=1.6pt,fill=poporange,rounded corners=4pt,minimum size=22pt,inner sep=0]{\popdisplay\fontsize{12}{12}\selectfont\color{popcream}\thepopsec};%
  \hspace{8pt}{\popdisplay\fontsize{15}{17}\selectfont\color{popink}\MakeUppercase{#1}}\par
  \vspace{4pt}\noindent{\color{poporange}\rule{36pt}{3.6pt}}
  \end{minipage}\par\nopagebreak\vspace{8pt}
}
\newcounter{popsub}
\newcommand{\courssub}[1]{%
  \stepcounter{popsub}%
  \par\vspace{9pt}\noindent{\popxbold\fontsize{11.5}{13}\selectfont\color{popink}{\color{poporange}\thepopsec.\thepopsub}\hspace{6pt}#1}\par\nopagebreak\vspace{4pt}
}

% ============================================================
% Boîtes pédagogiques — même recette : fond blanc, bordure épaisse colorée,
% ombre dure encre, badge plein. Titre optionnel [#1].
% ============================================================
\tcbset{popbase/.style={breakable,enhanced,colback=white,boxrule=2.3pt,arc=9pt,
  shadow={3pt}{-3pt}{0pt}{popink},left=14pt,right=14pt,top=11pt,bottom=11pt,before skip=11pt,after skip=15pt,
  fontupper=\popsemi\fontsize{10.6}{14.5}\selectfont\color{popink},
  fontlower=\popsemi\fontsize{10.6}{14.5}\selectfont\color{popink}}}
% Coche dessinée (le glyphe ✓ n'existe pas dans les polices du document)
\newcommand{\popcheck}[1]{\tikz[baseline=-0.5ex]\draw[#1,line width=1.6pt,line cap=round,line join=round](-0.13,0.0)--(-0.04,-0.09)--(0.13,0.1);}
\newcommand{\popboxhead}[3]{\poppill{#1}{#2}{white}\ifx\relax#3\relax\else\hspace{6pt}{\popxbold\fontsize{10.6}{12}\selectfont\color{popink}#3}\fi\par\vspace{7pt}}

\newtcolorbox{aretenir}[1][]{popbase,colframe=popgreen,before upper={\popboxhead{À retenir}{popgreen}{#1}}}
\newtcolorbox{exemplebox}[1][]{popbase,colframe=poporange,before upper={\popboxhead{Exemple}{poporange}{#1}}}
\newtcolorbox{definition}[1][]{popbase,colframe=popblue,before upper={\popboxhead{Définition}{popblue}{#1}}}
\newtcolorbox{propriete}[1][]{popbase,colframe=poppurple,before upper={\popboxhead{Propriété}{poppurple}{#1}}}
\newtcolorbox{methode}[1][]{popbase,colframe=popgold,before upper={\popboxhead{Méthode}{popgold}{#1}}}
\newtcolorbox{attention}[1][]{popbase,colframe=poppink,before upper={\popboxhead{Attention}{poppink}{#1}}}
\newtcolorbox{experience}[1][]{popbase,colframe=popblue,colback=popblueL,before upper={\popboxhead{Expérience}{popblue}{#1}}}
% « Le savais-tu ? » : carte violette pleine (contexte, culture, Cameroun)
\newtcolorbox{savaistu}[1][]{popbase,colback=poppurple,colframe=popink,
  fontupper=\popsemi\fontsize{10.4}{14.2}\selectfont\color{white},
  before upper={\poppill{Le savais-tu\,?}{white}{poppurple}\ifx\relax#1\relax\else\hspace{6pt}{\popxbold\color{white}#1}\fi\par\vspace{7pt}}}
% Exercice résolu : énoncé en haut, \tcblower puis la solution sur fond crème
\newcounter{popexo}\newcounter{popexr}
\newtcolorbox{exoresolu}[1][]{popbase,colframe=poporange,colbacklower=popcream,
  segmentation style={popink,line width=1.2pt,dashed},
  before upper={\stepcounter{popexr}\popboxhead{Exercice résolu \thepopexr}{poporange}{#1}},
  before lower={\poppill{Solution}{popink}{popcream}\par\vspace{6pt}}}
% Exercice (non résolu en place) — la correction est dans la partie Corrigés
\newtcolorbox{exercice}[1][]{popbase,colframe=popink,
  before upper={\stepcounter{popexo}\popboxhead{Exercice \thepopexo}{popink}{#1}}}
% Corrigé
\newtcolorbox{corrige}[1][]{popbase,colframe=popgreen,colback=popgreenL,
  before upper={\popboxhead{Corrigé}{popgreen}{#1}}}
% Objectifs / prérequis / bilan
\newtcolorbox{objectifs}{popbase,colframe=popink,colback=poporangeL,
  before upper={\popboxhead{Objectifs de la leçon}{popink}{}}}
\newtcolorbox{prerequis}{popbase,colframe=popink,colback=popblueL,
  before upper={\popboxhead{Prérequis}{popblue}{}}}
\newtcolorbox{bilan}{popbase,colframe=popink,colback=white,boxrule=2.8pt,
  before upper={\popboxhead{Fiche bilan}{poporange}{L'essentiel à connaître pour l'examen}}}
% Figure encadrée avec légende
\newtcolorbox{popfigure}[1][]{enhanced,breakable=false,colback=white,colframe=popline,boxrule=1.4pt,arc=8pt,
  left=10pt,right=10pt,top=10pt,bottom=8pt,before skip=10pt,after skip=12pt,halign=center,
  lower separated=false,halign lower=center,
  fontlower=\popsemi\fontsize{9}{11.5}\selectfont\color{popmuted},
  #1}
\newcommand{\legende}[1]{\par\vspace{6pt}{\popsemi\fontsize{9}{11.5}\selectfont\color{popmuted}#1}}

% ============================================================
% Signature OnBuch+
% ============================================================
\newcommand{\poplogoinline}[1]{{\poplogo\fontsize{#1}{#1}\selectfont\color{popink}OnBuch\color{poporange}+}}
\newcommand{\signatureonbuch}{%
  \begin{tcolorbox}[enhanced,colback=popink,colframe=popink,boxrule=2.2pt,arc=10pt,
      drop shadow=poporange,left=14pt,right=14pt,top=10pt,bottom=10pt,before skip=0pt,after skip=0pt]
    \tikz[baseline=-0.7ex]\node[circle,fill=popgreen,draw=popcream,line width=1.3pt,minimum size=22pt,inner sep=0]{\popcheck{white}};%
    \hspace{9pt}\begin{minipage}[c]{0.62\linewidth}
      {\popxbold\fontsize{12}{13}\selectfont\color{popcream}Signé\ \textbullet\ {\poplogo OnBuch\color{poporange}+}}\par\vspace{2pt}
      {\raggedright\popsemi\fontsize{8}{10.5}\selectfont\color{popline}\MakeUppercase{Équipe pédagogique OnBuch+}\\\MakeUppercase{Conforme au programme officiel MINESEC}\par}
    \end{minipage}\hfill
    {\poplogo\fontsize{22}{22}\selectfont\color{popcream}OnBuch\color{poporange}+}
  \end{tcolorbox}%
}

% ============================================================
% Page de garde
% #1 titre · #2 matière · #3 niveau · #4 module · #5 n° leçon · #6 durée ·
% #7 description / objectifs (texte ou itemize)
% ============================================================
\newcommand{\pagegarde}[7]{%
  \thispagestyle{empty}
  \newgeometry{a4paper,margin=1.7cm,top=1.6cm,bottom=1.4cm}%
  \begin{tikzpicture}[remember picture,overlay]
    \fill[poporange!18] ([xshift=-3.2cm,yshift=-3.6cm]current page.north east) circle (4.6cm);
    \fill[poppurple!14] ([xshift=2.4cm,yshift=7.6cm]current page.south west) circle (3.8cm);
  \end{tikzpicture}%
  {\poplogo\fontsize{30}{30}\selectfont\color{popink}OnBuch\color{poporange}+}\hfill
  \raisebox{0.6ex}{\poppill{Cours signé \textbullet\ Premium}{popink}{popcream}}\par
  \vspace{1pt}\popsquiggle{poppurple}\par
  \vspace{8pt}
  \begin{tcolorbox}[enhanced,colback=poporange,colframe=popink,boxrule=2.8pt,arc=13pt,
      drop shadow=popink,left=18pt,right=18pt,top=12pt,bottom=13pt,after skip=10pt]
    \poppill{#2 \textbullet\ #3}{popink}{popcream}\hspace{5pt}\poppill{#4}{white}{popink}\par\vspace{8pt}
    {\popdisplay\fontsize{14}{14}\selectfont\color{popink}LEÇON #5}\par\vspace{4pt}
    {\raggedright\hyphenpenalty=10000\exhyphenpenalty=10000\popdisplay\fontsize{25}{27}\selectfont\color{popcream}\MakeUppercase{#1}\par}
    \vspace{8pt}
    {\popxbold\fontsize{9.5}{9.5}\selectfont\color{popink}\MakeUppercase{Durée conseillée : #6}}
  \end{tcolorbox}
  \begin{tcolorbox}[enhanced,colback=white,colframe=popink,boxrule=2.2pt,arc=10pt,
      drop shadow=popink,left=16pt,right=16pt,top=10pt,bottom=10pt,after skip=8pt]
    \poppill{Dans cette leçon}{poppurple}{white}\par\vspace{5pt}
    \popsemi\fontsize{9.3}{12.6}\selectfont\color{popink}#7
  \end{tcolorbox}
  \noindent\begin{minipage}[t]{0.487\linewidth}
    \begin{tcolorbox}[enhanced,colback=popgreen,colframe=popink,boxrule=2.2pt,arc=10pt,
        drop shadow=popink,left=11pt,right=11pt,top=10pt,bottom=10pt,height=3.1cm,valign=center]
      \tcbox[colback=white,colframe=popink,boxrule=1.3pt,arc=5pt,boxsep=0pt,nobeforeafter,
        left=3pt,right=3pt,top=3pt,bottom=3pt,valign=center]{\qrcode[height=1.9cm]{https://play.google.com/store/apps/details?id=com.luvvix.onbuch&hl=fr}}%
      \hspace{8pt}\begin{minipage}[c]{0.5\linewidth}
        {\popdisplay\fontsize{11}{12.5}\selectfont\color{white}\MakeUppercase{Télécharge OnBuch}}\par\vspace{4pt}
        {\raggedright\popsemi\fontsize{8.4}{11}\selectfont\color{white}Gratuit \textbullet\ Google Play\\Tuteur IA, annales, concours, résultats\par}
      \end{minipage}
    \end{tcolorbox}
  \end{minipage}\hfill
  \begin{minipage}[t]{0.487\linewidth}
    \begin{tcolorbox}[enhanced,colback=poppurple,colframe=popink,boxrule=2.2pt,arc=10pt,
        drop shadow=popink,left=11pt,right=11pt,top=10pt,bottom=10pt,height=3.1cm,valign=center]
      \tikz[baseline=-0.7ex]\node[circle,fill=white,draw=popink,line width=1.3pt,minimum size=1.6cm,inner sep=0]{\popdisplay\fontsize{24}{24}\selectfont\color{poppurple}+};%
      \hspace{8pt}\begin{minipage}[c]{0.6\linewidth}
        {\popdisplay\fontsize{11}{12.5}\selectfont\color{white}\MakeUppercase{Passe à OnBuch+}}\par\vspace{4pt}
        {\raggedright\popsemi\fontsize{8.4}{11}\selectfont\color{white}Tous les cours signés, Tuteur IA illimité, fiches PDF — onglet OnBuch+ de l'app\par}
      \end{minipage}
    \end{tcolorbox}
  \end{minipage}
  \vspace{8pt}
  \popcontenu
  \vfill
  \signatureonbuch
  \restoregeometry
  \newpage
}

% Rangée « Ce cours contient » (page de garde)
\newcommand{\poptile}[3]{% #1 couleur #2 gros texte #3 légende
  \begin{tcolorbox}[enhanced,colback=#1,colframe=popink,boxrule=2pt,arc=9pt,shadow={2.5pt}{-2.5pt}{0pt}{popink},
      left=6pt,right=6pt,top=9pt,bottom=9pt,halign=center,valign=center,height=1.75cm,width=\linewidth,nobeforeafter]
    {\popdisplay\fontsize{12.5}{13}\selectfont\color{white}\MakeUppercase{#2}}\par\vspace{3pt}
    {\centering\popsemi\fontsize{7.8}{9.5}\selectfont\color{white}#3\par}
  \end{tcolorbox}}
\newcommand{\popcontenu}{%
  \par\noindent{\popxboldls\fontsize{7.5}{7.5}\selectfont\color{popmuted}CE COURS SIGNÉ CONTIENT}\par\vspace{7pt}
  \noindent\begin{minipage}[t]{0.235\linewidth}\poptile{popblue}{Cours}{complet, illustré, pas à pas}\end{minipage}\hfill
  \begin{minipage}[t]{0.235\linewidth}\poptile{poporange}{Méthodes}{et exercices résolus}\end{minipage}\hfill
  \begin{minipage}[t]{0.235\linewidth}\poptile{poppink}{Exercices}{type examen + corrigés}\end{minipage}\hfill
  \begin{minipage}[t]{0.235\linewidth}\poptile{popgreen}{Fiche bilan}{l'essentiel à retenir}\end{minipage}\par}

% Sommaire
\newtcolorbox{sommaire}{enhanced,colback=white,colframe=popink,boxrule=2.2pt,arc=10pt,
  drop shadow=popink,left=18pt,right=18pt,top=13pt,bottom=13pt,before skip=6pt,after skip=20pt,
  before upper={\poppill{Sommaire}{poppurple}{white}\par\vspace{9pt}\popsemi\fontsize{10.6}{14.5}\selectfont\color{popink}}}

% Signature de fin de cours — poussée tout en bas de la dernière page
\newcommand{\findecours}{%
  \par\vfill\nopagebreak
  \begin{center}\popsquiggle{poporange}\end{center}\vspace{4pt}
  \signatureonbuch
}

%% FIGFIX-BEGIN v5 — correctifs globaux des figures (docs/onbuch-generation/tools/figfix)
\usepackage{adjustbox}\usepackage{etoolbox}\tcbuselibrary{raster}
% toute figure TikZ/pgfplots plus large que la ligne est réduite à la largeur disponible
\BeforeBeginEnvironment{tikzpicture}{\begin{adjustbox}{max width=\linewidth}}
\AfterEndEnvironment{tikzpicture}{\end{adjustbox}}
% légendes pgfplots lisibles par-dessus les courbes
\pgfplotsset{every axis/.append style={legend style={fill=white,fill opacity=0.92,text opacity=1,draw=black!15,inner sep=3pt}}}
% étiquettes d'arêtes (syntaxe edge["..."]) avec fond blanc
\tikzset{every edge quotes/.append style={fill=white,inner sep=1.2pt}}
%% FIGFIX-END
\begin{document}

\pagegarde{Acides et bases en solution aqueuse : eau, autoprotolyse et pH}{Chimie}{Tle D}{Module 2 · Acides et bases}{06}{3 h (cours + TP)}{Cette leçon explore le rôle central de l'eau comme solvant ionisant, introduit la théorie acido-basique de Brønsted et construit l'échelle de pH à partir de l'autoprotolyse de l'eau. Elle donne les outils de calcul (Ke, pKe, concentrations, dilutions) et les techniques de mesure (papier pH, pH-mètre, indicateurs colorés naturels) indispensables pour le Baccalauréat C, D et E.
\vspace{4pt}
\begin{multicols}{2}\small
\begin{itemize}
\item À la fin de cette leçon, je sais décrire la structure de la molécule d'eau et expliquer ses propriétés de solvant (dissolvant, dispersant, ionisant, solvatant)
\item À la fin de cette leçon, je sais définir un acide et une base selon Brønsted et identifier les couples acide/base dans une réaction
\item À la fin de cette leçon, je sais écrire l'équation-bilan de l'autoprotolyse de l'eau et utiliser le produit ionique Ke
\item À la fin de cette leçon, je sais calculer [H3O+] et [OH-] à partir du pH et du pKe, en faisant les bonnes approximations
\item À la fin de cette leçon, je sais calculer le pH d'une solution et conclure sur son caractère acide, neutre ou basique à 25 °C
\item À la fin de cette leçon, je sais mesurer un pH avec le papier pH, un pH-mètre ou un indicateur coloré, en précisant la précision de chaque méthode
\end{itemize}\end{multicols}}

\begin{sommaire}
\begin{multicols}{2}
\begin{enumerate}
\item L'eau, un solvant exceptionnel
\item Acides et bases selon Brønsted
\item Autoprotolyse de l'eau et produit ionique
\item Le pH : définition, calcul et caractère des solutions
\item Mesure du pH et indicateurs colorés
\item Préparation de solutions par dilution
\item Travaux pratiques
\item Méthodes et exercices résolus
\item Exercices
\item Corrigés des exercices
\item Fiche bilan
\end{enumerate}
\end{multicols}
\end{sommaire}

\begin{objectifs}
\begin{itemize}
\item À la fin de cette leçon, je sais décrire la structure de la molécule d'eau et expliquer ses propriétés de solvant (dissolvant, dispersant, ionisant, solvatant)
\item À la fin de cette leçon, je sais définir un acide et une base selon Brønsted et identifier les couples acide/base dans une réaction
\item À la fin de cette leçon, je sais écrire l'équation-bilan de l'autoprotolyse de l'eau et utiliser le produit ionique Ke
\item À la fin de cette leçon, je sais calculer [H3O+] et [OH-] à partir du pH et du pKe, en faisant les bonnes approximations
\item À la fin de cette leçon, je sais calculer le pH d'une solution et conclure sur son caractère acide, neutre ou basique à 25 °C
\item À la fin de cette leçon, je sais mesurer un pH avec le papier pH, un pH-mètre ou un indicateur coloré, en précisant la précision de chaque méthode
\item À la fin de cette leçon, je sais préparer une solution diluée à partir d'une solution concentrée en utilisant le facteur de dilution
\item À la fin de cette leçon, je sais préparer un indicateur coloré naturel (chou rouge, bissap) et interpréter sa zone de virage
\end{itemize}
\end{objectifs}

\begin{prerequis}
\begin{itemize}
\item Notion de solution, soluté, solvant et concentration molaire (Première C/D)
\item Écriture et équilibrage d'équations-bilans de réactions chimiques (Seconde)
\item Ionisation des acides forts et bases fortes en solution aqueuse (début de Terminale, leçon sur les réactions acido-basiques)
\item Fonction logarithme décimal et notation scientifique (mathématiques)
\item Protocoles de base du laboratoire : pipette jaugée, fiole jaugée, bécher
\end{itemize}
\end{prerequis}

\newpage

\coursec{L'eau, un solvant exceptionnel}

\emph{Situation d'accroche.} Amina, élève en Terminale C à Yaoundé, récupère à l'hôpital le bilan médical de sa mère : le pH sanguin indique 7,38. À la maison, son père ajoute du vinaigre dans l'eau de rinçage du linge, et sa grand-mère prépare du jus de bissap qui vire au rouge vif dès qu'on y presse un citron. Trois scènes ordinaires, trois mystères. Pourquoi le sang doit-il garder un pH presque constant, sous peine de danger de mort ? Pourquoi le bissap change-t-il de couleur avec le citron ? Et comment le laborantin prépare-t-il ses solutions diluées avec précision ?

Toutes ces questions trouvent leur réponse dans un phénomène invisible à l'œil nu : le comportement de la molécule d'eau, son \cle{autoprotolyse} et la notion de \cle{pH}. Mais avant de parler de pH, il faut comprendre pourquoi l'eau, ce liquide si banal, est en réalité un solvant absolument exceptionnel. C'est l'objet de cette première section.

\textbf{Problématique :} quelles propriétés microscopiques de la molécule d'eau expliquent son extraordinaire aptitude à dissoudre, disperser et ioniser les espèces chimiques ?

\courssub{Structure de la molécule d'eau : une molécule coudée et polarisée}

\textbf{Intuition.} Imagine deux ballons (les atomes d'hydrogène) attachés à un pivot central (l'atome d'oxygène). Si la molécule d'eau était droite, comme \ce{CO2}, les deux ballons seraient diamétralement opposés. Or ce n'est pas le cas : la molécule d'eau est \emph{pliée}, et ce pli change tout.

L'atome d'oxygène possède 6 électrons sur sa couche externe. Deux de ces électrons servent à former deux \cle{liaisons covalentes} \ce{O-H} avec les deux atomes d'hydrogène. Les quatre électrons restants forment deux \cle{doublets non liants}, qui repoussent les liaisons \ce{O-H} vers le bas. Résultat : la molécule adopte une géométrie \cle{coudée}, avec un angle $\widehat{\mathrm{HOH}} \approx 104,5^\circ$ (inférieur aux $109,5^\circ$ du tétraèdre parfait, car les doublets non liants « poussent » plus fort que les liaisons).

Mais il y a plus. L'oxygène est beaucoup plus \cle{électronégatif} que l'hydrogène : il attire vers lui le doublet d'électrons de chaque liaison \ce{O-H}. Chaque liaison est donc \cle{polarisée} : l'oxygène porte une charge partielle négative, notée $2\,\delta^-$, et chaque hydrogène une charge partielle positive $\delta^+$.

\begin{popfigure}
\begin{center}
\begin{tikzpicture}[scale=1.2]
  % Atome O
  \shade[ball color=popblue] (0,0) circle (0.5);
  \node[white,font=\bfseries] at (0,0) {O};
  % Atomes H : angle HOH = 104.5 deg -> demi-angle = 52.25 deg par rapport à la bissectrice (axe y)
  \def\rH{1.5}
  \foreach \a in {37.75, 142.25} {
     \shade[ball color=poporange!80] (\a:\rH) circle (0.3);
     \node[white,font=\bfseries\small] at (\a:\rH) {H};
     \node[poporange,font=\bfseries] at (\a:\rH+0.4) {$\delta^+$};
  }
  % Charge partielle O
  \node[popblue,font=\bfseries] at (0,-0.85) {$2\,\delta^-$};
  % Doublets non liants (projection 2D approximative tétraédrique)
  \fill[poppurple!60] (220:0.7) ellipse (0.22 and 0.13);
  \fill[poppurple!60] (320:0.7) ellipse (0.22 and 0.13);
  \node[poppurple,font=\small] at (270:1.25) {doublets non liants};
  % Angle HOH
  \draw[popdark,thick] (37.75:0.8) arc[start angle=37.75,end angle=142.25,radius=0.8];
  \node[popdark,font=\small\bfseries] at (90:0.5) {$104,5^\circ$};
  % Moment dipolaire : du milieu des H (delta+) vers O (delta-)
  % Milieu des H sur l'axe y à rH*cos(52.25) ~ 0.92
  \draw[-{Stealth[length=3mm]},popgreen,very thick] (0,0.85) -- (0,0.15);
  \node[popgreen,font=\small\bfseries,right] at (0.2,0.5) {$\vec{\mu}$};
  \node[popgreen,font=\small,right] at (0.2,0.3) {(de $+$ vers $-$)};
\end{tikzpicture}
\end{center}
\legende{La molécule d'eau : géométrie coudée (angle $\approx 104,5^\circ$), liaisons \ce{O-H} polarisées et moment dipolaire résultant dirigé des charges positives (H) vers la charge négative (O).}
\end{popfigure}

\begin{attention}[Erreur fréquente : la molécule d'eau n'est pas linéaire]
Beaucoup d'élèves dessinent \ce{H-O-H} en ligne droite. C'est faux ! Si la molécule était linéaire, les deux moments dipolaires des liaisons \ce{O-H} s'annuleraient et l'eau serait \emph{apolaire} : elle ne dissoudrait ni le sel, ni le sucre, et la vie telle que nous la connaissons serait impossible. Retiens : \textbf{coudée $\Rightarrow$ polaire}.
\end{attention}

\courssub{Polarité de la molécule : l'eau est un solvant polaire}

Puisque la molécule est coudée, le \cle{barycentre des charges partielles positives} (situé entre les deux H) ne coïncide pas avec le \cle{barycentre des charges partielles négatives} (situé sur l'O). La molécule possède donc un \cle{moment dipolaire} $\vec{\mu}$ non nul : c'est un \emph{dipôle électrique} permanent.

\begin{propriete}[Conséquence fondamentale]
L'eau est un \cle{solvant polaire}. Elle interagit fortement, par des attractions électrostatiques, avec toutes les espèces chargées (ions) ou polarisées (molécules polaires comme l'éthanol, le glucose, le chlorure d'hydrogène). À l'inverse, elle n'interagit presque pas avec les espèces apolaires (huile, essence, graisses) : c'est pourquoi l'huile de palme et l'eau ne se mélangent pas.
\end{propriete}

\begin{savaistu}[L'eau, « solvant universel » ?]
On surnomme parfois l'eau « solvant universel ». C'est un abus de langage : elle ne dissout pas les corps gras ni la plupart des plastiques ! Mais elle dissout plus d'espèces ioniques et polaires que n'importe quel autre solvant courant. Sans sa polarité, le sang — qui est une solution aqueuse — ne pourrait pas transporter les sels minéraux (\ce{Na+}, \ce{K+}, \ce{Ca^2+}, \ce{Cl-}) indispensables au fonctionnement de ton cœur et de tes muscles. Le pH sanguin de la maman d'Amina, 7,38, n'a de sens que parce que le sang est une \emph{solution aqueuse}.
\end{savaistu}

\courssub{Pouvoir dissolvant et pouvoir dispersant : le cas du sel de cuisine}

\textbf{Intuition.} Verse une cuillère de sel dans un verre d'eau et agite : le sel « disparaît ». Il n'a pas disparu — il a été \emph{démantelé} ion par ion par les molécules d'eau. Décomposons ce phénomène en trois actes.

\textbf{Acte 1 : la dissociation (pouvoir dissolvant).} Le cristal de chlorure de sodium \ce{NaCl} est un empilement régulier d'ions \ce{Na+} et \ce{Cl-} retenus par de fortes attractions électrostatiques (liaison ionique). À la surface du cristal, les molécules d'eau polaires s'approchent : leurs pôles négatifs (O) attirent les ions \ce{Na+}, leurs pôles positifs (H) attirent les ions \ce{Cl-}. Quand l'attraction exercée par les molécules d'eau dépasse l'attraction du réseau cristallin, l'ion est \emph{arraché}. L'eau \cle{dissocie} le cristal.

\textbf{Acte 2 : la dispersion (pouvoir dispersant).} Une fois arrachés, les ions ne restent pas collés au cristal : agités par le mouvement thermique des molécules d'eau, ils sont \cle{dispersés} dans tout le volume de la solution. C'est pourquoi, après agitation, chaque gorgée d'eau salée a exactement le même goût : la solution est \emph{homogène}.

\textbf{Acte 3 : l'hydratation (pouvoir solvatant).} Chaque ion en solution s'entoure d'une « couronne » de molécules d'eau orientées : c'est la \cle{solvatation}, appelée \cle{hydratation} quand le solvant est l'eau. L'ion \ce{Na+} (positif) est entouré par les pôles \emph{négatifs} (atomes O) des molécules d'eau ; l'ion \ce{Cl-} (négatif) est entouré par les pôles \emph{positifs} (atomes H). On écrit parfois \ce{Na+(aq)} et \ce{Cl-(aq)}, où « aq » signifie \emph{aqueux}, c'est-à-dire hydraté.

L'équation de dissolution s'écrit :
\[ \ce{NaCl(s) ->[eau] Na+(aq) + Cl-(aq)} \]
Cette transformation est totale, rapide, à température ambiante, sans catalyseur. Remarque l'écriture : le solide \ce{NaCl} n'existe plus comme tel en solution ; il n'y a que des ions hydratés.

\begin{popfigure}
\begin{center}
\begin{tikzpicture}[scale=0.92]
  % Étape 1 : cristal
  \begin{scope}
    \foreach \x in {0,0.7,1.4}{
      \foreach \y in {0,0.7}{
        \shade[ball color=popblue] (\x,\y) circle (0.16);
        \shade[ball color=popgreen] (\x+0.35,\y+0.35) circle (0.2);
      }}
    \node[font=\small\bfseries,popdark] at (1.05,-0.7) {1. Dissociation};
    \node[font=\scriptsize,popmuted] at (1.05,-1.15) {l'eau arrache les ions};
    \draw[-{Stealth},poporange,thick] (2.1,0.7) -- (2.9,0.7);
    \node[font=\scriptsize,poporange] at (2.5,1.0) {eau};
  \end{scope}
  % Étape 2 : dispersion
  \begin{scope}[xshift=3.6cm]
    \shade[ball color=popblue] (0.2,1.2) circle (0.16);
    \shade[ball color=popgreen] (1.6,1.4) circle (0.2);
    \shade[ball color=popgreen] (0.5,0.2) circle (0.2);
    \shade[ball color=popblue] (1.8,0.4) circle (0.16);
    \shade[ball color=popblue] (1.0,0.9) circle (0.16);
    \shade[ball color=popgreen] (1.1,1.9) circle (0.2);
    \node[font=\small\bfseries,popdark] at (1.05,-0.7) {2. Dispersion};
    \node[font=\scriptsize,popmuted] at (1.05,-1.15) {ions répartis dans le volume};
    \draw[-{Stealth},poporange,thick] (2.3,0.7) -- (3.1,0.7);
  \end{scope}
  % Étape 3 : hydratation
  \begin{scope}[xshift=7.6cm,yshift=0.55cm]
    % Na+ hydraté
    \shade[ball color=popblue] (0,0) circle (0.16);
    \foreach \a in {30,90,150,210,270,330}{
      \draw[popblue!50,thick] (\a:0.2) -- (\a:0.5);
      \fill[popblue!35] (\a:0.62) circle (0.13);
      \node[font=\tiny,popblue] at (\a:0.62) {O};
    }
    \node[font=\scriptsize\bfseries,popblue] at (0,-1.05) {\ce{Na+} entouré par les O ($\delta^-$)};
    % Cl- hydraté
    \begin{scope}[xshift=2.6cm]
      \shade[ball color=popgreen] (0,0) circle (0.2);
      \foreach \a in {30,90,150,210,270,330}{
        \draw[poporange!60,thick] (\a:0.24) -- (\a:0.5);
        \fill[poporange!45] (\a:0.62) circle (0.11);
        \node[font=\tiny,popdark] at (\a:0.62) {H};
      }
      \node[font=\scriptsize\bfseries,popgreen!70!black] at (0,-1.05) {\ce{Cl-} entouré par les H ($\delta^+$)};
    \end{scope}
    \node[font=\small\bfseries,popdark] at (1.3,-1.75) {3. Hydratation (solvatation)};
  \end{scope}
\end{tikzpicture}
\end{center}
\legende{Les trois étapes de la dissolution de \ce{NaCl} dans l'eau : dissociation du cristal, dispersion des ions, puis hydratation (les molécules d'eau s'orientent selon la charge de l'ion).}
\end{popfigure}

\begin{attention}[Dissolution $\neq$ fusion : ne confonds pas !]
\textbf{Fondre} du sel, c'est le chauffer jusqu'à \SI{801}{\celsius} : on obtient du \ce{NaCl} liquide, composé d'ions mobiles mais \emph{sans eau}. \textbf{Dissoudre} du sel, c'est le mettre dans l'eau à température ambiante : c'est \emph{l'eau} qui dissocie le cristal, pas la chaleur. De même, dans \ce{NaCl(s) -> Na+(aq) + Cl-(aq)}, les ions existaient déjà dans le cristal ; l'eau ne les a pas « créés », elle les a \emph{séparés et hydratés}. Au Baccalauréat, écrire que l'eau « transforme NaCl en ions » sans préciser que le cristal était déjà ionique est une erreur classique.
\end{attention}

\courssub{Pouvoir ionisant : l'eau crée des ions}

\textbf{Intuition.} Jusqu'ici, l'eau séparait des ions préexistants. Mais elle sait faire mieux : elle peut \emph{fabriquer} des ions à partir de molécules qui n'en contenaient pas.

Le chlorure d'hydrogène \ce{HCl} est un gaz composé de \emph{molécules} : la liaison \ce{H-Cl} est covalente mais fortement polarisée (Cl est très électronégatif, il porte $\delta^-$ et H porte $\delta^+$). Quand on barbote ce gaz dans l'eau, une molécule d'eau capte le proton \ce{H+} de \ce{HCl} : la liaison covalente \ce{H-Cl} se \emph{rompt}, et il se forme deux ions là où il n'y avait que des molécules :
\[ \ce{HCl(g) + H2O(l) -> H3O+(aq) + Cl-(aq)} \]
Cette réaction est \textbf{totale}, rapide et exothermique. L'ion \ce{H3O+} est appelé ion \cle{oxonium} (ou hydronium) : c'est \emph{le} responsable de l'acidité en solution aqueuse, celui-là même dont dépendra le pH du sang d'Amina.

\begin{definition}[Solvant ionisant]
Un \cle{solvant ionisant} est un solvant capable de transformer un composé moléculaire (constitué de molécules) en ions en solution. L'eau est ionisante vis-à-vis des composés à liaison covalente fortement polarisée, comme \ce{HCl}, \ce{H2SO4} ou \ce{NH3}. Il faut distinguer :
\begin{itemize}
  \item la \cle{dissociation} : l'eau sépare des ions \emph{déjà présents} dans un solide ionique (cas de \ce{NaCl}) ;
  \item l'\cle{ionisation} : l'eau \emph{crée} des ions par rupture d'une liaison covalente polarisée (cas de \ce{HCl}).
\end{itemize}
\end{definition}

\begin{popfigure}
\begin{center}
\begin{tikzpicture}[scale=1.0]
  % HCl
  \shade[ball color=poporange!80] (0,0) circle (0.28);
  \node[white,font=\small\bfseries] at (0,0) {H};
  \shade[ball color=popgreen] (0.85,0) circle (0.42);
  \node[white,font=\small\bfseries] at (0.85,0) {Cl};
  \node[font=\scriptsize,poporange] at (0,0.5) {$\delta^+$};
  \node[font=\scriptsize,popgreen!60!black] at (0.85,0.65) {$\delta^-$};
  \node[font=\small,popdark] at (0.4,-0.85) {\ce{HCl} (molécule)};
  % H2O
  \begin{scope}[xshift=2.6cm]
    \shade[ball color=popblue] (0,0) circle (0.4);
    \node[white,font=\small\bfseries] at (0,0) {O};
    \shade[ball color=poporange!80] (115:0.85) circle (0.22);
    \shade[ball color=poporange!80] (65:0.85) circle (0.22);
    \node[font=\scriptsize,popblue] at (0,-0.62) {$2\,\delta^-$};
    \node[font=\small,popdark] at (0,-1.1) {\ce{H2O}};
  \end{scope}
  % Flèche de transfert de H+
  \draw[-{Stealth[length=3mm]},poppink,very thick] (0.15,0.35) .. controls (1.2,1.5) and (2.0,1.5) .. (2.55,0.5);
  \node[poppink,font=\small\bfseries] at (1.35,1.55) {transfert de \ce{H+}};
  % Flèche réaction
  \draw[-{Stealth},popdark,very thick] (3.6,0) -- (4.6,0);
  % Produits
  \begin{scope}[xshift=5.2cm]
    \shade[ball color=popblue] (0.5,0) circle (0.4);
    \node[white,font=\small\bfseries] at (0.5,0) {O};
    \shade[ball color=poporange!80] (0.5,0) ++(100:0.85) circle (0.2);
    \shade[ball color=poporange!80] (0.5,0) ++(20:0.85) circle (0.2);
    \shade[ball color=poporange!80] (0.5,0) ++(200:0.85) circle (0.2);
    \node[font=\small\bfseries,popblue] at (0.5,-1.1) {\ce{H3O+}};
    \node[font=\large,popdark] at (1.9,0) {+};
    \shade[ball color=popgreen] (3.0,0) circle (0.42);
    \node[white,font=\small\bfseries] at (3.0,0) {Cl};
    \node[font=\small\bfseries,popgreen!60!black] at (3.0,-1.1) {\ce{Cl-}};
  \end{scope}
\end{tikzpicture}
\end{center}
\legende{Ionisation de \ce{HCl} par l'eau : la molécule d'eau capte le proton \ce{H+} de la liaison polarisée \ce{H-Cl}. Deux ions, \ce{H3O+} et \ce{Cl-}, apparaissent là où il n'y avait que des molécules.}
\end{popfigure}

\courssub{Conséquence : les solutions aqueuses conduisent le courant}

Une solution contenant des ions mobiles est appelée un \cle{électrolyte}. Les ions hydratés, porteurs de charges électriques, se déplacent sous l'action d'un champ électrique : les cations (\ce{Na+}, \ce{K+}, \ce{H3O+}) migrent vers la cathode (pôle $-$), les anions (\ce{Cl-}, \ce{MnO4-}, \ce{OH-}) vers l'anode (pôle $+$). La solution conduit donc le courant électrique, alors que l'eau pure ne conduit quasiment pas (nous verrons pourquoi dans la section 3 : ses ions sont ultra-minoritaires).

\begin{experience}[Mise en évidence : le test de conductivité]
\textbf{Dispositif :} un circuit série comprenant un générateur, une lampe et deux électrodes plongeant dans un bécher.\\
\textbf{Protocole et observations :}
\begin{itemize}
  \item avec de l'eau distillée : la lampe ne brille pas (ou très faiblement) ;
  \item avec de l'eau salée ou une solution de \ce{KMnO4} : la lampe brille nettement ;
  \item avec une solution de sucre : la lampe ne brille pas.
\end{itemize}
\textbf{Interprétation :} l'eau salée et la solution de permanganate contiennent des ions mobiles (électrolytes). Le sucre, molécule polaire, se dissout bien mais \emph{sans s'ioniser} : la solution de glucose ne conduit pas le courant. \textbf{Dissolution ne signifie pas forcément ionisation !}
\end{experience}

\begin{exemplebox}[Dissolution du permanganate de potassium \ce{KMnO4}]
Le permanganate de potassium, utilisé en pharmacie comme antiseptique (les « gouttes violettes » pour désinfecter), est un solide ionique \ce{K+MnO4-}. Plongé dans l'eau, il se dissout selon :
\[ \ce{KMnO4(s) ->[eau] K+(aq) + MnO4-(aq)} \]
La solution obtenue est \textbf{violette} (couleur caractéristique de l'ion permanganate hydraté \ce{MnO4-}) et \textbf{conductrice}, car elle contient des ions \ce{K+} et \ce{MnO4-} hydratés et mobiles. Si l'on dissout $m = \SI{1,58}{g}$ de \ce{KMnO4} ($M = \SI{158}{g.mol^{-1}}$) dans $V = \SI{1,00}{L}$ d'eau, la concentration de chaque ion est :
\[ [\ce{K+}] = [\ce{MnO4-}] = \frac{n}{V} = \frac{m}{M \times V} = \frac{1,58}{158 \times 1{,}00} = \SI{1,0e-2}{mol.L^{-1}} \]
\end{exemplebox}

\begin{exoresolu}[Conductivité d'une solution de chlorure de calcium]
\textbf{Énoncé.} On dissout $n = \SI{0,020}{mol}$ de chlorure de calcium \ce{CaCl2} (solide ionique) dans de l'eau distillée pour obtenir $V = \SI{500}{mL}$ de solution. (1) Écrire l'équation de dissolution. (2) Calculer la concentration de chaque ion en solution. (3) Cette solution conduit-elle le courant ? Justifier.
\tcblower
\textbf{Solution détaillée.}
(1) \ce{CaCl2} est un cristal ionique : l'eau le dissocie, disperse et hydrate ses ions :
\[ \ce{CaCl2(s) ->[eau] Ca^2+(aq) + 2 Cl-(aq)} \]
Attention au coefficient 2 devant \ce{Cl-} : chaque motif \ce{CaCl2} libère \emph{deux} ions chlorure.

(2) Concentration apportée : $c = \dfrac{n}{V} = \dfrac{0{,}020}{0{,}500} = \SI{4,0e-2}{mol.L^{-1}}$.
D'après la stœchiométrie de l'équation :
\[ [\ce{Ca^2+}] = c = \SI{4,0e-2}{mol.L^{-1}} \qquad ; \qquad [\ce{Cl-}] = 2c = \SI{8,0e-2}{mol.L^{-1}} \]
Vérification de l'électroneutralité : $(+2)\times 0{,}040 + (-1)\times 0{,}080 = 0$. La solution est bien électriquement neutre.

(3) Oui : la solution contient des ions hydratés \ce{Ca^2+} et \ce{Cl-} mobiles, c'est un électrolyte. Elle conduit le courant d'autant mieux que les ions sont nombreux et chargés.
\end{exoresolu}

\courssub{L'eau dans la vie quotidienne : du sel de cuisine aux médicaments}

Le pouvoir dissolvant, dispersant, solvatant et ionisant de l'eau explique d'innombrables gestes du quotidien camerounais :

\begin{itemize}
  \item \textbf{En cuisine :} le sel et le cube d'assaisonnement se dissolvent dans la sauce grâce à la dissociation de leurs cristaux ioniques ; le goût se répartit uniformément grâce à la dispersion.
  \item \textbf{En agriculture :} les engrais (nitrates, phosphates, sels de potassium) sont épandus puis dissous par l'eau de pluie ; les racines absorbent les ions \ce{NO3-}, \ce{K+}, \ce{PO4^3-} hydratés directement depuis la solution du sol.
  \item \textbf{En médecine :} les sirops et solutions buvables contiennent des principes actifs dissous et donc rapidement absorbés par l'organisme ; les perfusions (sérum physiologique : \ce{NaCl} à \SI{9}{g.L^{-1}}) sont des solutions aqueuses isotoniques injectées directement dans le sang.
  \item \textbf{À l'hôpital d'Amina :} le pH sanguin de 7,38 n'existe que parce que le sang est une solution aqueuse dont les ions \ce{H3O+} — issus de l'autoprotolyse de l'eau et des acides du métabolisme — sont rigoureusement contrôlés.
\end{itemize}

\begin{aretenir}[L'essentiel de la section]
\begin{itemize}
  \item La molécule d'eau est \cle{coudée} ($\widehat{\mathrm{HOH}} \approx 104,5^\circ$) et \cle{polaire} : charges partielles $\delta^+$ sur les H, $2\,\delta^-$ sur l'O, moment dipolaire non nul.
  \item L'eau possède quatre pouvoirs complémentaires : \cle{dissolvant} (elle dissocie les cristaux ioniques), \cle{dispersant} (les ions se répartissent dans tout le volume), \cle{solvatant} (chaque ion est hydraté : \ce{Na+} par les pôles O, \ce{Cl-} par les pôles H) et \cle{ionisant} (elle crée des ions par rupture de liaisons covalentes polarisées : \ce{HCl + H2O -> H3O+ + Cl-}).
  \item Les solutions aqueuses contenant des ions sont des \cle{électrolytes} : elles conduisent le courant électrique.
  \item Ne jamais confondre \textbf{dissolution} (action de l'eau) et \textbf{fusion} (action de la chaleur), ni \textbf{dissociation} (ions préexistants) et \textbf{ionisation} (ions créés).
\end{itemize}
\end{aretenir}

Maintenant que tu comprends pourquoi l'eau est le théâtre de toute la chimie en solution, il nous reste à définir les acteurs principaux de cette leçon : les acides et les bases selon Brønsted. C'est l'objet de la section suivante.

\coursec{Acides et bases selon Brønsted}

Dans la section précédente, nous avons vu que l'eau, grâce à sa molécule coudée et polaire, est un solvant exceptionnel : elle dissocie les cristaux ioniques, solvate les ions et peut même ioniser certaines molécules. Mais une question demeure : pourquoi le jus de citron pique-t-il la langue, pourquoi la potasse de cuisine rend-elle les mains glissantes, et pourquoi ces deux produits semblent-ils « s'opposer » chimiquement ? La réponse tient en deux mots : \cle{acides} et \cle{bases}. C'est le chimiste danois Johannes \textbf{Brønsted} qui, en 1923, a proposé la définition la plus utile pour nous : tout se joue autour d'une particule minuscule, le \cle{proton} \ce{H+}.

\courssub{Les définitions de Brønsted}

\textbf{Intuition d'abord.} Imagine un match de football : l'acide est le joueur qui \emph{passe} le ballon, la base est celui qui le \emph{reçoit}. Le ballon, c'est le proton \ce{H+}. Sans receveur, pas de passe ; sans base, un acide ne peut pas se manifester. L'acidité et la basicité sont donc des \emph{comportements relatifs}, qui s'expriment toujours face à un partenaire.

\begin{definition}[Acide et base selon Brønsted]
\begin{itemize}
\item Un \cle{acide} (au sens de Brønsted) est une espèce chimique (atome, molécule ou ion) capable de \textbf{céder un ou plusieurs protons} \ce{H+}.
\item Une \cle{base} (au sens de Brønsted) est une espèce chimique capable de \textbf{capter un ou plusieurs protons} \ce{H+}.
\end{itemize}
\end{definition}

Quelques exemples familiers au Cameroun :
\begin{itemize}
\item le \textbf{citron} et le \textbf{tamarin} contiennent des acides (acide citrique, acide tartrique) ;
\item le \textbf{vinaigre} utilisé pour assaisonner le poisson braisé contient de l'\cle{acide éthanoïque} \ce{CH3COOH} ;
\item la \textbf{potasse de cuisine} et la \textbf{cendre de bois} utilisée en savonnerie traditionnelle contiennent des bases (\ce{KOH}, \ce{K2CO3}) qui captent des protons et libèrent des ions \ce{OH-} en solution.
\end{itemize}

\begin{attention}[Le proton n'existe pas libre dans l'eau]
Erreur très fréquente au Baccalauréat : écrire qu'une solution contient des ions \ce{H+} « libres ». C'est \textbf{faux}. Le proton \ce{H+} est un noyau d'hydrogène nu, extrêmement réactif : en solution aqueuse, il se fixe immédiatement sur le doublet non liant d'une molécule d'eau pour former l'ion \cle{oxonium} \ce{H3O+} :
\[ \ce{H+ + H2O -> H3O+} \]
On écrit parfois \ce{H+} par abréviation, mais il faut toujours garder à l'esprit qu'il s'agit en réalité de \ce{H3O+}.
\end{attention}

\courssub{Couples acide/base conjugués et demi-équations}

Lorsqu'un acide cède son proton, il se transforme en une nouvelle espèce qui, elle, pourrait le reprendre : c'est sa \cle{base conjuguée}. Les deux espèces forment un \cle{couple acide/base}, relié par une \cle{demi-équation acido-basique} :

\[ \text{acide} \; \ce{<=>} \; \text{base} + \ce{H+} \]

\begin{definition}[Couple conjugué]
Deux espèces chimiques forment un \textbf{couple acide/base conjugués} si l'on passe de l'une à l'autre par gain ou perte d'un proton \ce{H+}. On note le couple acide/base, l'acide étant écrit en premier : \ce{CH3COOH}/\ce{CH3COO-}.
\end{definition}

Exemples fondamentaux :
\begin{itemize}
\item \ce{HCl}/\ce{Cl-} : \quad \ce{HCl <=> Cl- + H+} \quad (acide chlorhydrique / ion chlorure) ;
\item \ce{NH4+}/\ce{NH3} : \quad \ce{NH4+ <=> NH3 + H+} \quad (ion ammonium / ammoniac) ;
\item \ce{CH3COOH}/\ce{CH3COO-} : \quad \ce{CH3COOH <=> CH3COO- + H+} \quad (acide éthanoïque / ion éthanoate).
\end{itemize}

Remarque essentielle : la double flèche \ce{<=>} signifie que la transformation est possible \emph{dans les deux sens}. L'ion éthanoate \ce{CH3COO-} peut capter un proton : c'est bien une base, même si elle provient d'un acide.

\begin{popfigure}
\begin{center}
\begin{tabularx}{\linewidth}{|l|l|X|}
\hline
\rowcolor{popblueL} \textbf{Couple acide/base} & \textbf{Demi-équation} & \textbf{Exemple d'usage} \\
\hline
\ce{HCl}/\ce{Cl-} & \ce{HCl <=> Cl- + H+} & Acide chlorhydrique (détartrant, estomac) \\
\hline
\ce{H3O+}/\ce{H2O} & \ce{H3O+ <=> H2O + H+} & Ion oxonium, responsable de l'acidité \\
\hline
\ce{H2O}/\ce{OH-} & \ce{H2O <=> OH- + H+} & Eau jouant le rôle d'acide \\
\hline
\ce{NH4+}/\ce{NH3} & \ce{NH4+ <=> NH3 + H+} & Ammoniac des produits ménagers \\
\hline
\ce{CH3COOH}/\ce{CH3COO-} & \ce{CH3COOH <=> CH3COO- + H+} & Vinaigre de table \\
\hline
\end{tabularx}
\end{center}
\legende{Couples acide/base usuels et leurs demi-équations acido-basiques.}
\end{popfigure}

\courssub{La réaction acido-basique : un transfert de proton}

\textbf{Idée directrice.} Une \cle{réaction acido-basique} est un transfert de proton \ce{H+} entre l'\textbf{acide d'un couple} (qui le cède) et la \textbf{base d'un autre couple} (qui le capte). Deux couples sont donc toujours mis en jeu.

Pour construire l'équation-bilan, on écrit les deux demi-équations, puis on les additionne membre à membre : le proton \ce{H+} « passeur » disparaît de l'écriture finale.

\begin{exemplebox}[Réaction entre l'ammoniac et l'eau]
L'ammoniac \ce{NH3} (base du couple \ce{NH4+}/\ce{NH3}) capte un proton cédé par l'eau (acide du couple \ce{H2O}/\ce{OH-}) :
\begin{align*}
\ce{NH3 + H+ &<=> NH4+} &&\text{(base qui capte } \ce{H+}\text{)} \\
\ce{H2O &<=> OH- + H+} &&\text{(acide qui cède } \ce{H+}\text{)} \\
\ce{NH3 + H2O &<=> NH4+ + OH-} &&\text{équation-bilan}
\end{align*}
La réaction est \textbf{partielle} (d'où la double flèche) : l'ammoniac est une base faible. La formation d'ions \ce{OH-} explique le caractère basique des solutions d'ammoniac.
\end{exemplebox}

\begin{popfigure}
\begin{center}
\begin{tikzpicture}[scale=1.0, every node/.style={font=\small}]
% Espèces
\node[font=\large\bfseries, popblue] (nh3) at (0,0) {\ce{NH3}};
\node[font=\large\bfseries, poporange] (h2o) at (3.2,0) {\ce{H2O}};
\node[font=\large\bfseries, popblue] (nh4) at (7.2,0) {\ce{NH4+}};
\node[font=\large\bfseries, poporange] (oh) at (10.4,0) {\ce{OH-}};
% Flèche de transfert du proton
\draw[-{Stealth[length=3mm]}, poppink, very thick, bend left=45] (h2o.north) to node[above, font=\small\bfseries, poppink]{transfert de \ce{H+}} (nh3.north);
% Flèche de réaction
\draw[-{Stealth[length=3mm]}, popdark, thick] (4.4,0) -- (6.2,0);
% Rôles
\node[popblue, font=\small\itshape] at (0,-0.8) {base (capte \ce{H+})};
\node[poporange, font=\small\itshape] at (3.2,-0.8) {acide (cède \ce{H+})};
\node[popblue, font=\small\itshape] at (7.2,-0.8) {acide conjugué};
\node[poporange, font=\small\itshape] at (10.4,-0.8) {base conjuguée};
% Couples
\draw[popline, thick, rounded corners] (-1.0,-1.4) rectangle (8.0,-2.1);
\node[popmuted, font=\small] at (3.5,-1.75) {couple 1 : \ce{NH4+}/\ce{NH3}};
\draw[popline, thick, rounded corners] (2.2,-2.5) rectangle (11.4,-3.2);
\node[popmuted, font=\small] at (6.8,-2.85) {couple 2 : \ce{H2O}/\ce{OH-}};
\end{tikzpicture}
\end{center}
\legende{Transfert de proton entre l'eau (acide) et l'ammoniac (base) : \ce{NH3 + H2O <=> NH4+ + OH-}.}
\end{popfigure}

\begin{methode}[Identifier les deux couples d'une réaction acido-basique]
\begin{enumerate}
\item Repérer l'espèce qui \textbf{perd un proton} entre les réactifs et les produits : c'est l'\textbf{acide} ; ce qu'il devient est sa base conjuguée.
\item Repérer l'espèce qui \textbf{gagne un proton} : c'est la \textbf{base} ; ce qu'elle devient est son acide conjugué.
\item Écrire les deux couples (acide en premier), puis les deux demi-équations.
\item Additionner les demi-équations : les \ce{H+} se simplifient ; vérifier la conservation des atomes et des charges.
\end{enumerate}
\end{methode}

\courssub{L'eau, une espèce ampholyte}

Nous venons de voir l'eau céder un proton face à \ce{NH3} : elle s'est comportée en \emph{acide}. Mais face à l'acide éthanoïque du vinaigre, c'est elle qui \emph{cape} le proton :
\[ \ce{CH3COOH + H2O <=> CH3COO- + H3O+} \]
Ici, \ce{H2O} est la base (elle devient \ce{H3O+}) et \ce{CH3COOH} est l'acide. L'eau peut donc jouer les deux rôles !

\begin{definition}[Ampholyte]
Un \cle{ampholyte} (ou espèce \cle{amphotère}) est une espèce chimique capable de se comporter \textbf{tantôt en acide, tantôt en base}, selon le partenaire rencontré. L'eau est l'ampholyte par excellence : elle appartient à deux couples, \ce{H3O+}/\ce{H2O} (où \ce{H2O} est la base) et \ce{H2O}/\ce{OH-} (où \ce{H2O} est l'acide). L'ion hydrogénocarbonate \ce{HCO3-}, présent dans le bicarbonate de soude, est un autre ampholyte classique.
\end{definition}

Cette double personnalité de l'eau aura une conséquence capitale dans la section suivante : l'eau peut réagir \emph{avec elle-même} — c'est l'\cle{autoprotolyse}.

\courssub{Complément Bac : acides forts, acides faibles, bases fortes, bases faibles}

Tous les acides ne cèdent pas leur proton avec la même facilité. On distingue :

\begin{itemize}
\item les \cle{acides forts} : leur réaction avec l'eau est \textbf{totale} (flèche simple \ce{->}). Exemple : \ce{HCl + H2O -> H3O+ + Cl-}. Toutes les molécules de \ce{HCl} sont transformées ; la base conjuguée \ce{Cl-} est dite \emph{indifférente} (elle ne capte plus de proton dans l'eau) ;
\item les \cle{acides faibles} : leur réaction avec l'eau est \textbf{partielle} (double flèche \ce{<=>}), un équilibre s'établit. Exemple : l'acide éthanoïque du vinaigre, \ce{CH3COOH + H2O <=> CH3COO- + H3O+} ;
\item les \cle{bases fortes} : réaction totale avec l'eau. Exemple : la potasse \ce{KOH} se dissocie entièrement : \ce{KOH -> K+ + OH-} ; l'ion éthanoate de sodium réagit totalement aussi dans certains cas ;
\item les \cle{bases faibles} : réaction partielle. Exemple : l'ammoniac \ce{NH3 + H2O <=> NH4+ + OH-}.
\end{itemize}

\begin{aretenir}[L'essentiel]
\begin{itemize}
\item Acide de Brønsted = \textbf{donneur} de \ce{H+} ; base de Brønsted = \textbf{accepteur} de \ce{H+}.
\item Toute réaction acido-basique met en jeu \textbf{deux couples} et un \textbf{transfert de proton}.
\item L'eau est un \textbf{ampholyte} : couples \ce{H3O+}/\ce{H2O} et \ce{H2O}/\ce{OH-}.
\item Fort = réaction \textbf{totale} avec l'eau ; faible = réaction \textbf{partielle} (équilibre).
\end{itemize}
\end{aretenir}

\begin{exoresolu}[Identifier acides, bases et couples]
\textbf{Énoncé.} On dissout du chlorure d'ammonium \ce{NH4Cl} dans l'eau ; l'ion ammonium réagit selon : \ce{NH4+ + H2O <=> NH3 + H3O+}. 1) Identifier l'acide et la base parmi les réactifs. 2) Écrire les deux couples mis en jeu. 3) Cette solution est-elle acide ou basique ? Justifier.
\tcblower
\textbf{Solution.} 1) \ce{NH4+} perd un proton pour devenir \ce{NH3} : c'est l'\textbf{acide}. \ce{H2O} gagne un proton pour devenir \ce{H3O+} : c'est la \textbf{base}. \par
2) Couples : \ce{NH4+}/\ce{NH3} et \ce{H3O+}/\ce{H2O}. \par
3) La réaction produit des ions oxonium \ce{H3O+}, responsables de l'acidité : la solution est donc \textbf{acide}. C'est pourquoi le chlorure d'ammonium, utilisé comme engrais, acidifie progressivement les sols.
\end{exoresolu}

\begin{savaistu}[La cendre de bois, base ancestrale]
Dans de nombreuses savonneries artisanales du Cameroun, on filtre de l'eau à travers de la cendre de bois pour obtenir une lessive basique riche en potasse \ce{K2CO3}. Les ions carbonate \ce{CO3^2-} sont des bases de Brønsted : \ce{CO3^2- + H2O <=> HCO3- + OH-}. Les ions \ce{OH-} formés permettent ensuite de transformer l'huile de palme en savon : c'est la saponification, que tu étudieras en chimie organique. Nos grands-mères pratiquaient donc la chimie de Brønsted bien avant 1923 !
\end{savaistu}

\begin{exercice}[À toi de jouer]
Le jus de bissap (infusion de fleurs d'hibiscus) a un goût acidulé dû à l'acide citrique. On écrira cet acide simplement \ce{H3A}. 1) Écrire la demi-équation du couple \ce{H3A}/\ce{H2A-}. 2) Écrire l'équation-bilan de la réaction de \ce{H3A} avec l'eau. 3) Préciser le rôle (acide ou base) joué par l'eau dans cette réaction.
\end{exercice}

\begin{corrige}[Exercice n° 1]
1) \ce{H3A <=> H2A- + H+}. \par
2) L'eau capte le proton : \ce{H2O + H+ <=> H3O+}. En additionnant : \ce{H3A + H2O <=> H2A- + H3O+}. \par
3) L'eau gagne un proton : elle joue le rôle de \textbf{base} (couple \ce{H3O+}/\ce{H2O}). L'acide citrique est un acide faible, d'où la double flèche.
\end{corrige}

\coursec{Autoprotolyse de l'eau et produit ionique}

Dans la section précédente, nous avons découvert la théorie de Brønsted : un \cle{acide} est un donneur de proton \ce{H+}, une \cle{base} est un accepteur de proton, et toute réaction acido-basique est un transfert de proton entre deux couples. Nous avons aussi remarqué que l'eau est \emph{amphotère} : elle peut jouer le rôle d'acide (couple \ce{H2O}/\ce{OH-}) ou celui de base (couple \ce{H3O+}/\ce{H2O}). Cette double personnalité de l'eau a une conséquence extraordinaire que nous allons maintenant explorer : \textbf{l'eau peut réagir avec elle-même}. Cette réaction, appelée \cle{autoprotolyse}, est la clé de voûte de toute l'acido-basie en solution aqueuse : c'est elle qui explique l'existence du pH et qui permet de relier entre elles les concentrations de tous les ions des solutions.

\courssub{Constat expérimental : l'eau pure conduit (très faiblement) le courant}

\begin{experience}[Conductivité de l'eau pure]
\textbf{Dispositif et protocole.} On plonge deux électrodes de platine reliées à un générateur et à un ampèremètre très sensible (ou à un conductimètre) dans de l'eau distillée, soigneusement débarrassée des gaz dissous.

\textbf{Observations.}
\begin{itemize}
\item L'ampèremètre détecte un courant, mais \emph{extrêmement faible} : la conductivité de l'eau pure à \SI{25}{\celsius} vaut environ \SI{5,5e-6}{S.m^{-1}}, soit environ un million de fois moins que celle d'une solution de sel de cuisine à \SI{1}{g.L^{-1}}.
\item L'eau du robinet, elle, conduit nettement mieux le courant, car elle contient des ions dissous (\ce{Ca^2+}, \ce{HCO3-}, \ce{Na+}, \ce{Cl-}\dots).
\end{itemize}

\textbf{Interprétation.} Un courant électrique dans une solution, c'est un déplacement d'\cle{ions}. Si l'eau pure conduit, même très faiblement, c'est qu'elle contient des ions\dots mais d'où viennent-ils ? Il n'y a que des molécules \ce{H2O} ! La seule explication possible : \textbf{les molécules d'eau s'ionisent spontanément entre elles}, en très faible proportion.
\end{experience}

Ce constat est capital. L'eau pure n'est pas un milieu électriquement « vide » : elle contient naturellement des ions \ce{H3O+} et \ce{OH-} en quantités infimes mais parfaitement mesurables. Toute la chimie des solutions aqueuses repose sur ce fait.

\courssub{L'autoprotolyse de l'eau : l'eau réagit avec elle-même}

\textbf{Le mécanisme : un transfert de proton entre deux molécules d'eau}\par

Imaginons deux molécules d'eau qui se rencontrent. L'une d'elles joue le rôle d'\textbf{acide} : elle cède un proton \ce{H+}. L'autre joue le rôle de \textbf{base} : grâce à ses doublets non liants, elle capte ce proton. Il se forme alors simultanément un ion \cle{oxonium} \ce{H3O+} (l'eau qui a gagné un proton) et un ion \cle{hydroxyde} \ce{OH-} (l'eau qui a perdu un proton).

\begin{popfigure}
\begin{center}
\begin{tikzpicture}[scale=1.0, every node/.style={font=\small}]
% Molécule 1 (base)
\draw[fill=popblue!25, draw=popblue, thick] (0,0) circle (0.45) node {O};
\draw[fill=white, draw=popdark, thick] (-0.55,0.55) circle (0.22) node[font=\scriptsize] {H};
\draw[fill=white, draw=popdark, thick] (0.55,0.55) circle (0.22) node[font=\scriptsize] {H};
\draw[thick, popdark] (-0.35,0.35) -- (-0.2,0.2);
\draw[thick, popdark] (0.35,0.35) -- (0.2,0.2);
\node[below, popblue] at (0,-0.55) {\ce{H2O} (base)};
% Molécule 2 (acide)
\draw[fill=poporange!25, draw=poporange, thick] (3.2,0) circle (0.45) node {O};
\draw[fill=white, draw=popdark, thick] (2.65,0.55) circle (0.22) node[font=\scriptsize] {H};
\draw[fill=popgreen!30, draw=popgreen, thick] (3.75,0.55) circle (0.22) node[font=\scriptsize, popdark] {H};
\draw[thick, popdark] (2.85,0.35) -- (3.0,0.2);
\draw[thick, popdark] (3.55,0.35) -- (3.4,0.2);
\node[below, poporange] at (3.2,-0.55) {\ce{H2O} (acide)};
% Transfert du proton
\draw[-{Stealth[length=3mm]}, very thick, popgreen] (3.6,0.85) .. controls (2.4,1.7) and (1.2,1.5) .. (0.35,0.75);
\node[above, popgreen, font=\small\bfseries] at (1.9,1.55) {transfert de \ce{H+}};
% Flèche réactionnelle
\draw[-{Stealth[length=3mm]}, very thick, popdark] (4.5,0) -- (5.7,0);
% Produits : H3O+
\draw[fill=popblue!25, draw=popblue, thick] (6.9,0) circle (0.45) node {O};
\draw[fill=white, draw=popdark, thick] (6.35,0.55) circle (0.2) node[font=\scriptsize] {H};
\draw[fill=white, draw=popdark, thick] (7.45,0.55) circle (0.2) node[font=\scriptsize] {H};
\draw[fill=popgreen!30, draw=popgreen, thick] (6.9,-0.62) circle (0.2) node[font=\scriptsize, popdark] {H};
\draw[thick, popdark] (6.55,0.35) -- (6.7,0.2);
\draw[thick, popdark] (7.25,0.35) -- (7.1,0.2);
\draw[thick, popdark] (6.9,-0.45) -- (6.9,-0.42);
\node[right, popblue] at (7.55,-0.15) {\ce{H3O+}};
\node[below, popblue] at (6.9,-1.0) {ion oxonium};
% +
\node[font=\Large\bfseries, popdark] at (8.9,0) {+};
% OH-
\draw[fill=poporange!25, draw=poporange, thick] (10.1,0) circle (0.45) node {O};
\draw[fill=white, draw=popdark, thick] (9.55,0.55) circle (0.2) node[font=\scriptsize] {H};
\draw[thick, popdark] (9.75,0.35) -- (9.9,0.2);
\node[right, poporange] at (10.6,0) {\ce{OH-}};
\node[below, poporange] at (10.1,-0.6) {ion hydroxyde};
\end{tikzpicture}
\end{center}
\legende{Schéma de l'autoprotolyse de l'eau : une molécule d'eau (acide) cède un proton \ce{H+} à une autre molécule d'eau (base). Il se forme un ion oxonium \ce{H3O+} et un ion hydroxyde \ce{OH-}.}
\end{popfigure}

\textbf{Équation-bilan et caractéristiques de la réaction}\par

L'équation-bilan de l'autoprotolyse s'écrit :

\[ \ce{2 H2O(l) <=> H3O+(aq) + OH-(aq)} \]

Commentons cette équation avec soin :
\begin{itemize}
\item Elle met en jeu les deux couples de l'eau : \ce{H3O+}/\ce{H2O} (l'eau base) et \ce{H2O}/\ce{OH-} (l'eau acide). C'est bien une réaction acido-basique au sens de Brønsted.
\item La double flèche \ce{<=>} traduit un \cle{équilibre chimique} : la réaction inverse (recombinaison de \ce{H3O+} et \ce{OH-} pour redonner de l'eau) a lieu en permanence.
\item La réaction est \textbf{très limitée} : à l'équilibre, il reste énormément plus de molécules \ce{H2O} que d'ions formés. On dit que l'équilibre est « très déplacé vers la gauche ».
\item Elle ne nécessite \textbf{aucun catalyseur} : elle est spontanée et quasi instantanée dans les deux sens.
\item Elle est \textbf{endothermique} : elle absorbe de la chaleur. Ce point aura une conséquence majeure sur l'effet de la température (voir plus loin).
\end{itemize}

\begin{attention}[Écriture de l'équation]
On écrit parfois l'autoprotolyse sous la forme simplifiée \ce{H2O <=> H+ + OH-}. Cette écriture est \textbf{à proscrire} au Baccalauréat : le proton \ce{H+} n'existe pas librement en solution aqueuse, il est toujours fixé sur une molécule d'eau sous la forme \ce{H3O+}. L'équation correcte fait donc apparaître \textbf{deux} molécules d'eau à gauche : \ce{2 H2O <=> H3O+ + OH-}.
\end{attention}

\begin{savaistu}[Une molécule sur 555 millions]
À \SI{25}{\celsius}, dans l'eau pure, seulement \textbf{2 molécules sur environ 555 millions} sont ionisées à un instant donné. Autrement dit, si l'eau d'une piscine olympique représentait les molécules d'eau, les ions formés tiendraient dans un dé à coudre ! Et pourtant, cette ionisation infime suffit à expliquer toute l'acido-basie : les ions \ce{H3O+} et \ce{OH-} sont extrêmement réactifs, et ils se renouvellent en permanence (chaque ion ne « vit » qu'une fraction de seconde avant de se recombiner, pendant que d'autres se forment ailleurs).
\end{savaistu}

\courssub{Le produit ionique de l'eau \texorpdfstring{$K_e$}{Ke}}

\textbf{Définition}\par

Puisque l'autoprotolyse est un équilibre chimique, on peut lui associer une \cle{constante d'équilibre}. Comme l'eau est le solvant (sa « concentration » est constante et absorbée dans la constante), cette constante ne fait intervenir que les deux ions :

\begin{definition}[Produit ionique de l'eau]
Dans \textbf{toute solution aqueuse}, il existe un équilibre entre les ions oxonium et hydroxyde, caractérisé par le \cle{produit ionique de l'eau}, noté $K_e$ :
\[ K_e = [\ce{H3O+}] \times [\ce{OH-}] \]
où les concentrations s'expriment en \si{mol.L^{-1}}. $K_e$ est une constante \textbf{sans unité} qui ne dépend \textbf{que de la température}.

À \SI{25}{\celsius} : $K_e = 1{,}0 \times 10^{-14}$.

On définit aussi : $\mathrm{p}K_e = -\log K_e$, soit à \SI{25}{\celsius} : $\mathrm{p}K_e = 14$.
\end{definition}

\begin{attention}[Unités et nature de Ke]
$K_e$ est \textbf{sans unité} : par convention, chaque concentration est en réalité divisée par la concentration de référence $C^0 = \SI{1}{mol.L^{-1}}$. On écrit donc $K_e = 1{,}0 \times 10^{-14}$ et non « \SI{e-14}{mol^2.L^{-2}} ». Autre piège : $K_e$ n'est \emph{pas} une constante universelle fixe : elle change avec la température (voir plus loin). La valeur $10^{-14}$ n'est valable qu'à \SI{25}{\celsius}.
\end{attention}

\textbf{Une relation universelle : valable pour TOUTE solution aqueuse}\par

Voici le point le plus important de cette section, et l'un des plus importants de tout le module :

\begin{propriete}[Universalité du produit ionique]
Dans \textbf{toute solution aqueuse}, qu'elle soit acide, neutre ou basique, quels que soient les solutés qu'elle contient, la relation suivante est \emph{toujours} vérifiée à l'équilibre :
\[ [\ce{H3O+}] \times [\ce{OH-}] = K_e \]
\begin{itemize}
\item Dans une solution \textbf{acide} : $[\ce{H3O+}] > [\ce{OH-}]$, mais le produit des deux reste égal à $K_e$.
\item Dans une solution \textbf{basique} : $[\ce{OH-}] > [\ce{H3O+}]$, mais le produit reste égal à $K_e$.
\item Dans une solution \textbf{neutre} : $[\ce{H3O+}] = [\ce{OH-}]$.
\end{itemize}
Ajouter un acide ou une base ne « supprime » jamais un des deux ions : il y a toujours des ions \ce{H3O+} \emph{et} des ions \ce{OH-} dans toute solution aqueuse ; seules leurs proportions changent.
\end{propriete}

C'est une idée contre-intuitive mais fondamentale : même dans une solution concentrée d'acide chlorhydrique, il reste des ions \ce{OH-} (en quantité infime) ; même dans une solution de soude concentrée, il reste des ions \ce{H3O+}. L'équilibre d'autoprotolyse se déplace, mais ne disparaît jamais.

\textbf{Cas de l'eau pure à 25 °C}\par

Dans l'eau pure, la stœchiométrie de l'équation \ce{2 H2O <=> H3O+ + OH-} impose que les deux ions soient formés en \textbf{quantités égales} :
\[ [\ce{H3O+}] = [\ce{OH-}] \]
En injectant dans le produit ionique à \SI{25}{\celsius} :
\[ [\ce{H3O+}]^2 = K_e = 1{,}0 \times 10^{-14} \quad \Longrightarrow \quad [\ce{H3O+}] = [\ce{OH-}] = \sqrt{10^{-14}} = 1{,}0 \times 10^{-7}~\si{mol.L^{-1}} \]

\begin{aretenir}[Les valeurs de référence à 25 °C]
Dans l'eau pure à \SI{25}{\celsius} :
\[ [\ce{H3O+}] = [\ce{OH-}] = 1{,}0 \times 10^{-7}~\si{mol.L^{-1}} \qquad K_e = 1{,}0 \times 10^{-14} \qquad \mathrm{p}K_e = 14 \]
Ces valeurs sont à connaître par cœur : elles servent de point de repère pour toute la suite (pH, neutralité, dosages).
\end{aretenir}

\textbf{Exploitation : calculer une concentration à partir de l'autre}\par

\begin{methode}[Passer de \texorpdfstring{$[\ce{H3O+}]$}{[H3O+]} à \texorpdfstring{$[\ce{OH-}]$}{[OH-]} et inversement]
Valable pour \textbf{toute} solution aqueuse, à température connue :
\begin{enumerate}
\item Identifier la température et en déduire la valeur de $K_e$ (à \SI{25}{\celsius} : $K_e = 10^{-14}$).
\item Écrire le produit ionique : $[\ce{H3O+}] \times [\ce{OH-}] = K_e$.
\item Isoler l'inconnue :
\[ [\ce{OH-}] = \frac{K_e}{[\ce{H3O+}]} \qquad \text{ou} \qquad [\ce{H3O+}] = \frac{K_e}{[\ce{OH-}]} \]
\item Vérifier la cohérence : si la solution est acide, on doit trouver $[\ce{OH-}] < 10^{-7}~\si{mol.L^{-1}}$ ; si elle est basique, $[\ce{OH-}] > 10^{-7}~\si{mol.L^{-1}}$ (à \SI{25}{\celsius}).
\end{enumerate}
\end{methode}

\begin{exemplebox}[Exemple chiffré détaillé]
À \SI{25}{\celsius}, une solution aqueuse (par exemple du jus de citron dilué) a une concentration en ions oxonium $[\ce{H3O+}] = 2{,}0 \times 10^{-3}~\si{mol.L^{-1}}$. Calculons $[\ce{OH-}]$.

\textbf{Données :} $K_e = 1{,}0 \times 10^{-14}$ à \SI{25}{\celsius}.

\textbf{Calcul :}
\[ [\ce{OH-}] = \frac{K_e}{[\ce{H3O+}]} = \frac{1{,}0 \times 10^{-14}}{2{,}0 \times 10^{-3}} = 5{,}0 \times 10^{-12}~\si{mol.L^{-1}} \]

\textbf{Vérification :} $[\ce{OH-}] = 5{,}0 \times 10^{-12} \ll 10^{-7}~\si{mol.L^{-1}}$ : les ions hydroxyde sont ultra-minoritaires, ce qui est cohérent avec une solution acide. Le produit $2{,}0 \times 10^{-3} \times 5{,}0 \times 10^{-12} = 1{,}0 \times 10^{-14} = K_e$ : la relation est bien vérifiée.
\end{exemplebox}

\begin{exoresolu}[Concentrations et hiérarchie des espèces dans une solution de potasse]
Une solution aqueuse d'hydroxyde de potassium (potasse, \ce{KOH}) à \SI{25}{\celsius} a une concentration en ions hydroxyde $[\ce{OH-}] = 4{,}0 \times 10^{-2}~\si{mol.L^{-1}}$.
\begin{enumerate}
\item Calculer la concentration en ions oxonium.
\item Comparer $[\ce{H3O+}]$ et $[\ce{OH-}]$ et conclure sur le caractère de la solution.
\end{enumerate}
\tcblower
\textbf{1. Concentration en ions oxonium.} La solution est aqueuse, donc le produit ionique s'applique :
\[ [\ce{H3O+}] = \frac{K_e}{[\ce{OH-}]} = \frac{1{,}0 \times 10^{-14}}{4{,}0 \times 10^{-2}} = 2{,}5 \times 10^{-13}~\si{mol.L^{-1}} \]

\textbf{2. Comparaison et conclusion.}
\[ [\ce{OH-}] = 4{,}0 \times 10^{-2}~\si{mol.L^{-1}} \gg [\ce{H3O+}] = 2{,}5 \times 10^{-13}~\si{mol.L^{-1}} \]
Les ions hydroxyde sont largement majoritaires, les ions oxonium sont ultra-minoritaires : la solution est \textbf{basique}. Remarquons que les ions \ce{H3O+} ne sont pas absents : il y en a toujours, même dans une base concentrée. C'est toute la force du produit ionique.
\end{exoresolu}

\courssub{Influence de la température sur \texorpdfstring{$K_e$}{Ke}}

\textbf{L'autoprotolyse est endothermique}\par

Nous l'avons signalé : la réaction \ce{2 H2O <=> H3O+ + OH-} est \textbf{endothermique} : elle absorbe de la chaleur pour briser une liaison \ce{O-H}. D'après le principe de Le Chatelier (loi de modération), si l'on \textbf{élève la température}, l'équilibre se déplace dans le sens qui \emph{consomme} la chaleur apportée, c'est-à-dire le \textbf{sens direct} : davantage d'ions \ce{H3O+} et \ce{OH-} se forment.

\begin{propriete}[Variation de Ke avec la température]
L'autoprotolyse étant endothermique, $K_e$ \textbf{augmente} quand la température augmente (et donc $\mathrm{p}K_e$ \textbf{diminue}). Quelques valeurs :

\begin{center}
\begin{tabularx}{\linewidth}{|c|c|c|c|}
\hline
\rowcolor{popblueL} $\theta$ (\si{\celsius}) & $K_e$ & $\mathrm{p}K_e$ & $[\ce{H3O+}]$ dans l'eau pure (\si{mol.L^{-1}}) \\
\hline
0 & $1{,}1 \times 10^{-15}$ & $14{,}96$ & $3{,}3 \times 10^{-8}$ \\
\hline
25 & $1{,}0 \times 10^{-14}$ & $14{,}00$ & $1{,}0 \times 10^{-7}$ \\
\hline
50 & $5{,}5 \times 10^{-14}$ & $13{,}26$ & $2{,}3 \times 10^{-7}$ \\
\hline
100 & $5{,}1 \times 10^{-13}$ & $12{,}29$ & $7{,}1 \times 10^{-7}$ \\
\hline
\end{tabularx}
\end{center}
\end{propriete}

\begin{popfigure}
\begin{center}
\begin{tikzpicture}
\begin{axis}[
width=0.85\linewidth, height=6.2cm,
xlabel={Température $\theta$ (\si{\celsius})},
ylabel={$\mathrm{p}K_e$},
xmin=0, xmax=100, ymin=12, ymax=15.2,
grid=major, grid style={popline!40},
axis line style={popdark},
label style={popdark},
tick label style={font=\small},
]
\addplot[popcurve, domain=0:100, samples=100] {14.96 - 0.035*x + 0.000083*x*x};
\addplot[only marks, mark=*, mark size=2pt, poporange] coordinates {(0,14.96) (25,14.00) (50,13.26) (100,12.29)};
\node[pin={[font=\scriptsize]80:{$\mathrm{p}K_e = 14$ à \SI{25}{\celsius}}}] at (axis cs:25,14.00) {};
\node[pin={[font=\scriptsize]-60:{$\mathrm{p}K_e \approx 13{,}26$ à \SI{50}{\celsius}}}] at (axis cs:50,13.26) {};
\end{axis}
\end{tikzpicture}
\end{center}
\legende{Évolution de $\mathrm{p}K_e$ en fonction de la température : l'autoprotolyse étant endothermique, $K_e$ croît et $\mathrm{p}K_e$ décroît quand la température augmente.}
\end{popfigure}

\textbf{Conséquence capitale : la neutralité dépend de la température}\par

Dans l'eau pure, on a toujours $[\ce{H3O+}] = [\ce{OH-}]$ (la stœchiométrie l'impose), donc l'eau pure est \textbf{toujours neutre}, quelle que soit la température. Mais la \emph{valeur} de ces concentrations égales dépend de $K_e$ :
\[ [\ce{H3O+}] = [\ce{OH-}] = \sqrt{K_e} \]
Anticipant sur la section suivante (où $\mathrm{pH} = -\log[\ce{H3O+}]$), le pH de l'eau pure vaut donc :
\[ \mathrm{pH}_{\text{neutre}} = -\log\sqrt{K_e} = \frac{\mathrm{p}K_e}{2} \]
\begin{itemize}
\item À \SI{25}{\celsius} : $\mathrm{pH}_{\text{neutre}} = \dfrac{14{,}00}{2} = 7{,}00$ ;
\item À \SI{50}{\celsius} : $\mathrm{pH}_{\text{neutre}} = \dfrac{13{,}26}{2} \approx 6{,}6$.
\end{itemize}

\begin{attention}[Erreur fréquente : « neutre = pH 7 » ? Pas toujours !]
L'affirmation « une solution neutre a un pH de 7 » n'est vraie \textbf{qu'à \SI{25}{\celsius}}. À \SI{50}{\celsius}, l'eau pure a un pH d'environ $6{,}6$ tout en restant \textbf{parfaitement neutre}, car $[\ce{H3O+}] = [\ce{OH-}]$. Le vrai critère de neutralité n'est pas une valeur de pH, mais l'\textbf{égalité des concentrations} :
\[ \text{solution neutre} \iff [\ce{H3O+}] = [\ce{OH-}] \iff \mathrm{pH} = \frac{\mathrm{p}K_e}{2} \]
De même, à \SI{50}{\celsius}, une solution de pH $= 7{,}0$ est en réalité \emph{légèrement basique} (car $7{,}0 > 6{,}6$) ! Au Baccalauréat, vérifie toujours la température indiquée dans l'énoncé avant de conclure sur le caractère acide ou basique d'une solution.
\end{attention}

\begin{savaistu}[Pourquoi les analyses médicales précisent la température]
Lorsqu'un laboratoire d'analyses médicales à Douala ou Yaoundé mesure le pH d'un échantillon (sang, urine), l'appareil affiche la température de mesure et applique une \emph{correction de température}. En effet, l'électrode du pH-mètre et les équilibres ioniques dépendent tous deux de la température : un même échantillon mesuré à \SI{20}{\celsius} ou à \SI{37}{\celsius} (température du corps humain) ne donne pas exactement la même lecture. Le pH sanguin normal, compris entre $7{,}35$ et $7{,}45$ à \SI{37}{\celsius}, est d'ailleurs légèrement \emph{supérieur} au pH de neutralité à cette température ($\mathrm{p}K_e/2 \approx 6{,}8$) : le sang est donc très légèrement basique, et c'est vital !
\end{savaistu}

\courssub{Bilan de la section}

\begin{aretenir}[L'essentiel sur l'autoprotolyse et le produit ionique]
\begin{itemize}
\item L'eau pure conduit très faiblement le courant : elle contient des ions \ce{H3O+} et \ce{OH-} formés par \cle{autoprotolyse} :
\[ \ce{2 H2O(l) <=> H3O+(aq) + OH-(aq)} \quad \text{(équilibre très limité, endothermique)} \]
\item Le \cle{produit ionique de l'eau} : $K_e = [\ce{H3O+}] \times [\ce{OH-}]$, sans unité, ne dépend que de la température. À \SI{25}{\celsius} : $K_e = 1{,}0 \times 10^{-14}$ et $\mathrm{p}K_e = 14$.
\item Cette relation est vérifiée dans \textbf{toute solution aqueuse}, acide, neutre ou basique : $[\ce{OH-}] = \dfrac{K_e}{[\ce{H3O+}]}$.
\item Dans l'eau pure à \SI{25}{\celsius} : $[\ce{H3O+}] = [\ce{OH-}] = 1{,}0 \times 10^{-7}~\si{mol.L^{-1}}$.
\item $K_e$ augmente avec la température (réaction endothermique) : à \SI{50}{\celsius}, $K_e \approx 5{,}5 \times 10^{-14}$ et $\mathrm{p}K_e \approx 13{,}26$.
\item La neutralité correspond à $[\ce{H3O+}] = [\ce{OH-}]$, soit $\mathrm{pH} = \dfrac{\mathrm{p}K_e}{2}$ : elle vaut 7 \textbf{uniquement à \SI{25}{\celsius}}.
\end{itemize}
\end{aretenir}

\begin{exercice}[À toi de jouer]
\begin{enumerate}
\item À \SI{25}{\celsius}, une solution de vinaigre a une concentration $[\ce{H3O+}] = 1{,}6 \times 10^{-3}~\si{mol.L^{-1}}$. Calculer $[\ce{OH-}]$ et préciser si les ions hydroxyde sont majoritaires, minoritaires ou ultra-minoritaires.
\item À \SI{50}{\celsius}, on mesure dans une solution $[\ce{OH-}] = 2{,}0 \times 10^{-6}~\si{mol.L^{-1}}$. Calculer $[\ce{H3O+}]$ (on prendra $K_e = 5{,}5 \times 10^{-14}$). Cette solution est-elle acide, neutre ou basique ?
\item Montrer que dans l'eau pure à \SI{100}{\celsius} ($K_e = 5{,}1 \times 10^{-13}$), le pH vaut environ $6{,}15$. L'eau pure bouillante est-elle acide ? Justifier.
\end{enumerate}
\end{exercice}

\begin{corrige}[Exercice « À toi de jouer »]
\textbf{1.} $[\ce{OH-}] = \dfrac{K_e}{[\ce{H3O+}]} = \dfrac{1{,}0 \times 10^{-14}}{1{,}6 \times 10^{-3}} = 6{,}3 \times 10^{-12}~\si{mol.L^{-1}}$. Cette valeur est très inférieure à $10^{-7}~\si{mol.L^{-1}}$ : les ions \ce{OH-} sont \textbf{ultra-minoritaires} (solution acide).

\textbf{2.} $[\ce{H3O+}] = \dfrac{5{,}5 \times 10^{-14}}{2{,}0 \times 10^{-6}} = 2{,}75 \times 10^{-8}~\si{mol.L^{-1}}$. On compare : $[\ce{OH-}] = 2{,}0 \times 10^{-6} > [\ce{H3O+}] = 2{,}75 \times 10^{-8}$ : la solution est \textbf{basique}. (Attention : ne pas utiliser $K_e = 10^{-14}$ ici, la température est de \SI{50}{\celsius} !)

\textbf{3.} Dans l'eau pure : $[\ce{H3O+}] = \sqrt{K_e} = \sqrt{5{,}1 \times 10^{-13}} = 7{,}1 \times 10^{-7}~\si{mol.L^{-1}}$, d'où $\mathrm{pH} = -\log(7{,}1 \times 10^{-7}) \approx 6{,}15$. L'eau pure bouillante \textbf{n'est pas acide} : elle reste neutre car $[\ce{H3O+}] = [\ce{OH-}]$. Son pH est simplement inférieur à 7 parce que $\mathrm{p}K_e$ a diminué avec la température.
\end{corrige}

Nous disposons désormais de l'outil fondamental qui relie $[\ce{H3O+}]$ et $[\ce{OH-}]$ dans toute solution aqueuse. Mais manipuler des concentrations comme $2{,}5 \times 10^{-13}~\si{mol.L^{-1}}$ est peu pratique au quotidien\dots C'est précisément pour cela que les chimistes ont inventé une échelle plus commode : le \cle{pH}, que nous allons définir et apprendre à calculer dans la section suivante.

\coursec{Le pH : définition, calcul et caractère des solutions}

Nous venons de voir que toute solution aqueuse contient à la fois des ions \ce{H3O+} et des ions \ce{OH-}, liés par le produit ionique de l'eau $K_e = [\ce{H3O+}][\ce{OH-}]$. Mais ces concentrations sont extrêmement petites : de l'ordre de \SI{1,0e-7}{mol.L^{-1}} dans l'eau pure, et elles peuvent descendre jusqu'à \SI{1,0e-14}{mol.L^{-1}} ! Manier de tels nombres au quotidien — au laboratoire, à l'hôpital, dans une savonnerie artisanale — serait fastidieux. Les chimistes ont donc inventé une \emph{échelle pratique}, plus parlante : le \cle{pH}. C'est cette grandeur, omniprésente sur les résultats d'analyses médicales et les étiquettes des produits usuels, que nous allons maintenant apprivoiser.

\courssub{Définition du pH}

\textbf{Intuition.} Imagine que tu veuilles comparer l'acidité d'un jus de citron et celle d'un jus de bissap. Dire « le citron contient \SI{1,0e-2}{mol.L^{-1}} d'ions \ce{H3O+} et le bissap \SI{1,0e-3}{mol.L^{-1}} » est correct mais peu commode. Le pH compresse cette immense plage de concentrations (de \SI{1}{mol.L^{-1}} à \SI{1,0e-14}{mol.L^{-1}}) en une échelle simple, allant usuellement de 0 à 14.

\begin{definition}[Le pH d'une solution aqueuse]
Le \cle{pH} (potentiel hydrogène) d'une solution aqueuse est défini par :
\[ \mathrm{pH} = -\log\,[\ce{H3O+}] \]
où $[\ce{H3O+}]$ est la concentration molaire des ions oxonium, exprimée en \si{mol.L^{-1}}. Le pH est un nombre \textbf{sans unité}.

Relation inverse, tout aussi importante :
\[ [\ce{H3O+}] = 10^{-\mathrm{pH}} \ \si{mol.L^{-1}} \]

Conséquence immédiate : \textbf{plus le pH est faible, plus la concentration en \ce{H3O+} est grande, donc plus la solution est acide.}
\end{definition}

\begin{attention}[Le pH est sans unité, et c'est une échelle logarithmique]
Deux erreurs classiques au Baccalauréat :
\begin{itemize}
\item Écrire « pH = 2,3 mol/L » : \textbf{faux}, le pH n'a pas d'unité.
\item Croire que passer de pH 3 à pH 2 « enlève un peu d'acidité » : en réalité, la concentration en \ce{H3O+} est \textbf{multipliée par 10}. Une variation d'\textbf{une unité de pH correspond à un facteur 10} sur $[\ce{H3O+}]$ ; deux unités, un facteur 100. Le sang (pH $\approx$ 7,4) est donc dix fois moins concentré en \ce{H3O+} qu'une solution de pH 6,4.
\end{itemize}
\end{attention}

\courssub{L'échelle de pH et le caractère des solutions à \SI{25}{\celsius}}

Pour les solutions diluées courantes ($[\ce{H3O+}]$ comprise entre \SI{1}{} et \SI{1,0e-14}{mol.L^{-1}}), le pH varie entre \textbf{0 et 14}. À \SI{25}{\celsius}, où $pK_e = 14$, on distingue trois domaines :

\begin{aretenir}[Caractère d'une solution à \SI{25}{\celsius}]
\begin{itemize}
\item Solution \cle{acide} : $\mathrm{pH} < 7$, c'est-à-dire $[\ce{H3O+}] > [\ce{OH-}]$ ;
\item Solution \cle{neutre} : $\mathrm{pH} = 7$, c'est-à-dire $[\ce{H3O+}] = [\ce{OH-}] = \SI{1,0e-7}{mol.L^{-1}}$ ;
\item Solution \cle{basique} : $\mathrm{pH} > 7$, c'est-à-dire $[\ce{H3O+}] < [\ce{OH-}]$.
\end{itemize}
Attention : à une température différente de \SI{25}{\celsius}, le pH de neutralité n'est plus 7 mais $pK_e/2$ (voir section 3).
\end{aretenir}

\begin{popfigure}
\begin{center}
\begin{tikzpicture}[x=0.82cm]
% Zones colorées
\fill[poporangeL] (0,0) rectangle (7,0.7);
\fill[popgreenL] (7,0) rectangle (7.15,0.7);
\fill[popblueL] (7.15,0) rectangle (14,0.7);
% Graduations
\foreach \x in {0,1,...,14}{
  \draw[popdark] (\x,0) -- (\x,-0.15) node[below,font=\small] {\x};
}
\draw[thick,popdark] (0,0) rectangle (14,0.7);
% Labels des zones
\node[font=\bfseries,poporange!70!black] at (3.5,0.35) {ACIDE};
\node[font=\bfseries,popgreen!60!black] at (7.07,1.05) {NEUTRE};
\node[font=\bfseries,popblue!70!black] at (10.5,0.35) {BASIQUE};
% Exemples de produits usuels
\draw[->,popdark,thick] (2,1.6) -- (2,0.75);
\node[above,font=\small,align=center] at (2,1.6) {jus de citron\\ pH $\approx$ 2};
\draw[->,popdark,thick] (3,-1.1) -- (3,-0.35);
\node[below,font=\small,align=center] at (3,-1.1) {jus de bissap\\ pH $\approx$ 3};
\draw[->,popdark,thick] (5.5,1.6) -- (5.5,0.75);
\node[above,font=\small,align=center] at (5.5,1.6) {vinaigre\\ pH $\approx$ 2,5--3};
\draw[->,popdark,thick] (7,-1.1) -- (7,-0.35);
\node[below,font=\small,align=center] at (7,-1.1) {eau minérale\\ pH $\approx$ 7};
\draw[->,popdark,thick] (10,1.6) -- (10,0.75);
\node[above,font=\small,align=center] at (10,1.6) {savon\\ pH $\approx$ 10};
\draw[->,popdark,thick] (13,-1.1) -- (13,-0.35);
\node[below,font=\small,align=center] at (13,-1.1) {potasse\\ pH $\approx$ 13};
\end{tikzpicture}
\end{center}
\legende{Frise de l'échelle de pH à \SI{25}{\celsius} : zones acide, neutre et basique, avec quelques produits usuels du quotidien camerounais.}
\end{popfigure}

\begin{savaistu}[Le pH autour de nous]
Le jus de citron (pH $\approx$ 2) doit son acidité à l'acide citrique ; le jus de bissap, préparé à partir des fleurs d'hibiscus très consommées au Cameroun, a un pH voisin de 3 grâce aux acides organiques qu'il contient — c'est d'ailleurs cette acidité qui donne sa belle couleur rouge aux anthocyanes de la fleur ! À l'autre bout de l'échelle, la potasse traditionnelle, utilisée en savonnerie artisanale et en cuisine, atteint pH $\approx$ 13 : c'est une solution fortement basique, à manipuler avec précaution.
\end{savaistu}

\courssub{La relation complémentaire : pH + pOH = pKe}

Par analogie avec le pH, on définit le \cle{pOH} : $\mathrm{pOH} = -\log\,[\ce{OH-}]$. En prenant l'opposé du logarithme des deux membres de $K_e = [\ce{H3O+}][\ce{OH-}]$, on obtient une relation très utile :

\begin{propriete}[Relation entre pH et pOH]
\[ \mathrm{pH} + \mathrm{pOH} = pK_e \]
À \SI{25}{\celsius} : $\mathrm{pH} + \mathrm{pOH} = 14$.
\end{propriete}

Cette relation permet de passer instantanément de l'acidité à la basicité : une solution de pH 3 a un pOH de 11, donc $[\ce{OH-}] = \SI{1,0e-11}{mol.L^{-1}}$.

\begin{popfigure}
\begin{center}
\begin{tikzpicture}
\begin{axis}[
  width=0.85\linewidth, height=6.5cm,
  xlabel={pH}, ylabel={$\log c$ (mol.L$^{-1}$)},
  xmin=0, xmax=14, ymin=-14.5, ymax=-0.5,
  xtick={0,2,4,6,7,8,10,12,14},
  ytick={-14,-12,-10,-8,-6,-4,-2,0},
  grid=both, grid style={popline!40},
  legend style={at={(0.02,0.05)},anchor=south west,font=\small},
  axis line style={popdark},
]
\addplot[popcurve,poporange,domain=0:14,samples=100] {-x};
\addlegendentry{$\log[\ce{H3O+}]$}
\addplot[popcurve,popblue,domain=0:14,samples=100] {x-14};
\addlegendentry{$\log[\ce{OH-}]$}
\draw[dashed,popdark] (axis cs:7,-14) -- (axis cs:7,0);
\node[above,font=\small,popdark] at (axis cs:7,0) {pH = 7};
\fill[popgreen!70!black] (axis cs:7,-7) circle (2pt);
\node[right,font=\small,popgreen!50!black] at (axis cs:7.2,-7) {$[\ce{H3O+}]=[\ce{OH-}]=10^{-7}$};
\end{axis}
\end{tikzpicture}
\end{center}
\legende{Évolution de $\log[\ce{H3O+}]$ et $\log[\ce{OH-}]$ en fonction du pH à \SI{25}{\celsius}. Les deux droites se coupent à pH = 7 : c'est la neutralité.}
\end{popfigure}

\courssub{Calcul du pH d'une solution d'acide fort}

Un \cle{acide fort} (comme \ce{HCl}, l'acide chlorhydrique) réagit \emph{totalement} avec l'eau :
\[ \ce{HCl + H2O -> H3O+ + Cl-} \]
Cette transformation est totale (réaction quantitative, quasi instantanée) : chaque mole de \ce{HCl} apportée produit une mole d'ions \ce{H3O+}. Pour une solution de concentration apportée $C$, on a donc $[\ce{H3O+}] = C$.

\begin{methode}[Calculer le pH d'une solution d'acide fort]
\begin{enumerate}
\item \textbf{Vérifier le domaine de validité} : la formule suppose que les ions \ce{H3O+} venant de l'autoprotolyse de l'eau (\SI{1,0e-7}{mol.L^{-1}}) sont négligeables devant $C$, et que la solution reste diluée. On exige donc $\SI{1e-6}{mol.L^{-1}} < C < \SI{1e-1}{mol.L^{-1}}$.
\item \textbf{Écrire} $[\ce{H3O+}] = C$.
\item \textbf{Conclure} : $\mathrm{pH} = -\log C$.
\end{enumerate}
\end{methode}

\begin{exemplebox}[Solution d'acide chlorhydrique]
Soit une solution de \ce{HCl} de concentration $C = \SI{5,0e-3}{mol.L^{-1}}$.

Domaine de validité : $\SI{1e-6}{} < \SI{5,0e-3}{} < \SI{1e-1}{}$ : vérifié.

\[ \mathrm{pH} = -\log(5{,}0\times 10^{-3}) = 3 - \log 5{,}0 \approx 3 - 0{,}70 = 2{,}3 \]

pH $= 2{,}3 < 7$ : la solution est bien \textbf{acide}. Vérification inverse : $[\ce{H3O+}] = 10^{-2{,}3} \approx \SI{5,0e-3}{mol.L^{-1}}$ : cohérent.
\end{exemplebox}

\begin{attention}[La formule pH = $-\log C$ n'est valable QUE pour les acides forts]
Erreur très fréquente au Bac : appliquer $\mathrm{pH} = -\log C$ à un acide \emph{faible} comme l'acide éthanoïque \ce{CH3COOH} (celui du vinaigre). Un acide faible ne réagit que \emph{partiellement} avec l'eau :
\[ \ce{CH3COOH + H2O <=> CH3COO- + H3O+} \]
donc $[\ce{H3O+}] < C$ et le pH réel est \textbf{supérieur} à $-\log C$. Une solution d'acide éthanoïque à \SI{1,0e-2}{mol.L^{-1}} a un pH voisin de 3,4 et non de 2. Nous étudierons ce calcul en détail dans la leçon sur les acides faibles.
\end{attention}

\courssub{Calcul du pH d'une solution de base forte}

Une \cle{base forte} (comme l'hydroxyde de sodium \ce{NaOH}, la soude) se dissocie totalement dans l'eau :
\[ \ce{NaOH ->[H_2O] Na+ + OH-} \]
Pour une solution de concentration $C$, on a donc $[\ce{OH-}] = C$. On remonte alors aux ions oxonium grâce au produit ionique :
\[ [\ce{H3O+}] = \frac{K_e}{[\ce{OH-}]} = \frac{K_e}{C} \]
d'où, en passant aux logarithmes à \SI{25}{\celsius} :
\[ \mathrm{pH} = pK_e + \log C = 14 + \log C \]

\begin{methode}[Calculer le pH d'une solution de base forte]
\begin{enumerate}
\item Vérifier le domaine de validité : $\SI{1e-6}{mol.L^{-1}} < C < \SI{1e-1}{mol.L^{-1}}$.
\item Écrire $[\ce{OH-}] = C$.
\item Conclure : $\mathrm{pH} = pK_e + \log C$, soit $\mathrm{pH} = 14 + \log C$ à \SI{25}{\celsius}.
\end{enumerate}
\end{methode}

\begin{exemplebox}[Solution de soude]
Solution de \ce{NaOH} à $C = \SI{2,0e-2}{mol.L^{-1}}$ :
\[ \mathrm{pH} = 14 + \log(2{,}0\times 10^{-2}) = 14 - 2 + \log 2{,}0 \approx 12 + 0{,}30 = 12{,}3 \]
pH $= 12{,}3 > 7$ : solution fortement \textbf{basique}, comme la lessive de soude utilisée en savonnerie.
\end{exemplebox}

\courssub{Retrouver les concentrations à partir du pH ; hiérarchie des espèces}

Démarche inverse, très demandée au Baccalauréat : on mesure le pH, et on veut les concentrations de toutes les espèces.

\begin{methode}[Des concentrations à partir du pH]
\begin{enumerate}
\item Calculer $[\ce{H3O+}] = 10^{-\mathrm{pH}}$.
\item En déduire $[\ce{OH-}] = \dfrac{K_e}{[\ce{H3O+}]} = 10^{\mathrm{pH}-pK_e}$.
\item \textbf{Classer les espèces} : l'espèce de concentration la plus élevée est \cle{majoritaire} ; celle qui est nettement plus faible est \cle{minoritaire} ; celle qui est extrêmement faible (souvent \ce{OH-} en milieu acide, ou \ce{H3O+} en milieu basique) est dite \cle{ultra-minoritaire}.
\item \textbf{Vérifier l'électroneutralité} : la somme des charges positives égale la somme des charges négatives. C'est un excellent contrôle de cohérence.
\end{enumerate}
\end{methode}

\begin{exoresolu}[Analyse complète d'une solution acide]
Une solution aqueuse de chlorure de sodium contenant de l'acide chlorhydrique a un pH de 3,0 à \SI{25}{\celsius}, avec $[\ce{Na+}] = \SI{1,0e-2}{mol.L^{-1}}$. Déterminer $[\ce{H3O+}]$, $[\ce{OH-}]$ et $[\ce{Cl-}]$, puis classer les espèces et vérifier l'électroneutralité.
\tcblower
\textbf{Étape 1 :} $[\ce{H3O+}] = 10^{-3{,}0} = \SI{1,0e-3}{mol.L^{-1}}$.

\textbf{Étape 2 :} $[\ce{OH-}] = \dfrac{K_e}{[\ce{H3O+}]} = \dfrac{10^{-14}}{10^{-3}} = \SI{1,0e-11}{mol.L^{-1}}$.

\textbf{Étape 3 : électroneutralité.} Les ions présents sont \ce{Na+}, \ce{H3O+}, \ce{Cl-}, \ce{OH-} :
\[ [\ce{Cl-}] + [\ce{OH-}] = [\ce{Na+}] + [\ce{H3O+}] \]
\[ [\ce{Cl-}] = \SI{1,0e-2}{} + \SI{1,0e-3}{} - \SI{1,0e-11}{} \approx \SI{1,1e-2}{mol.L^{-1}} \]

\textbf{Étape 4 : hiérarchie.}
\begin{itemize}
\item Espèces \textbf{majoritaires} : \ce{Cl-} (\SI{1,1e-2}{}) et \ce{Na+} (\SI{1,0e-2}{mol.L^{-1}}) ;
\item Espèce \textbf{minoritaire} : \ce{H3O+} (\SI{1,0e-3}{mol.L^{-1}}), dix fois moins concentrée ;
\item Espèce \textbf{ultra-minoritaire} : \ce{OH-} (\SI{1,0e-11}{mol.L^{-1}}), un milliard de fois moins concentrée que \ce{Na+} !
\end{itemize}
Cohérence : charges positives $= 10^{-2} + 10^{-3} = \SI{1,1e-2}{}$ ; charges négatives $= 1{,}1\times10^{-2} + 10^{-11} \approx \SI{1,1e-2}{mol.L^{-1}}$. L'électroneutralité est vérifiée.
\end{exoresolu}

\begin{attention}[Ne jamais oublier les ions \ce{OH-}... même en milieu acide]
Dire « il n'y a pas d'ions \ce{OH-} dans une solution acide » est \textbf{faux} : le produit ionique impose leur présence, même en quantité infime. À pH 2, $[\ce{OH-}] = \SI{1,0e-12}{mol.L^{-1}}$ : ultra-minoritaires, mais présents. Toute solution aqueuse contient \emph{toujours} \ce{H3O+} et \ce{OH-}.
\end{attention}

\courssub{Effet de la dilution sur le pH}

Diluer dix fois une solution d'acide fort revient à diviser $C$ par 10, donc $[\ce{H3O+}]$ par 10 : le pH \textbf{augmente d'une unité}. De même, diluer dix fois une base forte fait \textbf{diminuer} le pH d'une unité.

\begin{exemplebox}[Dilutions successives d'un acide fort]
\begin{center}
\begin{tabularx}{\linewidth}{|l|X|X|}
\hline
\rowcolor{popblueL} Concentration $C$ (\si{mol.L^{-1}}) & $[\ce{H3O+}]$ (\si{mol.L^{-1}}) & pH \\
\hline
$10^{-1}$ & $10^{-1}$ & 1,0 \\
\hline
$10^{-2}$ & $10^{-2}$ & 2,0 \\
\hline
$10^{-3}$ & $10^{-3}$ & 3,0 \\
\hline
$10^{-4}$ & $10^{-4}$ & 4,0 \\
\hline
\end{tabularx}
\end{center}
Chaque dilution au dixième ajoute une unité de pH. Mais attention : cette règle \textbf{cesse d'être valable} près de la neutralité ! Diluer une solution à $C = \SI{1e-8}{mol.L^{-1}}$ ne donne pas pH 8 (une solution d'acide ne peut pas devenir basique par dilution !) : les ions \ce{H3O+} de l'autoprotolyse de l'eau ne sont plus négligeables, et le pH tend simplement vers 7. D'où l'importance du domaine de validité $C > \SI{1e-6}{mol.L^{-1}}$.
\end{exemplebox}

\begin{aretenir}[L'essentiel de la section]
\begin{itemize}
\item $\mathrm{pH} = -\log[\ce{H3O+}]$ et $[\ce{H3O+}] = 10^{-\mathrm{pH}}$ ; le pH est sans unité.
\item À \SI{25}{\celsius} : acide si pH $< 7$, neutre si pH $= 7$, basique si pH $> 7$.
\item $\mathrm{pH} + \mathrm{pOH} = pK_e = 14$ à \SI{25}{\celsius}.
\item Acide fort : $\mathrm{pH} = -\log C$ ; base forte : $\mathrm{pH} = 14 + \log C$ (si $\SI{1e-6}{} < C < \SI{1e-1}{mol.L^{-1}}$).
\item Une unité de pH = un facteur 10 sur $[\ce{H3O+}]$ ; diluer dix fois un acide fort élève le pH de 1.
\item Toute solution contient \ce{H3O+} et \ce{OH-} : pense à classer les espèces (majoritaire, minoritaire, ultra-minoritaire) et à vérifier l'électroneutralité.
\end{itemize}
\end{aretenir}

Maintenant que tu sais \emph{calculer} un pH, une question naturelle se pose : comment le \emph{mesure}-t-on concrètement au laboratoire ou à l'hôpital ? C'est l'objet de la section suivante, consacrée au papier pH, au pH-mètre et aux indicateurs colorés — y compris celui que tu peux préparer toi-même avec des fleurs de bissap.

\coursec{Mesure du pH et indicateurs colorés}

Dans la section précédente, tu as appris à \emph{calculer} le pH d'une solution à partir de la concentration en ions \ce{H3O+} : c'est l'approche théorique. Mais au laboratoire, à l'hôpital ou dans une usine comme la SONARA, personne ne connaît à l'avance la concentration en ions oxonium du liquide qu'on tient entre les mains ! Il faut donc \emph{mesurer} le pH. Comment le technicien du laboratoire d'analyses médicales de l'Hôpital Central de Yaoundé vérifie-t-il le pH d'un échantillon ? Comment le brasseur contrôle-t-il l'acidité de son moût ? C'est toute la question de cette section : nous allons découvrir les trois grandes méthodes de mesure du pH — le papier pH, les indicateurs colorés et le pH-mètre — puis apprendre à fabriquer nous-mêmes un indicateur coloré naturel à partir de fleurs de bissap.

\courssub{Vue d'ensemble : trois méthodes, trois niveaux de précision}

Il existe trois méthodes usuelles pour déterminer le pH d'une solution aqueuse. Elles ne se valent pas : elles diffèrent par leur principe, leur coût, leur rapidité et surtout leur \cle{précision}.

\begin{center}
\begin{tabularx}{\linewidth}{|l|X|X|c|}
\hline
\rowcolor{popblueL} \textbf{Méthode} & \textbf{Principe} & \textbf{Avantage principal} & \textbf{Précision} \\
\hline
Papier pH & Mélange d'indicateurs imprégné sur papier, comparé à une échelle étalon & Rapide, bon marché & 0,5 à 1 unité \\
\hline
Indicateurs colorés & Couple acide/base dont les deux formes ont des couleurs différentes & Encadrement rapide du pH & 1 à 2 unités \\
\hline
pH-mètre & Électrode de verre mesurant une tension électrique liée à $[\ce{H3O+}]$ & Très précis, valeur numérique & 0,01 à 0,1 unité \\
\hline
\end{tabularx}
\end{center}

\begin{aretenir}[L'idée directrice]
Le choix de la méthode dépend du besoin : pour vérifier rapidement qu'une piscine n'est pas trop acide, le papier pH suffit ; pour un dosage acido-basique au baccalauréat ou un dosage sanguin, seul le \cle{pH-mètre} donne une valeur fiable à $0,05$ unité près.
\end{aretenir}

\courssub{Le papier pH : la mesure rapide}

\textbf{Principe.} Le papier pH est une bandelette de papier poreux imprégnée d'un \emph{mélange} de plusieurs indicateurs colorés. Comme chaque indicateur change de couleur dans une zone de pH différente, le mélange prend une teinte globale qui varie continûment du rouge (milieu très acide) au bleu-violet (milieu très basique), en passant par l'orange, le jaune et le vert. Il suffit alors de comparer la couleur obtenue à l'\cle{échelle étalon} imprimée sur la boîte : chaque teinte correspond à une valeur entière (ou demi-entière) de pH.

\begin{methode}[Mesurer un pH au papier pH]
\begin{enumerate}
\item À l'aide d'un agitateur en verre propre, \textbf{déposer une goutte} de la solution à étudier sur un petit morceau de papier pH posé sur une coupelle.
\item \textbf{Ne jamais tremper} le papier directement dans le flacon : cela contaminerait toute la solution et détremperait le papier.
\item \textbf{Comparer immédiatement} la couleur obtenue à l'échelle étalon (certains papiers virent en quelques secondes puis continuent d'évoluer à l'air).
\item Lire la valeur de pH correspondant à la teinte la plus proche.
\end{enumerate}
\end{methode}

\textbf{Précision.} La lecture se fait « à l'œil » entre deux teintes voisines : on obtient le pH à $0,5$ ou $1$ unité près. Dire « le pH vaut environ 4 » est honnête ; écrire « pH $= 4,37$ » avec du papier pH serait absurde.

\begin{attention}[Piège classique]
Une erreur fréquente consiste à \emph{mouiller} le papier pH avec de l'eau distillée avant de déposer la goutte, « pour que ça diffuse mieux ». Cela dilue la solution testée et fausse la mesure, surtout pour les solutions peu concentrées. Le papier pH s'utilise \textbf{sec}.
\end{attention}

\courssub{Les indicateurs colorés : des molécules qui changent de couleur}

\textbf{Intuition.} Certaines molécules ont une propriété remarquable : leur forme acide et leur forme basique n'ont \emph{pas la même couleur}. En versant quelques gouttes d'une telle molécule dans une solution, la couleur observée révèle dans quel état — acide ou basique — elle se trouve majoritairement, donc révèle le pH du milieu. C'est le principe des indicateurs colorés, déjà utilisés au XVII\textsuperscript{e} siècle par Robert Boyle avec des extraits de violettes !

\begin{definition}[Indicateur coloré acido-basique]
Un \cle{indicateur coloré} acido-basique est un couple acide/base, noté $\ce{HInd}/\ce{Ind-}$, dont les deux formes conjuguées ont des \textbf{couleurs différentes}. L'équilibre mis en jeu est :
\[ \ce{HInd(aq) + H2O(l) <=> Ind^-(aq) + H3O+(aq)} \]
Il est caractérisé par sa \cle{zone de virage}, intervalle de pH (en général centré sur le $\mathrm{p}K_a$ du couple, soit $\mathrm{p}K_a \pm 1$) sur lequel la couleur passe progressivement de la teinte acide à la teinte basique.
\end{definition}

\textbf{Comment ça marche ?} En milieu très acide ($\mathrm{pH} < \mathrm{p}K_a - 1$), la forme \ce{HInd} domine : la solution prend la \emph{teinte acide}. En milieu très basique ($\mathrm{pH} > \mathrm{p}K_a + 1$), la forme \ce{Ind-} domine : on observe la \emph{teinte basique}. Entre les deux, les deux formes coexistent en proportions comparables : on observe une couleur intermédiaire, la \emph{teinte sensible} (mélange des deux couleurs).

\begin{center}
\begin{tabularx}{\linewidth}{|l|X|X|X|c|}
\hline
\rowcolor{popblueL} \textbf{Indicateur} & \textbf{Teinte acide} & \textbf{Zone de virage} & \textbf{Teinte basique} & $\mathrm{p}K_a$ \\
\hline
Hélianthine (méthylorange) & rouge & 3,1 -- 4,4 & jaune orangé & 3,7 \\
\hline
Bleu de bromothymol (BBT) & jaune & 6,0 -- 7,6 & bleu & 7,1 \\
\hline
Phénolphtaléine & incolore & 8,2 -- 10,0 & rose (fuchsia) & 9,2 \\
\hline
\end{tabularx}
\end{center}

\begin{popfigure}
\begin{tikzpicture}[font=\small, scale=0.95]
% Axe pH
\draw[->,thick,popdark] (0,0) -- (14.5,0) node[right]{pH};
\foreach \x/\l in {0/0,1/1,2/2,3/3,4/4,5/5,6/6,7/7,8/8,9/9,10/10,11/11,12/12,13/13,14/14}{
  \draw (\x,0.08) -- (\x,-0.08) node[below]{\l};
}
% Hélianthine : rouge 0-3.1, virage 3.1-4.4, jaune 4.4-14
\fill[red!70] (0,0.6) rectangle (3.1,1.2);
\shade[left color=red!70,right color=yellow!80] (3.1,0.6) rectangle (4.4,1.2);
\fill[yellow!80] (4.4,0.6) rectangle (14,1.2);
\node[left] at (0,0.9) {\textbf{Hélianthine}};
% BBT : jaune 0-6, virage 6-7.6, bleu 7.6-14
\fill[yellow!85] (0,1.7) rectangle (6,2.3);
\shade[left color=yellow!85,right color=blue!70] (6,1.7) rectangle (7.6,2.3);
\fill[blue!70] (7.6,1.7) rectangle (14,2.3);
\node[left] at (0,2.0) {\textbf{BBT}};
% Phénolphtaléine : incolore 0-8.2, virage 8.2-10, rose 10-14
\fill[gray!10,draw=gray!40] (0,2.8) rectangle (8.2,3.4);
\shade[left color=gray!10,right color=poppink!80] (8.2,2.8) rectangle (10,3.4);
\fill[poppink!80] (10,2.8) rectangle (14,3.4);
\node[left] at (0,3.1) {\textbf{Phénolphtaléine}};
% Annotations zones de virage
\draw[<->,popink] (3.1,-0.7) -- (4.4,-0.7) node[midway,below,popink]{\scriptsize 3,1 -- 4,4};
\draw[<->,popgreen!60!black] (6,-1.5) -- (7.6,-1.5) node[midway,below,popgreen!60!black]{\scriptsize 6,0 -- 7,6};
\draw[<->,poppurple] (8.2,-2.3) -- (10,-2.3) node[midway,below,poppurple]{\scriptsize 8,2 -- 10,0};
\end{tikzpicture}
\legende{Zones de virage des trois indicateurs usuels sur l'axe des pH à \SI{25}{\celsius}. Le dégradé marque la transition progressive entre la teinte acide et la teinte basique.}
\end{popfigure}

\begin{exemplebox}[Exploiter la couleur observée]
On ajoute quelques gouttes de BBT dans trois solutions A, B et C :
\begin{itemize}
\item A devient \textbf{jaune} : teinte acide, donc $\mathrm{pH}_A < 6{,}0$ ;
\item B devient \textbf{bleue} : teinte basique, donc $\mathrm{pH}_B > 7{,}6$ ;
\item C devient \textbf{verte} (mélange jaune + bleu) : teinte sensible, donc $6{,}0 < \mathrm{pH}_C < 7{,}6$.
\end{itemize}
Remarque bien que l'indicateur \emph{encadre} le pH sans le donner précisément : c'est sa limite, mais aussi sa force (une seule goutte suffit pour un encadrement immédiat).
\end{exemplebox}

\begin{attention}[Erreur fréquente : mal choisir son indicateur]
La \textbf{phénolphtaléine} vire entre 8,2 et 10 : elle reste \emph{incolore} dans tout milieu acide \textbf{et} neutre ! L'utiliser pour étudier un acide fort (pH voisin de 1 ou 2) ne donne \emph{aucune} information utile : la solution reste incolore, exactement comme de l'eau pure. Règle pratique : choisir un indicateur dont la zone de virage \emph{contient} le pH attendu — pour un milieu acide, préférer l'hélianthine ; pour un milieu proche de la neutralité, le BBT.
\end{attention}

\courssub{Le pH-mètre : la mesure de précision}

\textbf{Principe.} Le pH-mètre est un appareil électronique muni d'une \cle{électrode de verre} (souvent combinée avec une électrode de référence dans une seule sonde). La membrane de verre, très fine, est le siège d'une différence de potentiel qui dépend de la concentration en ions \ce{H3O+} de la solution dans laquelle elle plonge. L'appareil mesure cette tension électrique et la convertit directement en une valeur de pH affichée à l'écran : la relation entre tension et pH est affine, de l'ordre de $-59$ mV par unité de pH à \SI{25}{\celsius}.

\begin{popfigure}
\begin{tikzpicture}[font=\small, scale=0.9]
% Bécher
\draw[thick,popdark] (0,0) -- (0,3.2);
\draw[thick,popdark] (3.4,0) -- (3.4,3.2);
\draw[thick,popdark] (0,0) -- (3.4,0);
% Liquide
\fill[popblue!25] (0.06,0.06) rectangle (3.34,2.1);
\draw[popblue] (0.06,2.1) -- (3.34,2.1);
\node[popdark] at (1.7,0.45) {\scriptsize solution à étudier};
% Électrode (sonde)
\fill[gray!30] (1.45,1.0) rectangle (1.95,4.3);
\draw[thick,popdark] (1.45,1.0) rectangle (1.95,4.3);
\fill[popgold!60] (1.5,1.0) rectangle (1.9,1.7);
\draw[thick,popdark] (1.5,1.0) rectangle (1.9,1.7);
\node[right,align=left] at (2.1,1.35) {\scriptsize membrane de verre};
\draw[->,thin,popdark] (2.1,1.55) -- (1.95,1.4);
% Bulbe
\draw[thick,popdark,fill=popgold!40] (1.7,1.0) circle (0.22);
% Câble
\draw[thick,popdark] (1.7,4.3) -- (1.7,4.9) -- (5.2,4.9) -- (5.2,4.35);
% Boîtier pH-mètre
\draw[thick,popdark,fill=popcream] (4.3,3.0) rectangle (7.6,4.35);
\draw[fill=white,draw=popdark] (4.6,3.7) rectangle (7.3,4.15);
\node[popink] at (5.95,3.92) {\textbf{pH = 4,52}};
\node[popmuted] at (5.95,3.3) {\scriptsize pH-mètre};
\node[popmuted] at (5.95,2.7) {\scriptsize (précision 0,01)};
\end{tikzpicture}
\legende{Schéma simplifié d'un pH-mètre : la sonde (électrode de verre) plongée dans la solution mesure une tension électrique proportionnelle à l'acidité ; le boîtier la convertit en pH affiché.}
\end{popfigure}

\begin{methode}[Mesurer un pH au pH-mètre]
\begin{enumerate}
\item \textbf{Étalonner} l'appareil : plonger l'électrode successivement dans deux \cle{solutions tampons} de pH connus (par exemple pH $= 7{,}00$ puis pH $= 4{,}00$) et valider l'affichage de ces valeurs. Sans étalonnage, la mesure n'a aucune valeur !
\item \textbf{Rincer} l'électrode à l'eau distillée et l'essuyer délicatement (papier Joseph) entre chaque solution.
\item \textbf{Plonger} l'électrode dans la solution à étudier, agiter doucement.
\item \textbf{Attendre la stabilisation} de l'affichage, puis lire le pH.
\end{enumerate}
\end{methode}

\textbf{Précision.} Un pH-mètre de laboratoire bien étalonné donne le pH à $0,01$ ou $0,05$ unité près ; un appareil portable, à $0,1$ unité près. C'est la seule méthode qui permette de suivre l'évolution du pH au cours d'un dosage (nous l'utiliserons dans les leçons suivantes).

\begin{attention}[Sécurité et bon sens]
L'extrémité de l'électrode est une membrane de verre \emph{extrêmement fragile} : elle coûte cher et se brise au moindre choc contre le fond du bécher. On la manipule avec précaution et on la conserve plongée dans une solution de stockage, jamais à l'air sec.
\end{attention}

\courssub{Les indicateurs naturels : la chimie du jardin et de la cuisine}

\textbf{Intuition.} Tu as peut-être déjà remarqué qu'en pressant du citron dans un jus de bissap, la boisson devient d'un rouge plus vif et plus éclatant. Ce n'est pas un hasard : les fleurs d'hibiscus contiennent des \cle{anthocyanes}, des pigments naturels dont la couleur dépend du pH. Le citron acidifie le milieu, les anthocyanes passent à leur forme acide rouge vif : le bissap est un indicateur coloré naturel !

De nombreux végétaux contiennent de tels pigments : le chou rouge (rouge en milieu acide, violet au neutre, bleu puis vert en milieu basique), le curcuma (jaune en milieu acide, brun-rouge en milieu basique), les pétales de rose, de pétunia, le thé, le jus de betterave\ldots

\begin{popfigure}
\begin{tikzpicture}[font=\small, scale=0.95]
\draw[->,thick,popdark] (0,0) -- (14.5,0) node[right]{pH};
\foreach \x in {0,2,4,6,8,10,12,14}{
  \draw (\x,0.08) -- (\x,-0.08) node[below]{\x};
}
\shade[left color=red!75,right color=red!60!poppurple] (0,0.5) rectangle (4,1.3);
\shade[left color=red!60!poppurple,right color=poppurple!70] (4,0.5) rectangle (7,1.3);
\shade[left color=poppurple!70,right color=blue!70] (7,0.5) rectangle (10,1.3);
\shade[left color=blue!70,right color=popgreen!70] (10,0.5) rectangle (14,1.3);
\node[white] at (2,0.9) {\scriptsize rouge};
\node[white] at (5.5,0.9) {\scriptsize violet};
\node[white] at (8.5,0.9) {\scriptsize bleu};
\node[white] at (12,0.9) {\scriptsize vert};
\node[popdark,align=left] at (2,1.9) {\scriptsize acide};
\node[popdark,align=left] at (7,1.9) {\scriptsize neutre};
\node[popdark,align=left] at (12,1.9) {\scriptsize basique};
\end{tikzpicture}
\legende{Dégradé des couleurs du jus de chou rouge en fonction du pH : rouge en milieu acide, violet vers la neutralité, bleu puis vert en milieu basique.}
\end{popfigure}

\begin{savaistu}[Le bissap, indicateur du terroir]
Les fleurs d'hibiscus (\emph{Hibiscus sabdariffa}), appelées \emph{foléré} au Nord-Cameroun, donnent le jus de bissap. Leurs anthocyanes virent au rouge vif en milieu acide : c'est pourquoi un zeste de citron ou un morceau de tamarin « réveille » la couleur de la boisson. En milieu basique, le même jus tirerait vers le vert-bleu. Les vendeuses de bissap pratiquent donc, sans le savoir, de la chimie acido-basique au quotidien !
\end{savaistu}

\begin{experience}[Préparer un indicateur coloré naturel (chou rouge ou bissap)]
\textbf{Protocole.}
\begin{enumerate}
\item Hacher finement quelques feuilles de chou rouge (ou prendre une poignée de fleurs séchées d'hibiscus).
\item Les faire \textbf{macérer} dans de l'eau chaude (\SI{60}{\celsius} environ) pendant 10 à 15 minutes — on peut aussi utiliser un peu d'alcool, qui dissout mieux les anthocyanes.
\item \textbf{Filtrer} le mélange (entonnoir + papier filtre, ou simple tamis propre).
\item Recueillir le \textbf{filtrat} coloré : c'est l'indicateur. Le conserver au frais, dans un flacon fermé, à l'abri de la lumière (les anthocyanes se dégradent en quelques jours).
\end{enumerate}
\textbf{Test.} Verser quelques gouttes de l'indicateur dans du vinaigre (acide), de l'eau pure (neutre) et une solution de bicarbonate de sodium ou d'eau de Javel diluée (basique) : observer trois couleurs nettement différentes.
\textbf{Interprétation.} Les anthocyanes extraites changent de structure selon qu'elles cèdent ou captent des ions \ce{H+} : ce sont bien des couples acide/base dont les deux formes ont des couleurs différentes — la définition même d'un indicateur coloré.
\end{experience}

\begin{exoresolu}[Choisir le bon indicateur et la bonne méthode]
Une infirmière d'un centre de santé de Bafoussam dispose de trois solutions incolores notées S\textsubscript{1}, S\textsubscript{2}, S\textsubscript{3} : du vinaigre dilué, de l'eau de source et une solution d'eau de Javel diluée. Elle teste chaque solution avec du BBT et obtient : S\textsubscript{1} jaune, S\textsubscript{2} verte, S\textsubscript{3} bleue.
\begin{enumerate}
\item Attribuer à chaque solution un encadrement de pH et l'identifier.
\item Pour connaître précisément le pH de l'eau de source, quelle méthode choisir et pourquoi ?
\end{enumerate}
\tcblower
\textbf{1.} D'après la zone de virage du BBT (jaune 6,0 -- 7,6 bleu) :
\begin{itemize}
\item S\textsubscript{1} est jaune $\Rightarrow \mathrm{pH} < 6{,}0$ : milieu acide, c'est le \textbf{vinaigre dilué} ;
\item S\textsubscript{2} est verte $\Rightarrow 6{,}0 < \mathrm{pH} < 7{,}6$ : milieu voisin de la neutralité, c'est l'\textbf{eau de source} ;
\item S\textsubscript{3} est bleue $\Rightarrow \mathrm{pH} > 7{,}6$ : milieu basique, c'est l'\textbf{eau de Javel diluée}.
\end{itemize}
\textbf{2.} Le BBT ne donne qu'un encadrement. Pour une valeur précise (par exemple distinguer pH $= 6{,}8$ de pH $= 7{,}3$), il faut utiliser un \textbf{pH-mètre préalablement étalonné} avec des solutions tampons : sa précision ($0,01$ à $0,1$ unité) est très supérieure à celle des indicateurs colorés et du papier pH.
\end{exoresolu}

\begin{exercice}[À toi de jouer]
On dispose de trois indicateurs : hélianthine, BBT et phénolphtaléine. Une solution de jus de citron a un pH voisin de 2,5 ; une solution de lessive de soude diluée a un pH voisin de 11.
\begin{enumerate}
\item Quelle couleur prend chacune de ces deux solutions avec chacun des trois indicateurs ?
\item Quel indicateur permettrait de distinguer le plus nettement ces deux solutions ? Justifier.
\item Le papier pH donne pour le citron « pH $\approx 2$ ou 3 ». Comment obtenir la valeur 2,5 avec certitude ?
\end{enumerate}
\end{exercice}

\begin{corrige}[Exercice 1]
\textbf{1.} Jus de citron (pH $= 2{,}5$) : hélianthine \textbf{rouge} (pH $< 3{,}1$), BBT \textbf{jaune} (pH $< 6{,}0$), phénolphtaléine \textbf{incolore} (pH $< 8{,}2$). Lessive de soude (pH $= 11$) : hélianthine \textbf{jaune} (pH $> 4{,}4$), BBT \textbf{bleu} (pH $> 7{,}6$), phénolphtaléine \textbf{rose} (pH $> 10$).
\textbf{2.} La \textbf{phénolphtaléine} distingue le plus nettement les deux solutions : incolore pour le citron, rose vif pour la soude — un contraste franc, sans teinte intermédiaire possible. Le BBT convient aussi (jaune/bleu), tandis que l'hélianthine (rouge/jaune) donne un contraste moins marqué.
\textbf{3.} Il faut un \textbf{pH-mètre étalonné} : le papier pH (précision 0,5 à 1 unité) ne peut pas trancher entre 2 et 3 ; seul le pH-mètre affiche une valeur à $0,01$--$0,1$ unité près.
\end{corrige}

\begin{aretenir}[L'essentiel de la section]
\begin{itemize}
\item Trois méthodes de mesure du pH : le \textbf{papier pH} (rapide, $\pm 0{,}5$ à 1 unité), les \textbf{indicateurs colorés} (encadrement du pH par la couleur), le \textbf{pH-mètre} (précis, $\pm 0{,}01$ à 0,1 unité, après \textbf{étalonnage} avec des solutions tampons).
\item Un indicateur coloré est un couple acide/base $\ce{HInd}/\ce{Ind-}$ dont les deux formes ont des couleurs différentes ; il est caractérisé par sa \textbf{zone de virage} ($\mathrm{p}K_a \pm 1$).
\item À connaître par cœur : hélianthine (rouge 3,1 -- 4,4 jaune), BBT (jaune 6,0 -- 7,6 bleu), phénolphtaléine (incolore 8,2 -- 10 rose).
\item Un indicateur ne donne jamais une valeur précise : il \emph{encadre} le pH. Il faut le choisir adapté au milieu étudié.
\item Les jus de chou rouge, de bissap ou le curcuma sont des indicateurs naturels, préparés par macération, filtration et conservation du filtrat.
\end{itemize}
\end{aretenir}

\coursec{Préparation de solutions par dilution}

Dans la section précédente, tu as appris à \emph{mesurer} le pH d'une solution : papier pH, indicateurs colorés, pH-mètre. Mais une question pratique demeure : comment \emph{fabriquer} une solution de concentration connue, par exemple la solution d'acide chlorhydrique à \SI{0,010}{mol.L^{-1}} que tu doseras au Baccalauréat ? Au laboratoire, on ne pèse presque jamais l'acide ou la base directement : on part d'une solution concentrée (la \cle{solution mère}) et on lui ajoute de l'eau distillée pour obtenir une solution moins concentrée (la \cle{solution fille}). Cette opération s'appelle la \cle{dilution}. C'est l'un des gestes les plus fréquents du chimiste, du biologiste et du technicien de laboratoire d'analyses médicales : maîtrisons-la rigoureusement.

\courssub{Le principe de la dilution}

\textbf{Une situation concrète.} Le sirop de bissap concentré que prépare ta maman est beaucoup trop sucré et trop coloré pour être bu tel quel. Que fait-elle ? Elle verse un peu de concentré dans une calebasse et \emph{complète avec de l'eau}. Le goût s'atténue, la couleur s'éclaircit, mais la quantité de « matière de bissap » (sucre, pigments, acides) n'a pas changé : elle est simplement \emph{répartie dans un volume plus grand}. Tu viens de comprendre l'essentiel de la dilution.

\begin{definition}[Dilution]
La \cle{dilution} est l'opération qui consiste à ajouter du \emph{solvant} (en général de l'eau distillée) à une solution appelée \cle{solution mère}, de concentration $C_0$, afin d'obtenir une solution moins concentrée, appelée \cle{solution fille}, de concentration $C_1 < C_0$.
\end{definition}

Le point fondamental est le suivant : lorsqu'on ajoute de l'eau, on n'ajoute \textbf{aucune} quantité de soluté et on n'en retire aucune. Les ions ou molécules dissous sont simplement « dispersés » dans un volume plus grand. Il y a donc \emph{conservation de la quantité de matière} de soluté :

\[ n_{\text{prélevée dans la mère}} = n_{\text{présente dans la fille}} \]

Or $n = C \times V$. On en déduit la relation fondamentale.

\begin{aretenir}[Relation de dilution]
Lors d'une dilution, la quantité de matière de soluté se conserve :
\[ C_0 \times V_0 = C_1 \times V_1 \]
où $C_0$ est la concentration de la solution mère, $V_0$ le volume de solution mère prélevé, $C_1$ la concentration de la solution fille et $V_1$ le volume \emph{final} de la solution fille. Les concentrations s'expriment dans la même unité (par exemple \si{mol.L^{-1}}) et les volumes dans la même unité (par exemple \si{mL}).
\end{aretenir}

\courssub{Le facteur de dilution}

Pour caractériser « combien de fois » on a dilué, on introduit un nombre sans unité, très pratique dans les énoncés.

\begin{definition}[Facteur de dilution]
Le \cle{facteur de dilution} $F$ est le rapport :
\[ F = \frac{C_0}{C_1} = \frac{V_1}{V_0} \]
C'est un nombre \emph{sans unité}, toujours \textbf{supérieur à 1}. On dit que la solution a été diluée « $F$ fois ». Une \emph{dilution au dixième} correspond à $F = 10$ : la concentration est divisée par 10. Une dilution au centième correspond à $F = 100$.
\end{definition}

\begin{attention}[Ne pas inverser le rapport]
Le facteur de dilution est $F = \dfrac{C_0}{C_1}$ (mère sur fille), jamais $\dfrac{C_1}{C_0}$. Vérifie toujours que ton résultat est supérieur à 1 : si tu trouves $F = 0{,}1$, tu as inversé le rapport !
\end{attention}

\courssub{La verrerie de précision et le protocole de dilution}

\textbf{Pourquoi une verrerie spéciale ?} Une dilution n'est utile que si les volumes $V_0$ et $V_1$ sont connus avec précision. Un bécher ou un erlenmeyer, dont les graduations sont approximatives (à \SI{10}{\%} près !), est \emph{interdit} pour cette opération. On utilise deux instruments de \emph{verrerie jaugée} :

\begin{itemize}
\item la \cle{pipette jaugée} (par exemple \SI{10,0}{mL}), qui permet de prélever \emph{un volume fixe} $V_0$ avec une précision de l'ordre de \SI{0,05}{mL} ; on aspire le liquide à l'aide d'une propipette (poire d'aspiration) ;
\item la \cle{fiole jaugée} (par exemple \SI{100,0}{mL}), qui possède un \emph{trait de jauge} gravé sur son col étroit : remplie exactement jusqu'à ce trait, elle contient le volume $V_1$ avec une précision de l'ordre de \SI{0,1}{mL}.
\end{itemize}

La \cle{pissette} d'eau distillée sert à ajuster goutte à goutte le niveau final jusqu'au trait de jauge, sans risque de dépassement. Pour prélever un volume $V_0$ qui ne correspond pas à une pipette jaugée disponible (ex. \SI{12,5}{mL} ou \SI{8,0}{mL}), on utilise une \cle{burette} (ou, à défaut, une pipette graduée), qui permet de mesurer un volume quelconque avec une bonne précision.

\begin{methode}[Protocole d'une dilution]
Pour préparer un volume $V_1$ de solution fille de concentration $C_1$ à partir d'une solution mère de concentration $C_0$ :
\begin{enumerate}
\item \textbf{Calculer} le volume à prélever : $V_0 = \dfrac{C_1 \times V_1}{C_0}$.
\item \textbf{Prélever} $V_0$ de solution mère à l'aide d'une pipette jaugée (si $V_0$ correspond à son volume nominal) munie d'une propipette, \textbf{ou à la burette} (pour tout autre volume) ; \emph{jamais à la bouche !}
\item \textbf{Introduire} ce prélèvement dans une fiole jaugée de volume $V_1$ contenant déjà un peu d'eau distillée.
\item \textbf{Compléter} avec de l'eau distillée jusqu'aux deux tiers, boucher et agiter ; puis ajuster au trait de jauge à la pissette, le bas du ménisque étant tangent au trait de jauge, à hauteur des yeux.
\item \textbf{Homogénéiser} : boucher la fiole et la retourner une dizaine de fois.
\end{enumerate}
\end{methode}

\begin{popfigure}
\begin{tikzpicture}[scale=0.95, every node/.style={font=\small}]
% Étape 1 : pipette dans bécher mère
\draw[thick, popblue] (-0.6,0) -- (-0.6,2.2) -- (0.6,2.2) -- (0.6,0) -- cycle;
\fill[popblue!25] (-0.55,0.05) rectangle (0.55,1.1);
\draw[thick, popdark] (-0.07,0.4) -- (-0.07,3.6);
\draw[thick, popdark] (0.07,0.4) -- (0.07,3.6);
\draw[thick, popdark] (0,3.6) ellipse (0.07 and 0.03);
\fill[popblue!45] (-0.06,0.45) rectangle (0.06,2.3);
\draw[popdark] (-0.12,2.3) -- (0.12,2.3);
\node[popdark, right] at (0.7,3.3) {pipette jaugée};
\node[popdark, below] at (0,-0.15) {solution mère $C_0$};
\node[popdark, right] at (0.15,1.7) {$V_0$};
\node[font=\bfseries\small, poppurple] at (0,4.0) {1. Prélèvement};
% Flèche
\draw[-{Stealth[length=3mm]}, very thick, poporange] (1.6,1.8) -- (2.8,1.8);
% Étape 2 : fiole partiellement remplie
\begin{scope}[xshift=4.6cm]
\draw[thick, popblue] (-0.9,0) -- (-0.9,1.0) -- (-0.18,1.9) -- (-0.18,3.1);
\draw[thick, popblue] (0.9,0) -- (0.9,1.0) -- (0.18,1.9) -- (0.18,3.1);
\draw[thick, popblue] (0,0) ellipse (0.9 and 0.12);
\fill[popblue!25] (-0.88,0.06) -- (-0.88,0.7) -- (0.88,0.7) -- (0.88,0.06) -- cycle;
\draw[popdark] (-0.18,2.6) -- (0.18,2.6);
\node[popdark, right, font=\footnotesize] at (0.2,2.6) {trait de jauge};
\draw[-{Stealth}, thick, poporange] (0,3.9) -- (0,3.2);
\node[poporange, above, font=\footnotesize] at (0,3.9) {verser $V_0$};
\node[popdark, below] at (0,-0.3) {fiole jaugée $V_1$};
\node[font=\bfseries\small, poppurple] at (0,4.3) {2. Transfert};
\end{scope}
% Flèche
\draw[-{Stealth[length=3mm]}, very thick, poporange] (6.2,1.8) -- (7.4,1.8);
% Étape 3 : fiole pleine au trait de jauge
\begin{scope}[xshift=9.4cm]
\draw[thick, popblue] (-0.9,0) -- (-0.9,1.0) -- (-0.18,1.9) -- (-0.18,3.1);
\draw[thick, popblue] (0.9,0) -- (0.9,1.0) -- (0.18,1.9) -- (0.18,3.1);
\draw[thick, popblue] (0,0) ellipse (0.9 and 0.12);
\fill[popblue!12] (-0.88,0.06) -- (-0.88,1.0) -- (-0.17,1.9) -- (-0.17,2.6) -- (0.17,2.6) -- (0.17,1.9) -- (0.88,1.0) -- (0.88,0.06) -- cycle;
\draw[popdark] (-0.18,2.6) -- (0.18,2.6);
\draw[-{Stealth}, thick, popgreen] (0.9,3.6) -- (0.15,2.75);
\node[popgreen, right, font=\footnotesize, align=left] at (0.95,3.6) {pissette : ajuster\\au trait de jauge};
\node[popdark, below] at (0,-0.3) {solution fille $C_1$};
\node[font=\bfseries\small, poppurple] at (0,4.3) {3. Ajustement};
\end{scope}
\end{tikzpicture}
\legende{Les trois étapes du protocole de dilution : prélèvement de $V_0$ à la pipette jaugée (ou burette), transfert dans la fiole jaugée de volume $V_1$, puis ajustement au trait de jauge à la pissette d'eau distillée, suivi de l'homogénéisation.}
\end{popfigure}

\begin{attention}[Compléter au trait de jauge, et non « ajouter $V_1$ d'eau »]
L'erreur la plus fréquente consiste à mesurer $V_1 = \SI{100,0}{mL}$ d'eau et à les ajouter aux $V_0 = \SI{10,0}{mL}$ de solution mère. Le volume final serait alors d'environ \SI{110}{mL} (les volumes ne sont d'ailleurs pas rigoureusement additifs) : la concentration obtenue serait \emph{fausse}. La bonne formulation est : « on complète avec de l'eau distillée \textbf{jusqu'au trait de jauge} », de sorte que le volume \emph{total} soit exactement $V_1$.
\end{attention}

\begin{attention}[Sécurité : toujours verser l'acide dans l'eau]
Pour diluer un acide très concentré (acide sulfurique commercial, par exemple), il ne faut \textbf{jamais} verser l'eau sur l'acide : la dissolution est fortement exothermique, les premières gouttes d'eau bouillent instantanément et projettent de l'acide brûlant. On verse \emph{lentement} l'acide dans l'eau, en agitant, avec blouse, gants et lunettes. Retiens l'ordre : \emph{« l'acide rejoint l'eau, jamais l'inverse »}.
\end{attention}

\courssub{Conservation de la quantité de matière : regard microscopique}

Que se passe-t-il à l'échelle des ions ? Prenons une solution d'acide chlorhydrique : le chlorure d'hydrogène s'y est ionisé totalement selon :

\[ \ce{HCl(g) + H2O(l) -> H3O+(aq) + Cl-(aq)} \]

(réaction totale avec l'eau, qui joue le rôle de base). Lors de la dilution, le nombre d'ions \ce{H3O+} et \ce{Cl-} prélevés ne change pas ; seul le volume dans lequel ils « nagent » augmente. La concentration, qui est un \emph{nombre d'ions par unité de volume}, diminue donc.

\begin{popfigure}
\begin{tikzpicture}[scale=1.0]
% Avant
\draw[thick, popdark, fill=popblue!15] (0,0) rectangle (3,3);
\node[popdark, below, font=\small] at (1.5,0) {avant : volume $V_0$};
\foreach \x/\y in {0.5/0.5, 1.5/0.7, 2.5/0.5, 0.7/1.6, 1.6/1.8, 2.4/1.5, 0.5/2.5, 1.5/2.6, 2.5/2.5}{
  \fill[popink] (\x,\y) circle (0.09);
  \fill[popgreen] (\x+0.22,\y+0.15) circle (0.09);
}
\node[popdark, font=\small] at (1.5,3.25) {9 paires d'ions};
% Flèche
\draw[-{Stealth[length=3mm]}, very thick, poporange] (3.4,1.5) -- (4.6,1.5);
\node[poporange, above, font=\footnotesize] at (4,1.5) {$+\,\ce{H2O}$};
% Après
\draw[thick, popdark, fill=popblue!8] (5,0) rectangle (10,3);
\node[popdark, below, font=\small] at (7.5,0) {après : volume $V_1 = 10\,V_0$};
\foreach \x/\y in {5.4/0.4, 7.2/0.6, 9.3/0.4, 5.8/1.5, 7.6/1.7, 9.4/1.4, 5.5/2.5, 7.4/2.6, 9.2/2.5}{
  \fill[popink] (\x,\y) circle (0.09);
  \fill[popgreen] (\x+0.22,\y+0.15) circle (0.09);
}
\node[popdark, font=\small] at (7.5,3.25) {toujours 9 paires d'ions};
\fill[popink] (5.2,-0.9) circle (0.09); \node[right, font=\footnotesize] at (5.4,-0.9) {\ce{H3O+}};
\fill[popgreen] (7.2,-0.9) circle (0.09); \node[right, font=\footnotesize] at (7.4,-0.9) {\ce{Cl-}};
\end{tikzpicture}
\legende{Interprétation microscopique d'une dilution au dixième ($F = 10$) : le nombre d'ions est conservé, mais ils occupent un volume dix fois plus grand ; la concentration est donc divisée par dix.}
\end{popfigure}

\courssub{Exemple chiffré et conséquence sur le pH}

\begin{exemplebox}[Préparation d'une solution au dixième]
On veut préparer $V_1 = \SI{100,0}{mL}$ d'une solution d'acide chlorhydrique de concentration $C_1 = \SI{0,010}{mol.L^{-1}}$ à partir d'une solution mère de concentration $C_0 = \SI{0,100}{mol.L^{-1}}$.

\textbf{Calcul du volume à prélever :}
\[ V_0 = \frac{C_1 \times V_1}{C_0} = \frac{0{,}010 \times 100{,}0}{0{,}100} = \SI{10,0}{mL} \]

\textbf{Facteur de dilution :} $F = \dfrac{C_0}{C_1} = \dfrac{0{,}100}{0{,}010} = 10$ (dilution au dixième).

\textbf{Protocole :} prélever \SI{10,0}{mL} de solution mère à la pipette jaugée de \SI{10,0}{mL}, les introduire dans une fiole jaugée de \SI{100,0}{mL} contenant un peu d'eau distillée, compléter au trait de jauge, homogénéiser.
\end{exemplebox}

\textbf{Et le pH dans tout cela ?} L'acide chlorhydrique est un \emph{acide fort} : il est totalement ionisé, donc $[\ce{H3O+}] = C$ et $\mathrm{pH} = -\log C$. La solution mère a pour pH :

\[ \mathrm{pH}_0 = -\log(0{,}100) = 1{,}0 \]

Après dilution au dixième :

\[ \mathrm{pH}_1 = -\log(0{,}010) = 2{,}0 \]

Le pH a \emph{augmenté de 1 unité} : la solution est devenue moins acide. Symétriquement, pour une base forte comme l'hydroxyde de sodium \ce{NaOH}, totalement dissociée selon $\ce{NaOH(s) ->[eau] Na+(aq) + HO-(aq)}$, on a $[\ce{HO-}] = C$, donc $[\ce{H3O+}] = \dfrac{K_e}{C}$ : diviser $C$ par 10 \emph{diminue} le pH de 1 unité.

\begin{propriete}[Dilution et pH, à \SI{25}{\celsius}]
\begin{itemize}
\item Diluer $10$ fois une solution d'\textbf{acide fort} : son pH \emph{augmente} de $1$ unité.
\item Diluer $10$ fois une solution de \textbf{base forte} : son pH \emph{diminue} de $1$ unité.
\end{itemize}
En revanche, on ne peut jamais, par dilution, faire passer un acide au-dessus de 7 ni une base en dessous de 7 : à très grande dilution, le pH tend vers 7, celui de l'eau pure, car les ions \ce{H3O+} et \ce{HO-} de l'autoprotolyse de l'eau ($\ce{2 H2O <=> H3O+ + HO-}$) deviennent non négligeables.
\end{propriete}

\begin{attention}[Approximation des espèces ultra-minoritaires]
Dans une solution acide, on néglige en général les ions \ce{HO-} devant les ions \ce{H3O+}. Mais si la solution acide est diluée jusqu'à $C \approx \SI{e-7}{mol.L^{-1}}$, cette approximation devient fausse : les \ce{H3O+} issus de l'autoprotolyse de l'eau ne sont plus négligeables. Le pH d'un acide fort dilué ne pourra jamais dépasser 7, quelle que soit la dilution.
\end{attention}

\begin{exoresolu}[Préparer une solution pour un dosage]
Au laboratoire du lycée, on dispose d'une solution mère d'hydroxyde de sodium de concentration $C_0 = \SI{1,0}{mol.L^{-1}}$. On souhaite préparer $V_1 = \SI{250,0}{mL}$ d'une solution de concentration $C_1 = \SI{0,050}{mol.L^{-1}}$ pour un dosage acido-basique.
\begin{enumerate}
\item Calculer le volume $V_0$ de solution mère à prélever.
\item Calculer le facteur de dilution.
\item Décrire le protocole et préciser la verrerie.
\item Quel est le pH de la solution fille à \SI{25}{\celsius} ? On donne $\mathrm{p}K_e = 14{,}0$.
\end{enumerate}
\tcblower
\textbf{1. Volume à prélever.} La conservation de la quantité de matière s'écrit $C_0 V_0 = C_1 V_1$, d'où :
\[ V_0 = \frac{C_1 \times V_1}{C_0} = \frac{0{,}050 \times 250{,}0}{1{,}0} = \SI{12,5}{mL} \]

\textbf{2. Facteur de dilution.}
\[ F = \frac{C_0}{C_1} = \frac{1{,}0}{0{,}050} = 20 \]
La solution est diluée 20 fois (on vérifie : $V_1/V_0 = 250{,}0/12{,}5 = 20$).

\textbf{3. Protocole.} À l'aide d'une \textbf{burette} (ou, à défaut, d'une pipette graduée de \SI{25}{mL}), prélever \SI{12,5}{mL} de solution mère ; les verser dans une fiole jaugée de \SI{250,0}{mL} contenant déjà de l'eau distillée ; compléter au trait de jauge à la pissette ; boucher et homogénéiser en retournant la fiole une dizaine de fois.

\textbf{4. pH de la solution fille.} \ce{NaOH} est une base forte, totalement dissociée : $[\ce{HO-}] = C_1 = \SI{0,050}{mol.L^{-1}}$. D'après le produit ionique de l'eau :
\[ [\ce{H3O+}] = \frac{K_e}{[\ce{HO-}]} = \frac{10^{-14{,}0}}{5{,}0\times 10^{-2}} = 2{,}0\times 10^{-13}\ \si{mol.L^{-1}} \]
\[ \mathrm{pH} = -\log(2{,}0\times 10^{-13}) = 12{,}7 \]
La solution est bien basique ($\mathrm{pH} > 7$), ce qui est cohérent.
\end{exoresolu}

\courssub{Pourquoi diluer ? Applications au quotidien et au Baccalauréat}

\begin{savaistu}[La dilution, un geste omniprésent]
\begin{itemize}
\item \textbf{Analyses médicales :} dans les laboratoires de Yaoundé ou de Douala, le sérum d'un patient est souvent trop concentré pour être mesuré directement ; le technicien le dilue 10 ou 100 fois avant le dosage, puis multiplie le résultat par $F$.
\item \textbf{Produits ménagers :} l'eau de Javel concentrée et les détergents vendus au marché portent la mention « à diluer » : l'utilisateur réalise, sans le savoir, une dilution de facteur $F$ donné par le fabricant.
\item \textbf{Agriculture :} les engrais liquides et les pesticides sont livrés concentrés ; le planteur prépare la bouillie de pulvérisation en respectant un facteur de dilution précis, sous peine de brûler les plants.
\item \textbf{Au Baccalauréat :} l'épreuve de TP de chimie exige presque toujours la préparation d'une solution fille avant un dosage : relation $C_0V_0 = C_1V_1$, choix de la verrerie jaugée et rédaction du protocole sont des points de notation classiques.
\end{itemize}
\end{savaistu}

\begin{exercice}[À toi de jouer]
On dispose d'une solution mère d'acide chlorhydrique de concentration $C_0 = \SI{0,50}{mol.L^{-1}}$ et de pH $= 0{,}3$.
\begin{enumerate}
\item Quel volume $V_0$ faut-il prélever pour préparer \SI{200,0}{mL} de solution de concentration $C_1 = \SI{0,020}{mol.L^{-1}}$ ?
\item Quel est le facteur de dilution ?
\item Quel est le pH de la solution fille ? Vérifier la règle « diluer 10 fois un acide fort augmente le pH de 1 ».
\end{enumerate}
\end{exercice}

\begin{corrige}[Exercice : dilution de l'acide chlorhydrique]
\textbf{1.} $V_0 = \dfrac{C_1 V_1}{C_0} = \dfrac{0{,}020 \times 200{,}0}{0{,}50} = \SI{8,0}{mL}$, prélevés à la \textbf{pipette graduée de 10 mL} (ou à la \textbf{burette}), versés dans une fiole jaugée de \SI{200,0}{mL}, complétée au trait de jauge.

\textbf{2.} $F = \dfrac{C_0}{C_1} = \dfrac{0{,}50}{0{,}020} = 25$. La solution est diluée 25 fois.

\textbf{3.} Acide fort : $\mathrm{pH}_1 = -\log C_1 = -\log(0{,}020) = 1{,}7$. La dilution est ici de facteur 25, soit $10 \times 2{,}5$ : le pH a augmenté de $\log 25 \approx 1{,}4$ (de $0{,}3$ à $1{,}7$). La règle « $+1$ par dilution au dixième » s'applique bien : après une première dilution au dixième ($C = \SI{0,050}{mol.L^{-1}}$), le pH vaudrait $1{,}3$, soit exactement $0{,}3 + 1$.
\end{corrige}

\begin{aretenir}[L'essentiel de la section]
\begin{itemize}
\item Diluer, c'est ajouter du solvant : la quantité de matière de soluté se conserve, d'où $C_0 V_0 = C_1 V_1$.
\item Le facteur de dilution $F = \dfrac{C_0}{C_1} = \dfrac{V_1}{V_0} > 1$ ; « au dixième » signifie $F = 10$.
\item Verrerie de précision : pipette jaugée (volume fixe) ou burette (volume quelconque) pour prélever $V_0$, fiole jaugée pour le volume final $V_1$, pissette pour ajuster au trait de jauge ; on \emph{complète} jusqu'au trait, on n'ajoute pas $V_1$ d'eau.
\item Sécurité : toujours verser l'acide concentré dans l'eau, jamais l'inverse.
\item Diluer 10 fois un acide fort : $\mathrm{pH} + 1$ ; diluer 10 fois une base forte : $\mathrm{pH} - 1$ ; le pH tend vers 7 à très grande dilution, sans jamais le franchir.
\end{itemize}
\end{aretenir}

\newpage

\coursec{Travaux pratiques}

\courssub{Objectifs du TP}

Ce travail pratique te place dans la peau d'un \cle{technicien de laboratoire} : comme l'infirmier qui lit un résultat d'analyse sanguine ou le contrôleur qualité d'une savonnerie artisanale, tu vas mesurer et interpréter le pH de produits du quotidien. À l'issue de la séance, tu seras capable de :

\begin{itemize}
\item mesurer le pH d'une solution par \cle{trois méthodes} différentes : papier pH, pH-mètre et indicateur coloré naturel ;
\item fabriquer un \cle{indicateur coloré naturel} à partir de fleurs de bissap (hibiscus), disponibles sur tous les marchés du Cameroun ;
\item préparer une solution diluée de facteur de dilution $F = 10$ et vérifier expérimentalement l'effet de la dilution sur le pH ;
\item calculer les concentrations $[\ce{H3O+}]$ et $[\ce{OH-}]$ à partir du pH mesuré ;
\item comparer la \cle{précision} des trois méthodes de mesure et classer des produits usuels du plus acide au plus basique.
\end{itemize}

\courssub{Matériel et produits}

\begin{center}
\begin{tabularx}{\linewidth}{|l|X|}
\hline
\rowcolor{popblue!40} \textbf{Matériel} & \textbf{Produits} \\
\hline
6 béchers de \SI{100}{mL} & Fleurs de bissap séchées (marché local) \\
Pipette jaugée de \SI{10}{mL} + poire & Jus de citron frais \\
Fiole jaugée de \SI{100}{mL} & Vinaigre du commerce \\
Papier pH (bandelettes) + baguettes de verre & Eau du robinet (ou eau de forage) \\
pH-mètre étalonné (si disponible) & Solution de savon (lessive de savon artisanal) \\
Entonnoir + papier filtre (ou tissu propre) & Eau de Javel diluée au 1/10 \\
Casserole ou bécher + source de chaleur & Solution d'acide chlorhydrique à \SI{0,1}{mol.L^{-1}} \\
Eprouvette graduée de \SI{50}{mL} & Eau distillée (ou eau bouillie refroidie) \\
\hline
\end{tabularx}
\end{center}

\begin{savaistu}[Pas de pH-mètre au lycée ?]
De nombreux lycées camerounais ne disposent pas de pH-mètre. Ce n'est pas grave : le papier pH et l'indicateur de bissap suffisent pour toutes les mesures de ce TP. À défaut de papier pH commercial, on peut utiliser du papier de curcuma (trempé dans une décoction de curcuma puis séché) qui vire du jaune au rouge-brun en milieu basique. L'eau distillée peut être remplacée par de l'eau de pluie bouillie puis refroidie, ou par l'eau du condensat d'un fer à repasser à vapeur.
\end{savaistu}

\courssub{Sécurité}

\begin{attention}[Consignes de sécurité obligatoires]
\begin{itemize}
\item \textbf{Port de la blouse et des lunettes} pendant toute la séance ; cheveux longs attachés.
\item \textbf{Acide chlorhydrique \SI{0,1}{mol.L^{-1}}} : solution \emph{irritante} (pictogramme : point d'exclamation sur fond blanc, losange rouge). Ne jamais pipeter à la bouche : utiliser systématiquement la \cle{poire à pipeter}. En cas de contact avec la peau, rincer abondamment à l'eau froide pendant plusieurs minutes.
\item \textbf{Eau de Javel} : solution \emph{corrosive} (pictogramme : liquide attaquant une surface et une main). Ne jamais la mélanger avec un acide : il se dégagerait du dichlore \ce{Cl2}, gaz toxique. Manipuler près d'une fenêtre ouverte.
\item \textbf{Potasse et savon} : solutions basiques irritantes pour les yeux ; se laver les mains en fin de séance.
\item \textbf{Eau chaude} pour la macération : manipuler le bécher avec des pincettes en bois pour éviter les brûlures.
\item Ne jamais goûter les produits, même le jus de citron utilisé au laboratoire.
\end{itemize}
\end{attention}

\courssub{Schéma du dispositif}

\begin{popfigure}
\begin{center}
\begin{tikzpicture}[scale=0.9, every node/.style={font=\small}]
% Bécher de macération
\draw[thick, popdark] (0,0) -- (0,2.2) (1.6,0) -- (1.6,2.2);
\draw[thick, popdark] (0,0) -- (1.6,0);
\draw[thick, popdark] (1.6,2.2) -- (1.9,2.5);
\fill[poppinkL] (0.06,0.06) rectangle (1.54,1.3);
\draw[popink] (0.06,1.3) -- (1.54,1.3);
\node[popdark] at (0.8,-0.4) {macération bissap};
\node[popdark, font=\footnotesize] at (0.8,0.65) {infusion rouge};
% Flèche filtration
\draw[-{Stealth}, thick, poporange] (2.1,1.4) -- (3.1,1.4) node[midway, above, font=\footnotesize]{filtrer};
% Entonnoir + bécher
\draw[thick, popdark] (3.4,2.2) -- (4.3,1.5) -- (4.3,1.1);
\draw[thick, popdark] (5.2,2.2) -- (4.3,1.5);
\draw[thick, popdark] (3.4,2.2) -- (5.2,2.2);
\draw[thick, popdark] (3.6,0) -- (3.6,1.0) (5.0,0) -- (5.0,1.0);
\draw[thick, popdark] (3.6,0) -- (5.0,0);
\fill[poppinkL] (3.66,0.06) rectangle (4.94,0.6);
\node[popdark] at (4.3,-0.4) {indicateur filtré};
% Mesure pH-mètre
\draw[thick, popdark] (6.4,0) -- (6.4,1.6) (8.0,0) -- (8.0,1.6);
\draw[thick, popdark] (6.4,0) -- (8.0,0);
\fill[popgoldL] (6.46,0.06) rectangle (7.94,0.9);
\draw[popink] (6.46,0.9) -- (7.94,0.9);
\draw[thick, popblue] (7.2,0.3) -- (7.2,2.4) -- (8.6,2.4);
\draw[thick, popblue, fill=popblueL] (8.3,2.4) rectangle (9.6,3.1);
\node[popdark, font=\footnotesize] at (8.95,2.75) {pH-mètre};
\node[popdark] at (7.2,-0.4) {mesure du pH};
% Bandelette papier pH
\draw[thick, popgreen, fill=popgreenL] (10.3,0.3) rectangle (11.6,0.7);
\node[popdark, font=\footnotesize] at (10.95,0.5) {papier pH};
\node[popdark] at (10.95,-0.4) {trempage + lecture};
\end{tikzpicture}
\end{center}
\legende{Chaîne des opérations du TP : préparation de l'indicateur de bissap par macération et filtration, puis mesure du pH des échantillons au pH-mètre ou au papier pH.}
\end{popfigure}

\courssub{Protocole}

\textbf{Partie 1 : fabrication de l'indicateur coloré naturel}\par
\begin{enumerate}
\item Verser environ \SI{100}{mL} d'eau distillée (ou bouillie refroidie) dans un bécher et porter à environ \SI{60}{\celsius}.
\item Introduire une petite poignée de fleurs de bissap séchées et laisser \cle{macérer 10 minutes} en agitant de temps en temps : l'eau prend une teinte rouge foncé.
\item Filtrer (papier filtre ou tissu propre) et recueillir le \cle{filtrat rouge} dans un bécher propre : c'est ton indicateur coloré. Laisser refroidir.
\end{enumerate}

\textbf{Partie 2 : mesure du pH des produits usuels}\par
\begin{enumerate}
\setcounter{enumi}{3}
\item Répartir environ \SI{20}{mL} de chaque produit (jus de citron, vinaigre, eau du robinet, solution de savon, eau de Javel diluée) dans des béchers étiquetés.
\item \textbf{Au papier pH} : déposer une goutte de chaque solution sur une bandelette à l'aide d'une baguette de verre propre ; comparer la couleur à l'échelle étalon ; noter le pH.
\item \textbf{Au pH-mètre} (si disponible) :
\begin{enumerate}
\item Vérifier l'étalonnage du pH-mètre avec les solutions tampons pH 4, 7 et 10.
\item Rincer l'électrode à l'eau distillée, l'essuyer délicatement avec du papier absorbant (sans frotter le bulbe).
\item Plonger l'électrode dans l'échantillon, agiter légèrement, attendre la stabilisation et noter le pH à $0,1$ unité près.
\item Rincer l'électrode à l'eau distillée entre chaque mesure.
\end{enumerate}
\item \textbf{À l'indicateur de bissap} : ajouter 3 gouttes d'indicateur dans chaque échantillon, agiter et noter la couleur obtenue.
\end{enumerate}

\begin{methode}[Mesure du pH au pH-mètre]
\begin{enumerate}
\item \textbf{Étalonnage} : immerger l'électrode successivement dans deux solutions tampons (ex. pH 4,01 et pH 7,00 ou pH 7,00 et pH 10,01) et valider les valeurs.
\item \textbf{Rinçage} : rincer abondamment l'électrode à l'eau distillée, sécher le pourtour par tamponnage (papier essuie-tout), \emph{ne pas essuyer le bulbe sensible}.
\item \textbf{Mesure} : plonger l'électrode dans la solution testée, homogénéiser, attendre la stabilisation de l'affichage (critère de dérive $< 0,01$ pH/min).
\item \textbf{Re-ralisation} : rincer à nouveau avant la mesure suivante.
\end{enumerate}
\end{methode}

\textbf{Partie 3 : dilution d'une solution d'acide chlorhydrique}\par
\begin{enumerate}
\setcounter{enumi}{7}
\item Mesurer le pH de la solution mère \ce{HCl} à $C_0 = \SI{0,1}{mol.L^{-1}}$.
\item Rincer la pipette jaugée de \SI{10}{mL} avec un peu de solution mère (opération de \cle{rinçage à l'essence}), puis prélever $V_0 = \SI{10,0}{mL}$ de solution mère à l'aide de la poire.
\item Verser dans une fiole jaugée de \SI{100}{mL} (préalablement rincée à l'eau distillée), compléter à l'eau distillée jusqu'au trait de jauge, boucher et homogénéiser par inversions successives : le facteur de dilution est $F = \dfrac{100}{10} = 10$, d'où $C_1 = \SI{0,01}{mol.L^{-1}}$.
\item Mesurer le pH de la solution diluée et comparer au pH de la solution mère.
\end{enumerate}

\begin{methode}[Préparation d'une solution par dilution (facteur $F$)]
\begin{enumerate}
\item Calculer le volume $V_0$ à prélever : $V_0 = \dfrac{V_1}{F}$ ($V_1$ volume de la fiole jaugée).
\item Rincer la pipette jaugée $V_0$ avec la solution mère (2 ou 3 rinçages).
\item Prélever $V_0$ de solution mère, la transférer dans la fiole jaugée $V_1$ (rincée à l'eau distillée).
\item Compléter à l'eau distillée jusqu'au trait de jauge (bas du ménisque au niveau du trait).
\item Boucher, mélanger par inversions répétées (au moins 10) pour assurer l'homogénéité.
\item La concentration finale est $C_1 = \dfrac{C_0}{F}$.
\end{enumerate}
\end{methode}

\courssub{Observations et résultats attendus}

\begin{center}
\begin{tabularx}{\linewidth}{|X|c|c|X|c|}
\hline
\rowcolor{popgreen!50} \textbf{Produit testé} & \textbf{pH (papier)} & \textbf{pH (pH-mètre)} & \textbf{Couleur au bissap} & \textbf{Caractère} \\
\hline
\rowcolor{popgreenL} Jus de citron & 2 & 2,3 & rouge vif & Acide \\
\rowcolor{white} Vinaigre & 3 & 2,9 & rouge rosé & Acide \\
\rowcolor{popgreenL} Eau du robinet & 7 & 7,1 & violet-rose pâle & Neutre \\
\rowcolor{white} Solution de savon & 10 & 9,8 & vert-bleu & Basique \\
\rowcolor{popgreenL} Eau de Javel diluée & 11 & 10,9 & vert-jaune & Basique \\
\rowcolor{white} \ce{HCl} \SI{0,1}{mol.L^{-1}} & 1 & 1,1 & rouge intense & Acide \\
\rowcolor{popgreenL} \ce{HCl} \SI{0,01}{mol.L^{-1}} & 2 & 2,0 & rouge vif & Acide \\
\hline
\end{tabularx}
\end{center}

On constate que le pH de la solution d'acide chlorhydrique \cle{augmente d'une unité} (de 1 à 2) lorsqu'on dilue dix fois : la concentration en ions \ce{H3O+} a été divisée par 10, et comme $\mathrm{pH} = -\log[\ce{H3O+}]$, le pH augmente de $\log 10 = 1$.

\courssub{Exploitation}

\begin{exoresolu}[Questions d'exploitation]
\textbf{1.} Écrire l'équation-bilan de l'autoprotolyse de l'eau et donner l'expression du produit ionique $K_e$.

\textbf{2.} Pour le jus de citron ($\mathrm{pH} = 2,3$) et pour la solution de savon ($\mathrm{pH} = 9,8$), calculer $[\ce{H3O+}]$ et $[\ce{OH-}]$ à \SI{25}{\celsius} ($pK_e = 14$).

\textbf{3.} Vérifier par le calcul que la dilution de la solution d'\ce{HCl} fait passer le pH de 1 à 2.

\textbf{4.} Classer les produits testés du plus acide au plus basique.

\textbf{5.} Comparer la précision des trois méthodes de mesure du pH.
\tcblower
\textbf{1.} L'autoprotolyse de l'eau s'écrit :
\[ \ce{2 H2O <=> H3O+ + OH-} \qquad K_e = [\ce{H3O+}][\ce{OH-}] = 10^{-14} \text{ à } \SI{25}{\celsius}. \]

\textbf{2.} On utilise $[\ce{H3O+}] = 10^{-\mathrm{pH}}$ et $[\ce{OH-}] = \dfrac{K_e}{[\ce{H3O+}]} = 10^{\mathrm{pH}-14}$.
\begin{align*}
\text{Citron : } & [\ce{H3O+}] = 10^{-2,3} = \SI{5,0e-3}{mol.L^{-1}} ; \quad [\ce{OH-}] = 10^{-11,7} = \SI{2,0e-12}{mol.L^{-1}}.\\
\text{Savon : } & [\ce{H3O+}] = 10^{-9,8} = \SI{1,6e-10}{mol.L^{-1}} ; \quad [\ce{OH-}] = 10^{-4,2} = \SI{6,3e-5}{mol.L^{-1}}.
\end{align*}
Dans le citron, \ce{H3O+} est \cle{majoritaire} et \ce{OH-} \cle{ultra-minoritaire} ; c'est l'inverse dans la solution de savon.

\textbf{3.} L'acide chlorhydrique est un acide fort, totalement dissocié : $[\ce{H3O+}] = C$.
\begin{itemize}
\item Solution mère : $\mathrm{pH} = -\log(0,1) = 1$.
\item Solution diluée : $C_1 = \dfrac{C_0}{F} = \dfrac{0,1}{10} = \SI{0,01}{mol.L^{-1}}$, donc $\mathrm{pH} = -\log(0,01) = 2$.
\end{itemize}
Le pH a bien augmenté d'une unité, conformément à la mesure.

\textbf{4.} Classement du plus acide (pH le plus bas) au plus basique (pH le plus haut) :
\begin{enumerate}
\item \ce{HCl} \SI{0,1}{mol.L^{-1}} (pH $\approx$ 1)
\item \ce{HCl} \SI{0,01}{mol.L^{-1}} (pH $\approx$ 2)
\item Jus de citron (pH $\approx$ 2,3)
\item Vinaigre (pH $\approx$ 2,9)
\item Eau du robinet (pH $\approx$ 7)
\item Solution de savon (pH $\approx$ 9,8)
\item Eau de Javel diluée (pH $\approx$ 10,9)
\end{enumerate}

\textbf{5.} Le \cle{pH-mètre} est la méthode la plus précise (à $0,1$ voire $0,01$ unité près) mais exige un étalonnage rigoureux. Le \cle{papier pH} donne une valeur à une unité près (demi-unité pour les papiers de précision) : rapide et bon marché, mais peu précis. L'\cle{indicateur de bissap} ne donne qu'une estimation par couleur (à 1 ou 2 unités près), mais il est gratuit, renouvelable et disponible partout : c'est un excellent outil pédagogique et de dépannage.
\end{exoresolu}

\courssub{Conclusion}

Ce TP t'a permis de mesurer le pH de produits familiers par trois méthodes complémentaires et de vérifier expérimentalement deux lois essentielles : le lien $\mathrm{pH} = -\log[\ce{H3O+}]$ et l'effet de la dilution (diviser la concentration par 10 augmente le pH d'une unité pour un acide fort dilué). Tu as aussi fabriqué un indicateur coloré naturel au bissap, dont le virage du rouge (milieu acide) au vert (milieu basique) illustre concrètement le principe des indicateurs colorés. Ces compétences — mesurer, diluer, interpréter — sont exactement celles du technicien d'analyse, du contrôleur qualité ou de l'infirmier lisant un bilan sanguin : la chimie du pH est partout autour de toi.

\newpage

\coursec{Méthodes et exercices résolus}

\coursec{L'eau, un solvant exceptionnel}

\begin{definition}[Molécule d'eau : géométrie et polarité]
La molécule d'eau \ce{H2O} possède une géométrie \textbf{coudée} (angle \ce{H-O-H} \approx \SI{104,5}{\degree}) due aux deux doublets non liants sur l'oxygène (répulsion VSEPR). L'oxygène étant plus électronégatif que l'hydrogène (\SI{3,44} contre \SI{2,20} sur l'échelle de Pauling), chaque liaison \ce{O-H} est polarisée \ce{O^{\delta-}-H^{\delta+}}. La molécule possède un \textbf{moment dipolaire} permanent (\SI{1,85}{D}), ce qui en fait une molécule \cle{polare}.
\end{definition}

\begin{popfigure}
\begin{tikzpicture}[scale=1.2, every node/.style={font=\small}]
% Atomes
\node[draw, circle, fill=poppink, text=white, minimum size=8mm] (O) at (0,0) {\ce{O}};
\node[draw, circle, fill=popblue, text=white, minimum size=6mm] (H1) at (1.2,0.8) {\ce{H}};
\node[draw, circle, fill=popblue, text=white, minimum size=6mm] (H2) at (1.2,-0.8) {\ce{H}};
% Liaisons
\draw[thick, popdark] (O) -- (H1);
\draw[thick, popdark] (O) -- (H2);
% Doublets non liants (représentation schématique)
\draw[thick, poppurple, decorate, decoration={snake, amplitude=1pt, segment length=4pt}] (-0.8,0.5) -- (-1.5,1.0);
\draw[thick, poppurple, decorate, decoration={snake, amplitude=1pt, segment length=4pt}] (-0.8,-0.5) -- (-1.5,-1.0);
\node[poppurple, anchor=east] at (-0.9,0.5) {2 doublets non liants};
% Dipôles de liaison
\draw[->, thick, poporange] (0.2,0.2) -- (0.9,0.6) node[midway, above, sloped, font=\tiny] {$\delta^- \leftarrow \delta^+$};
\draw[->, thick, poporange] (0.2,-0.2) -- (0.9,-0.6) node[midway, below, sloped, font=\tiny] {$\delta^- \leftarrow \delta^+$};
% Dipôle moléculaire résultant
\draw[->, very thick, popgreen] (0,0) -- (-1.5,0) node[midway, above] {$\vec{\mu}$ \SI{1,85}{D}};
% Angle
\draw[popmuted, <->] (0.5,0.3) arc (33.7:-33.7:0.5) node[midway, right] {$104,5^\circ$};
\end{tikzpicture}
\legende{Structure de la molécule d'eau : géométrie coudée, polarité des liaisons et dipôle résultant.}
\end{popfigure}

\begin{propriete}[Pouvoirs de l'eau : dissolvant, dispersant, solvatant/ionisant]
\begin{itemize}
\item \textbf{Dissolvant / Dispersant :} L'eau dissout les espèces ioniques (sel, soude) et polaires (sucre, éthanol, acides) par \cle{solvatation} : les ions \ce{H3O+} (ou \ce{H+}) sont solvatés par les pôles \ce{\delta-} (O), les anions par les pôles \ce{\delta+} (H). Elle forme des dispersions colloïdales (gel d'amidon, lait).
\item \textbf{Solvant ionisant :} Elle favorise la dissociation des acides et bases (ex. \ce{HCl -> H3O+ + Cl-}) en stabilisant les ions formés.
\item \textbf{Concentration de l'eau pure :} Masse molaire \SI{18,0}{g.mol^{-1}}, masse volumique \SI{1,00}{g.mL^{-1}} (à \SI{25}{\celsius}) $\Rightarrow$ $C_{\ce{H2O}} = \frac{1000}{18,0} \approx \SI{55,5}{mol.L^{-1}}$. L'eau est en \textbf{excès immense} par rapport aux solutés usuels ($< \SI{10}{mol.L^{-1}}$).
\end{itemize}
\end{propriete}

\begin{savaistu}[L'eau au Cameroun : ressources et enjeux]
Le Cameroun, « Afrique en miniature », possède d'immenses ressources en eau (fleuves Sanaga, Wouri, Bénoué ; lacs volcaniques Nyos, Monoun ; nappes phréatiques). La \textbf{SONARA} (Société Nationale de Raffinage) utilise l'eau comme fluide caloporteur et pour le traitement des effluents acides (neutralisation à la chaux). Les \textbf{brasseries} (Isenberg, Mutzig) et la production de \textbf{bili-bili} (bière de mil) ou de \textbf{vin de palme} surveillent le pH de l'eau de brassage (idéalement 5,2--5,6) pour l'activité enzymatique. La \textbf{CIMENCAM} exploite le calcaire (\ce{CaCO3}) dont la solubilité dépend du pH et du \ce{CO2} dissous.
\end{savaistu}

\coursec{Acides et bases selon Brønsted}

\begin{definition}[Acide et base de Brønsted]
Un \cle{acide} (noté \ce{AH}) est une espèce chimique capable de \textbf{céder un proton} \ce{H+} ; une \cle{base} (notée \ce{B}) est une espèce capable de \textbf{capturer un proton}.
Lorsqu'un acide \ce{AH} cède son proton, il se transforme en sa \cle{base conjuguée} \ce{A-}. Lorsqu'une base \ce{B} capture un proton, elle devient son \cle{acide conjugué} \ce{BH+}.
On note le \textbf{couple acide/base} \ce{AH}/\ce{A-} ou \ce{BH+}/\ce{B}.
\end{definition}

\begin{propriete}[Réaction acido-basique : transfert de proton]
La réaction entre un acide \ce{AH_1} et une base \ce{A_2^-} (base du couple \ce{AH_2}/\ce{A_2^-}) s'écrit :
\[ \ce{AH_1 + A_2^- <=> A_1^- + AH_2} \]
C'est un \textbf{transfert de proton} de \ce{AH_1} vers \ce{A_2^-}. L'équilibre se déplace vers le couple le plus faible (acide le plus faible / base la plus faible).
\end{propriete}

\begin{exemplebox}[Couples acide/base usuels]
\begin{center}
\begin{tabularx}{\linewidth}{|l|X|X|}\hline
\rowcolor{popblueL}
\textbf{Couple} & \textbf{Acide} & \textbf{Base} \\\hline
\ce{H3O+}/\ce{H2O} & \ce{H3O+} (ion oxonium) & \ce{H2O} \\\hline
\ce{H2O}/\ce{OH-} & \ce{H2O} & \ce{OH-} (ion hydroxyde) \\\hline
\ce{CH3COOH}/\ce{CH3COO-} & Acide éthanoïque & Ion éthanoate \\\hline
\ce{NH4+}/\ce{NH3} & Ion ammonium & Ammoniac \\\hline
\ce{H2CO3}/\ce{HCO3-} & Acide carbonique & Ion hydrogénocarbonate \\\hline
\ce{HCO3-}/\ce{CO3^2-} & Ion hydrogénocarbonate & Ion carbonate \\\hline
\ce{HSO4-}/\ce{SO4^2-} & Ion hydrogénosulfate & Ion sulfate \\\hline
\end{tabularx}
\end{center}
\textbf{Note :} \ce{H2O} apparaît dans deux couples : elle est \cle{ampholyte} (acide dans \ce{H2O}/\ce{OH-}, base dans \ce{H3O+}/\ce{H2O}).
\end{exemplebox}

\coursec{Autoprotolyse de l'eau et produit ionique}

\courssub{Méthode 1 : Écrire l'équation-bilan de l'autoprotolyse de l'eau}

\begin{methode}[Écrire l'équation de l'autoprotolyse de l'eau]
\begin{enumerate}
\item Identifier les \textbf{deux couples acide/base} de l'eau (espèce \cle{ampholyte}) : $\ce{H2O}/\ce{OH-}$ (l'eau acide, donne \ce{H+}) et $\ce{H3O+}/\ce{H2O}$ (l'eau base, capture \ce{H+}).
\item Écrire les deux demi-équations de transfert de proton :
\[ \ce{H2O <=> H+ + OH-} \quad \text{(eau acide)} \qquad \ce{H2O + H+ <=> H3O+} \quad \text{(eau base)} \]
\item Additionner membre à membre : le proton \ce{H+} n'existe pas libre en solution aqueuse, il est instantanément solvaté par une molécule d'eau en \ce{H3O+}.
\item Conclure à l'équation-bilan :
\[ \ce{2 H2O <=> H3O+ + OH-} \]
\item Préciser le caractère \textbf{limité} de la réaction (double flèche \ce{<=>}) et l'égalité stœchiométrique imposée dans l'eau pure : $[\ce{H3O+}] = [\ce{OH-}]$.
\end{enumerate}
\textbf{Réflexe Bac :} Toute réaction acido-basique est un \cle{transfert de proton} ; le proton ne voyage jamais seul, il est toujours fixé sur une molécule d'eau (formation de \ce{H3O+}).
\end{methode}

\begin{aretenir}[Produit ionique de l'eau $K_e$ et $pK_e$]
Pour l'équilibre $\ce{2 H2O <=> H3O+ + OH-}$, la constante d'équilibre (produit ionique) s'écrit :
\[ K_e = [\ce{H3O+}] [\ce{OH-}] \quad (\text{en } \si{mol^2.L^{-2}}) \]
L'activité de l'eau pure (solvant) est prise égale à 1 (constante).
On définit $pK_e = -\log K_e$.
\textbf{Variation avec la température :} L'autoprotolyse est \textbf{endothermique} ($\Delta_r H > 0$). D'après le principe de Le Chatelier, $K_e$ \textbf{augmente} avec la température ($pK_e$ diminue).
\begin{center}
\begin{tabular}{|c|c|c|}\hline
\rowcolor{popblueL}
\textbf{Température} & \textbf{$K_e$} & \textbf{$pK_e$} \\\hline
\SI{25}{\celsius} & $1,0 \times 10^{-14}$ & 14,0 \\\hline
\SI{60}{\celsius} & $1,0 \times 10^{-13}$ & 13,0 \\\hline
\SI{100}{\celsius} & $5,5 \times 10^{-13}$ & 12,3 \\\hline
\end{tabular}
\end{center}
\end{aretenir}

\begin{exoresolu}[Autoprotolyse et neutralité électrique]
\textbf{Énoncé.} (a) Écrire l'équation-bilan de l'autoprotolyse de l'eau. (b) En déduire la relation entre $[\ce{H3O+}]$ et $[\ce{OH-}]$ dans l'eau pure à \SI{25}{\celsius}, puis calculer ces concentrations sachant que $pK_e = 14{,}0$.
\tcblower
\textbf{Solution.}
(a) L'eau est un ampholyte : elle appartient aux couples $\ce{H3O+}/\ce{H2O}$ et $\ce{H2O}/\ce{OH-}$. Le transfert de proton d'une molécule d'eau vers une autre s'écrit :
\[ \ce{2 H2O <=> H3O+ + OH-} \]
Cette réaction est très limitée : dans l'eau pure, seulement 2 molécules d'eau sur 550 millions environ sont ionisées à \SI{25}{\celsius}.

(b) D'après la stœchiométrie de l'équation, il se forme autant d'ions $\ce{H3O+}$ que d'ions $\ce{OH-}$ :
\[ [\ce{H3O+}] = [\ce{OH-}] \]
Or le produit ionique s'écrit $K_e = [\ce{H3O+}][\ce{OH-}]$. Donc :
\[ [\ce{H3O+}]^2 = K_e = 10^{-pK_e} = 10^{-14{,}0} \]
\[ [\ce{H3O+}] = [\ce{OH-}] = \sqrt{10^{-14}} = \SI{1,0e-7}{mol.L^{-1}} \]
\textbf{Conclusion :} dans l'eau pure à \SI{25}{\celsius}, $[\ce{H3O+}] = [\ce{OH-}] = \SI{1,0e-7}{mol.L^{-1}}$, ce qui correspond à $\mathrm{pH} = 7{,}0$ : l'eau pure est \cle{neutre}.
\end{exoresolu}

\coursec{Le pH : définition, calcul et caractère des solutions}

\begin{definition}[pH d'une solution aqueuse]
Le \cle{pH} (potentiel hydrogène) est défini par :
\[ \mathrm{pH} = -\log_{10} \left( \frac{[\ce{H3O+}]}{\SI{1}{mol.L^{-1}}} \right) = -\log_{10} [\ce{H3O+}] \]
où $[\ce{H3O+}]$ est la concentration molaire en ions oxonium (en \si{mol.L^{-1}}). Le pH est une grandeur \textbf{sans unité}.
Réciproquement : $[\ce{H3O+}] = 10^{-\mathrm{pH}}$ \si{mol.L^{-1}}.
\end{definition}

\begin{propriete}[Caractère acide, neutre, basique à une température donnée]
À une température $\theta$ où le produit ionique vaut $K_e$ ($pK_e = -\log K_e$) :
\begin{itemize}
\item Solution \cle{neutre} : $[\ce{H3O+}] = [\ce{OH-}] \iff \mathrm{pH} = \dfrac{pK_e}{2}$.
\item Solution \cle{acide} : $[\ce{H3O+}] > [\ce{OH-}] \iff \mathrm{pH} < \dfrac{pK_e}{2}$.
\item Solution \cle{basique} : $[\ce{H3O+}] < [\ce{OH-}] \iff \mathrm{pH} > \dfrac{pK_e}{2}$.
\end{itemize}
À \SI{25}{\celsius} ($pK_e = 14,0$) : Neutre $\mathrm{pH}=7,0$ ; Acide $\mathrm{pH}<7,0$ ; Basique $\mathrm{pH}>7,0$.
\end{propriete}

\begin{popfigure}
\begin{tikzpicture}
\begin{axis}[
    width=\linewidth, height=6cm,
    axis lines=middle,
    xlabel={pH}, ylabel={Concentration relative},
    xmin=0, xmax=14, ymin=0, ymax=1.1,
    xtick={0,2,4,6,7,8,10,12,14},
    ytick=\empty,
    domain=0:14, samples=200,
    clip=false,
    axis line style={popdark},
    tick style={popdark},
    xlabel style={anchor=north, font=\small},
]
% H3O+ curve: 10^-pH normalized roughly for visibility? No, log scale needed.
% Let's draw schematic zones instead.
\addplot[poppink, thick, domain=0:7] {1-x/7} node[right, pos=0.9, font=\tiny, poppink] {$[\ce{H3O+}]$ décroit};
\addplot[popblue, thick, domain=7:14] {(x-7)/7} node[right, pos=0.9, font=\tiny, popblue] {$[\ce{OH-}]$ croît};
% Neutrality line
\draw[popgreen, dashed, thick] (7,0) -- (7,1.1) node[above, font=\small, popgreen] {Neutralité $\mathrm{pH}=7$};
% Zones
\node[poppink, opacity=0.3, font=\large, rotate=90] at (3.5, 0.5) {ACIDE};
\node[popblue, opacity=0.3, font=\large, rotate=90] at (10.5, 0.5) {BASIQUE};
\node[popgreen, font=\small] at (7, 1.15) {$\mathrm{pH} = pK_e/2$};
\end{axis}
\end{tikzpicture}
\legende{Échelle de pH à \SI{25}{\celsius} : variation inverse des concentrations $[\ce{H3O+}]$ et $[\ce{OH-}]$.}
\end{popfigure}

\courssub{Méthode 2 : Trouver les concentrations des ions connaissant le pH et le pKe}

\begin{methode}[Calculer [\ce{H3O+}] et [\ce{OH-}] à partir du pH]
\begin{enumerate}
\item Relever le pH et la température (qui fixe la valeur de $pK_e$).
\item Calculer la concentration en ions oxonium : \[ [\ce{H3O+}] = 10^{-\mathrm{pH}} \quad (\si{mol.L^{-1}}) \]
\item Calculer le produit ionique : $K_e = 10^{-pK_e}$.
\item En déduire la concentration en ions hydroxyde : \[ [\ce{OH-}] = \frac{K_e}{[\ce{H3O+}]} = \frac{10^{-pK_e}}{10^{-\mathrm{pH}}} = 10^{\mathrm{pH}-pK_e} \]
\item Vérifier la cohérence : $[\ce{H3O+}] \times [\ce{OH-}] = K_e$ et comparer les deux concentrations pour conclure sur le caractère acide ou basique.
\end{enumerate}
\textbf{Réflexe :} La relation $[\ce{OH-}] = K_e/[\ce{H3O+}]$ est valable pour \textbf{toute} solution aqueuse, qu'elle soit acide, neutre ou basique.
\end{methode}

\begin{exoresolu}[Analyse d'un jus de bissap]
\textbf{Énoncé.} Le jus de bissap (infusion de fleurs d'hibiscus) a un pH de $2{,}8$ à \SI{25}{\celsius} ($pK_e = 14{,}0$). (a) Calculer $[\ce{H3O+}]$ et $[\ce{OH-}]$ dans ce jus. (b) Le jus est-il acide, neutre ou basique ? Justifier.
\tcblower
\textbf{Solution.}
(a) Concentration en ions oxonium :
\[ [\ce{H3O+}] = 10^{-\mathrm{pH}} = 10^{-2{,}8} = \SI{1,6e-3}{mol.L^{-1}} \]
Concentration en ions hydroxyde :
\[ [\ce{OH-}] = \frac{K_e}{[\ce{H3O+}]} = \frac{10^{-14{,}0}}{10^{-2{,}8}} = 10^{-11{,}2} = \SI{6,3e-12}{mol.L^{-1}} \]
Vérification : $[\ce{H3O+}][\ce{OH-}] = 1{,}6\times 10^{-3} \times 6{,}3\times 10^{-12} \approx 1{,}0\times 10^{-14} = K_e$. \checkmark

(b) Comparaison : $[\ce{H3O+}] = \SI{1,6e-3}{mol.L^{-1}} \gg [\ce{OH-}] = \SI{6,3e-12}{mol.L^{-1}}$.
De plus $\mathrm{pH} = 2{,}8 < 7{,}0 = pK_e/2$.
\textbf{Conclusion :} Le jus de bissap est \cle{acide}. C'est cette acidité (acides organiques : citrique, malique, hibiscus) qui lui donne sa saveur acidulée et fait virer ses anthocyanes au rouge vif.
\end{exoresolu}

\courssub{Méthode 3 : Savoir faire les approximations (espèces majoritaires, minoritaires, ultra-minoritaires)}

\begin{methode}[Hiérarchiser les espèces en solution]
\begin{enumerate}
\item Faire l'\textbf{inventaire} de toutes les espèces présentes : solvant $\ce{H2O}$, ions $\ce{H3O+}$ et $\ce{OH-}$ (toujours présents), ions ou molécules apportés par le soluté.
\item Calculer ou estimer chaque concentration (à partir de la concentration apportée $C$, du pH, de $K_e$).
\item Classer par comparaison des ordres de grandeur (règle du \textbf{facteur 100}) :
\begin{itemize}
\item \cle{Majoritaire} : concentration nettement la plus grande (souvent $\geq \SI{e-3}{mol.L^{-1}}$) ;
\item \cle{Minoritaire} : concentration faible mais non négligeable (facteur $< 100$ devant la majoritaire) ;
\item \cle{Ultra-minoritaire} : concentration négligeable devant toutes les autres (facteur $\geq 100$), que l'on \textbf{élimine des bilans de matière et de l'électroneutralité}.
\end{itemize}
\item Appliquer les approximations classiques (à vérifier \textbf{systématiquement}) :
\begin{itemize}
\item En milieu \textbf{acide} ($\mathrm{pH} \leq 6$ à \SI{25}{\celsius}) : $[\ce{OH-}] \leq 10^{-8} \ll [\ce{H3O+}]$, on néglige $\ce{OH-}$.
\item En milieu \textbf{basique} ($\mathrm{pH} \geq 8$ à \SI{25}{\celsius}) : $[\ce{H3O+}] \leq 10^{-8} \ll [\ce{OH-}]$, on néglige $\ce{H3O+}$.
\end{itemize}
\item Vérifier \textbf{a posteriori} que l'approximation est justifiée (rapport des concentrations $\geq 100$).
\end{enumerate}
\textbf{Réflexe :} Une approximation se \emph{pose} puis se \emph{vérifie} explicitement ; une approximation non vérifiée coûte des points au Bac.
\end{methode}

\begin{exoresolu}[Espèces présentes dans une solution d'acide chlorhydrique]
\textbf{Énoncé.} On dissout du chlorure d'hydrogène gazeux dans l'eau ; la solution obtenue a une concentration apportée $C = \SI{1,0e-3}{mol.L^{-1}}$ et un pH mesuré de $3{,}0$ à \SI{25}{\celsius}. (a) Faire l'inventaire des espèces chimiques présentes. (b) Classer ces espèces en majoritaires, minoritaires et ultra-minoritaires. (c) Écrire l'équation d'électroneutralité simplifiée.
\tcblower
\textbf{Solution.}
(a) $\ce{HCl}$ est un acide fort, totalement ionisé : $\ce{HCl + H2O -> H3O+ + Cl-}$. Espèces présentes : $\ce{H2O}$, $\ce{H3O+}$, $\ce{OH-}$, $\ce{Cl-}$. Il ne reste pratiquement pas de $\ce{HCl}$ moléculaire.

(b) Calculs des concentrations ioniques :
\[ [\ce{H3O+}] = 10^{-\mathrm{pH}} = \SI{1,0e-3}{mol.L^{-1}} \qquad [\ce{OH-}] = \frac{10^{-14}}{10^{-3}} = \SI{1,0e-11}{mol.L^{-1}} \]
\[ [\ce{Cl-}] = C = \SI{1,0e-3}{mol.L^{-1}} \]
Classement (facteur 100 = $10^2$) :
\begin{itemize}
\item Majoritaires : $\ce{H2O}$ (solvant, $\approx \SI{55,5}{mol.L^{-1}}$), $\ce{H3O+}$ et $\ce{Cl-}$ ($\SI{1,0e-3}{mol.L^{-1}}$) ;
\item Ultra-minoritaire : $\ce{OH-}$ ($\SI{1,0e-11}{mol.L^{-1}}$), négligeable devant $\ce{H3O+}$ et $\ce{Cl-}$ (rapport $10^8 \gg 100$).
\end{itemize}

(c) Électroneutralité exacte : $[\ce{H3O+}] = [\ce{Cl-}] + [\ce{OH-}]$. Comme $[\ce{OH-}] = 10^{-11} \ll [\ce{Cl-}] = 10^{-3}$ \si{mol.L^{-1}}, on néglige $\ce{OH-}$ :
\[ [\ce{H3O+}] \approx [\ce{Cl-}] = \SI{1,0e-3}{mol.L^{-1}} \]
Vérification : $10^{-3}/10^{-11} = 10^{8} \gg 100$, l'approximation est excellente.
\textbf{Conclusion :} Dans cette solution acide, les ions hydroxyde sont ultra-minoritaires et n'interviennent pas dans les bilans.
\end{exoresolu}

\courssub{Méthode 4 : Calculer le pH d'une solution et conclure sur son caractère}

\begin{methode}[Calculer un pH et conclure (acide, neutre, basique)]
\begin{enumerate}
\item Identifier la nature du soluté : acide fort (ionisation totale), base forte, ou solution neutre.
\item Pour un \textbf{acide fort} de concentration apportée $C$ (ex. \ce{HCl}, \ce{HNO3}, \ce{H2SO4} 1er proton) : si $C \gg 10^{-7}$ \si{mol.L^{-1}} (pratique : $C \geq \SI{e-6}{mol.L^{-1}}$), $[\ce{H3O+}] = C$, donc \[ \mathrm{pH} = -\log C \]
\item Pour une \textbf{base forte} de concentration apportée $C$ (ex. \ce{NaOH}, \ce{KOH}) : si $C \geq \SI{e-6}{mol.L^{-1}}$, $[\ce{OH-}] = C$, donc \[ [\ce{H3O+}] = \frac{K_e}{C} \quad \Rightarrow \quad \mathrm{pH} = pK_e + \log C \]
\item Conclure sur le caractère à la température $\theta$ considérée :
\begin{itemize}
\item $\mathrm{pH} < \dfrac{pK_e}{2}$ : solution \cle{acide} ;
\item $\mathrm{pH} = \dfrac{pK_e}{2}$ : solution \cle{neutre} ;
\item $\mathrm{pH} > \dfrac{pK_e}{2}$ : solution \cle{basique}.
\end{itemize}
À \SI{25}{\celsius}, $\dfrac{pK_e}{2} = 7{,}0$.
\item Vérifier la validité : les formules $\mathrm{pH} = -\log C$ et $\mathrm{pH} = pK_e + \log C$ supposent $10^{-6} \lesssim C \lesssim 10^{-1}$ \si{mol.L^{-1}}. En dehors, l'autoprotolyse de l'eau ou les effets d'activité perturbent le résultat.
\end{enumerate}
\textbf{Attention :} La neutralité n'est pas toujours à pH 7 ! À \SI{60}{\celsius}, $pK_e \approx 13{,}0$ et la neutralité est à $\mathrm{pH} = 6{,}5$.
\end{methode}

\begin{exoresolu}[pH d'une soude diluée et effet de la température]
\textbf{Énoncé.} On dispose d'une solution d'hydroxyde de sodium (soude, base forte) de concentration $C = \SI{2,0e-3}{mol.L^{-1}}$. (a) Calculer son pH à \SI{25}{\celsius} ($pK_e = 14{,}0$) et conclure sur son caractère. (b) À \SI{60}{\celsius}, $pK_e = 13{,}0$. Une eau pure à cette température a un pH de $6{,}5$. Est-elle acide ? Justifier.
\tcblower
\textbf{Solution.}
(a) La soude est une base forte, totalement dissociée : $\ce{NaOH -> Na+ + OH-}$, donc $[\ce{OH-}] = C = \SI{2,0e-3}{mol.L^{-1}}$.
\[ [\ce{H3O+}] = \frac{K_e}{[\ce{OH-}]} = \frac{10^{-14{,}0}}{2{,}0\times 10^{-3}} = 5{,}0\times 10^{-12}\ \si{mol.L^{-1}} \]
\[ \mathrm{pH} = -\log(5{,}0\times 10^{-12}) = 11{,}30 \]
Formule directe : $\mathrm{pH} = pK_e + \log C = 14{,}0 + \log(2{,}0\times 10^{-3}) = 14{,}0 - 2{,}70 = 11{,}30$. \checkmark
Comme $\mathrm{pH} = 11{,}30 > 7{,}0$, la solution est \cle{basique}.

(b) À \SI{60}{\celsius}, la neutralité correspond à $\mathrm{pH} = \dfrac{pK_e}{2} = \dfrac{13{,}0}{2} = 6{,}5$.
L'eau pure a $\mathrm{pH} = 6{,}5$ : on a toujours $[\ce{H3O+}] = [\ce{OH-}]$ car l'autoprotolyse produit les deux ions en quantités égales.
\textbf{Conclusion :} L'eau pure à \SI{60}{\celsius} est \cle{neutre} bien que son pH soit inférieur à 7. Le caractère acide ou basique se juge par comparaison à $pK_e/2$, jamais par la valeur absolue 7 hors de \SI{25}{\celsius}.
\end{exoresolu}

\coursec{Mesure du pH et indicateurs colorés}

\begin{definition}[Indicateur coloré et zone de virage]
Un \cle{indicateur coloré} est un couple acide/base \ce{HIn}/\ce{In-} dont les deux formes ont des \textbf{couleurs distinctes} dans le visible. Sa \cle{zone de virage} est l'intervalle de pH où les deux formes coexistent en proportions comparables (généralement $\mathrm{p}K_a \pm 1$). Hors de cette zone, une seule forme prédomine et la couleur est stable.
\end{definition}

\begin{popfigure}
\begin{tikzpicture}
\begin{axis}[
    width=\linewidth, height=5cm,
    axis lines=middle,
    xlabel={pH}, ylabel={Fraction de forme basique $\alpha_{\ce{In-}}$},
    xmin=0, xmax=14, ymin=0, ymax=1.1,
    xtick={0,2,4,6,7,8,10,12,14},
    ytick={0,0.5,1},
    domain=0:14, samples=200,
    axis line style={popdark},
    tick style={popdark},
    xlabel style={anchor=north, font=\small},
    ylabel style={anchor=east, font=\small},
    clip=false, % Permet d'afficher le nœud au-dessus de ymax
]
% Courbe sigmoïde pour indicateur pKa=7
\addplot[poppurple, very thick] {1/(1+10^(7-x))};
% Zone de virage (remplissage manuel sans bibliothèque fillbetween)
\addplot[poppurple!30, fill opacity=0.3, draw=none, domain=6:8, samples=100] 
    {1/(1+10^(7-x))} -- (axis cs:8,0) -- (axis cs:6,0) -- cycle;
% Repères visuels
\draw[poppurple, dashed] (6,0) -- (6,1) node[above, font=\tiny] {pH$_1$};
\draw[poppurple, dashed] (8,0) -- (8,1) node[above, font=\tiny] {pH$_2$};
\node[poppurple, font=\small] at (7, 1.15) {Zone de virage $[\mathrm{p}K_a-1; \mathrm{p}K_a+1]$};
\end{axis}
\end{tikzpicture}
\legende{Courbe de virage d'un indicateur (ex. BBT, $\mathrm{p}K_a \approx 7$) : transition de la forme acide (jaune) à la forme basique (bleu).}
\end{popfigure}

\courssub{Méthode 5 : Mesurer le pH et utiliser les indicateurs colorés}

\begin{methode}[Choisir et exploiter une méthode de mesure du pH]
\begin{enumerate}
\item \textbf{Papier pH (universel) :} Déposer une goutte de solution sur le papier à l'aide d'un agitateur en verre (ne jamais tremper le papier dans le flacon). Comparer la couleur à l'échelle de référence. Précision : $\pm 1$ unité pH $\rightarrow$ estimation rapide.
\item \textbf{Indicateurs colorés (solution) :} Ajouter 2 à 3 gouttes d'indicateur à un échantillon de solution (\SI{5}{mL} env.). La teinte observée, comparée à la \cle{zone de virage}, donne un encadrement du pH :
\begin{itemize}
\item Teinte acide (forme \ce{HIn}) $\Rightarrow \mathrm{pH} < \mathrm{pH}_{\text{virage inf}}$ ;
\item Teinte basique (forme \ce{In-}) $\Rightarrow \mathrm{pH} > \mathrm{pH}_{\text{virage sup}}$ ;
\item Teinte sensible (mélange) $\Rightarrow$ pH dans la zone de virage.
\end{itemize}
\item \textbf{pH-mètre (électrode de verre) :} Rincer l'électrode à l'eau distillée, l'essuyer délicatement (papier buvard, sans frotter le bulbe), étalonner avec deux solutions tampons (ex. pH 4,00 et 7,00 ou 7,00 et 9,00), puis plonger l'électrode dans la solution, agiter doucement et attendre la stabilisation. Précision : $\pm 0{,}05$ à $\pm 0{,}1$ unité pH.
\item \textbf{Préparer un indicateur naturel (chou rouge, bissap) :} Hacher le végétal (feuilles de chou rouge ou fleurs d'hibiscus séchées), le faire macérer \SI{15}{min} dans l'eau chaude (\SI{60}{\celsius}) ou bouillir \SI{5}{min}, filtrer sur coton ou papier filtre ; le filtrat (anthocyanes) change de couleur selon le pH (rouge acide, violet neutre, vert/jaune basique).
\end{enumerate}
\textbf{Réflexe Bac :} Pour une valeur précise de pH, seul le \cle{pH-mètre} étalonné convient ; les indicateurs (papier ou solution) ne donnent qu'un encadrement.
\end{methode}

\begin{experience}[Préparation et étalonnage d'un indicateur naturel (Bissap)]
\textbf{Matériel :} Fleurs d'hibiscus séchées (bissap), eau distillée, bécher, source de chauffe, entonnoir, papier filtre, tubes à essai, solutions tests (acide chlorhydrique dilué, eau, soude diluée, vinaigre, jus de citron, lessive).
\textbf{Protocole :}
\begin{enumerate}
\item Mettre \SI{5}{g} de fleurs dans \SI{100}{mL} d'eau, chauffer \SI{5}{min} à frémissement.
\item Filtrer. Recueillir le filtrat violet foncé (indicateur brut).
\item Répartir \SI{3}{mL} de filtrat dans 5 tubes. Ajouter \SI{1}{mL} de chaque solution test.
\item Observer et noter les couleurs. Comparer aux solutions témoins de pH connu si possible.
\end{enumerate}
\textbf{Observations typiques :} Rouge vif (pH 1-3, acide fort), Rose (pH 4-5, acide faible/vinaigre), Violet (pH 6-7, neutre), Bleu/Vert (pH 8-10, basique faible), Jaune (pH > 11, basique fort).
\textbf{Interprétation :} Les anthocyanes du bissap sont des indicateurs polyvalents (multi-zones de virage), utiles pour un diagnostic rapide en l'absence de pH-mètre.
\end{experience}

\begin{exoresolu}[Exploitation de la zone de virage d'un indicateur]
\textbf{Énoncé.} L'hélianthine est rouge pour $\mathrm{pH} < 3{,}1$, orange (teinte sensible) entre $3{,}1$ et $4{,}4$, et jaune pour $\mathrm{pH} > 4{,}4$. Le bleu de bromothymol (BBT) est jaune pour $\mathrm{pH} < 6{,}0$, vert entre $6{,}0$ et $7{,}6$, bleu pour $\mathrm{pH} > 7{,}6$. On teste un vinaigre de table : il colore l'hélianthine en jaune et le BBT en jaune. (a) Donner l'encadrement du pH du vinaigre. (b) Un élève mesure ensuite le pH au pH-mètre : $\mathrm{pH} = 2{,}9$. Ce résultat est-il compatible avec les tests ? (c) Pourquoi le pH-mètre est-il indispensable ici ?
\tcblower
\textbf{Solution.}
(a) Hélianthine jaune $\Rightarrow \mathrm{pH} > 4{,}4$. BBT jaune $\Rightarrow \mathrm{pH} < 6{,}0$.
Encadrement : $4{,}4 < \mathrm{pH} < 6{,}0$.

(b) La mesure au pH-mètre donne $\mathrm{pH} = 2{,}9$, ce qui est \textbf{hors} de l'encadrement $]4{,}4\,;\,6{,}0[$. Il y a donc une incohérence : soit les tests colorés ont été mal réalisés (indicateur périmé, lecture erronée de la teinte, contamination), soit le pH-mètre est mal étalonné. En refaisant les tests avec soin, l'hélianthine devrait être \emph{rouge} ($\mathrm{pH} = 2{,}9 < 3{,}1$) et non jaune.

(c) Les indicateurs colorés ne fournissent qu'un \textbf{encadrement} lié à leur zone de virage ; ils ne peuvent pas donner une valeur précise. Seul le pH-mètre étalonné fournit une valeur à $\pm 0{,}1$ unité près, indispensable pour identifier un dépassement de norme ou suivre un dosage.
\textbf{Conclusion :} Les indicateurs servent à estimer, le pH-mètre sert à mesurer.
\end{exoresolu}

\coursec{Préparation de solutions par dilution}

\begin{propriete}[Loi de conservation de la quantité de matière lors d'une dilution]
Lors d'une dilution, la quantité de matière de soluté $n$ est conservée :
\[ n = C_0 V_0 = C_f V_f \]
où $C_0, V_0$ concernent la solution mère (concentrée) et $C_f, V_f$ la solution fille (diluée).
Le \cle{facteur de dilution} $F$ est défini par :
\[ F = \frac{C_0}{C_f} = \frac{V_f}{V_0} \quad (F > 1, \text{sans unité}) \]
\textbf{Attention :} « Diluer 10 fois » signifie $F = 10$, donc $V_f = 10 V_0$ (et non $V_0 = 10 V_f$).
\end{propriete}

\courssub{Méthode 6 : Réaliser et calculer une dilution}

\begin{methode}[Réaliser une dilution au laboratoire]
\begin{enumerate}
\item \textbf{Calculer} $V_0 = \dfrac{C_f V_f}{C_0} = \dfrac{V_f}{F}$. Choisir $V_f$ et $F$ pour que $V_0$ corresponde à une \textbf{pipette jaugée} disponible (\SI{5}{mL}, \SI{10}{mL}, \SI{20}{mL}, \SI{25}{mL}, \SI{50}{mL}).
\item \textbf{Verrerie :} Pipette jaugée de volume $V_0$ + propipette (jamais à la bouche) ; Fiole jaugée de volume $V_f$ ; Bécher pour la solution mère ; Bouteille de lavage (eau distillée).
\item \textbf{Protocole :}
\begin{enumerate}
\item Rincer la pipette avec la solution mère (2-3 fois, \SI{1}{mL} chaque fois).
\item Prélever $V_0$ au-dessus du trait, ajuster le ménisque au trait (œil au niveau du trait).
\item Verser le contenu dans la fiole jaugée \textbf{contenant déjà environ 1/3 à 1/2 de $V_f$ en eau distillée} (sécurité : dilution partielle avant versement).
\item Rincer la pipette (pointe contre paroi) avec un peu d'eau distillée vers la fiole.
\item Compléter à l'eau distillée jusqu'au \textbf{trait de jauge} (bas du ménisque tangent au trait).
\item Boucher, homogénéiser (retourner 10-20 fois).
\end{enumerate}
\item \textbf{Sécurité :} Pour les acides/bases concentrés : porter lunettes, gants, blouse. Toujours verser \textbf{l'acide (ou la base) dans l'eau}, jamais l'inverse (éclaboussures, chaleur). En cas de projection : rincer \SI{15}{min} à l'eau courante.
\end{enumerate}
\end{methode}

\begin{exoresolu}[Préparation d'une solution d'acide chlorhydrique diluée]
\textbf{Énoncé.} Au laboratoire du lycée, on dispose d'une solution commerciale d'acide chlorhydrique de concentration $C_0 = \SI{1,0}{mol.L^{-1}}$. On veut préparer $V_f = \SI{250,0}{mL}$ d'une solution fille de concentration $C_f = \SI{2,0e-2}{mol.L^{-1}}$. (a) Calculer le facteur de dilution et le volume $V_0$ à prélever. (b) Décrire le protocole expérimental en précisant la verrerie. (c) Calculer le pH de la solution fille à \SI{25}{\celsius}.
\tcblower
\textbf{Solution.}
(a) Facteur de dilution :
\[ F = \frac{C_0}{C_f} = \frac{1{,}0}{2{,}0\times 10^{-2}} = 50 \]
Volume de solution mère à prélever :
\[ V_0 = \frac{C_f V_f}{C_0} = \frac{2{,}0\times 10^{-2} \times 250{,}0}{1{,}0} = \SI{5,0}{mL} \]
Vérification : $V_f/V_0 = 250{,}0/5{,}0 = 50 = F$. \checkmark

(b) Protocole :
\begin{enumerate}
\item Verser un peu de solution mère dans un bécher propre et sec.
\item Prélever $V_0 = \SI{5,0}{mL}$ à l'aide d'une \textbf{pipette jaugée de 5,0 mL} munie d'une propipette.
\item Introduire ce volume dans une \textbf{fiole jaugée de 250,0 mL} contenant déjà environ \SI{100}{mL} d'eau distillée.
\item Rincer la pointe de la pipette contre la paroi de la fiole avec un filet d'eau distillée.
\item Agiter, puis compléter avec de l'eau distillée jusqu'au trait de jauge (bas du ménisque au trait).
\item Boucher, retourner plusieurs fois pour homogénéiser.
\end{enumerate}

(c) $\ce{HCl}$ est un acide fort : $[\ce{H3O+}] = C_f = \SI{2,0e-2}{mol.L^{-1}}$.
\[ \mathrm{pH} = -\log(2{,}0\times 10^{-2}) = 1{,}70 \]
\textbf{Conclusion :} La solution fille, de pH $1{,}70$, est très acide ; la dilution par 50 a fait passer le pH de $0{,}00$ (solution mère) à $1{,}70$, soit une augmentation de $\log 50 \approx 1{,}70$ unité.
\end{exoresolu}

\begin{attention}[Les 5 erreurs qui coûtent des points au Bac]
\begin{enumerate}
\item \textbf{Confondre neutralité et pH = 7.} La neutralité correspond à $[\ce{H3O+}] = [\ce{OH-}]$, c'est-à-dire $\mathrm{pH} = pK_e/2$, qui ne vaut 7,0 qu'à \SI{25}{\celsius}. Toujours regarder la température et le $pK_e$ donnés dans l'énoncé.
\item \textbf{Oublier les ions $\ce{OH-}$ dans une solution acide (et inversement).} Toute solution aqueuse contient \emph{à la fois} $\ce{H3O+}$ et $\ce{OH-}$, liés par $K_e = [\ce{H3O+}][\ce{OH-}]$. Dire « il n'y a pas d'ions $\ce{OH-}$ » est faux ; dire « ils sont ultra-minoritaires » est correct.
\item \textbf{Se tromper de sens dans $[\ce{H3O+}] = 10^{-\mathrm{pH}}$.} Un pH \emph{petit} correspond à une concentration en $\ce{H3O+}$ \emph{grande}. Vérifie toujours l'ordre de grandeur : pH 3 $\Rightarrow \SI{e-3}{mol.L^{-1}}$, pas $10^{3}$.
\item \textbf{Poser une approximation sans la vérifier.} Négliger $\ce{OH-}$ devant $\ce{H3O+}$ exige un rapport de concentrations d'au moins 100. Écris explicitement la vérification : le correcteur l'attend.
\item \textbf{Mal gérer la dilution.} Erreurs typiques : inverser $F = C_0/C_f$, oublier que $V_0$ se prélève à la \emph{pipette jaugée} (pas à l'éprouvette), compléter la fiole \emph{avant} d'y verser le prélèvement, ou oublier d'homogénéiser après le trait de jauge. Et pour la sécurité : toujours verser l'acide concentré dans l'eau, jamais l'inverse.
\end{enumerate}
\end{attention}

\newpage

\coursec{Exercices}

\coursec{L'eau, un solvant exceptionnel}

\begin{definition}[Structure et polarité de la molécule d'eau]
La molécule d'eau \ce{H2O} possède une géométrie \textbf{coudée} (angle \ce{H-O-H} \approx \SI{104,5}{\degree}) due aux deux doublets non liants sur l'oxygène. L'oxygène étant plus électronégatif que l'hydrogène, les liaisons \ce{O-H} sont polarisées \ce{O^{\delta-}-H^{\delta+}}. La molécule possède un \textbf{moment dipolaire} permanent (\SI{1,85}{D}) : c'est une molécule \cle{polaires}.
\end{definition}

\begin{propriete}[Pouvoir dissolvant, dispersant et solvatant]
Grâce à sa polarité et à ses doublets non liants, l'eau interagit fortement avec les espèces ioniques et polaires :
\begin{itemize}
    \item \textbf{Solvatation (hydratation) :} Les ions en solution sont entourés d'une coquille de molécules d'eau orientées (oxygène vers les cations, hydrogènes vers les anions). Cela stabilise les ions et explique la dissolution des solides ioniques (ex. \ce{NaCl}, \ce{KNO3}).
    \item \textbf{Pouvoir dispersant :} L'eau disperse les molécules polaires (sucre, éthanol, acides) par liaisons hydrogène.
    \item \textbf{Pouvoir ionisant :} L'eau favorise la dissociation des acides et bases (ex. \ce{HCl -> H+ + Cl-}, solvaté en \ce{H3O+ + Cl-}).
\end{itemize}
\end{propriete}

\begin{savaistu}[L'eau au Cameroun]
L'eau de source (Mont Cameroun, Monts Mandara) contient des ions minéraux (\ce{Ca^2+}, \ce{Mg^2+}, \ce{HCO3-}) dissous par érosion des roches. L'eau de mer (Kribi, Limbé) est une solution saline complexe (\approx \SI{0,6}{mol.L^{-1}} en \ce{NaCl}). La potabilité exige un pH entre 6,5 et 8,5 (norme OMS/Ministère de la Santé).
\end{savaistu}

\begin{experience}[Mise en évidence de la polarité de l'eau]
\textbf{Dispositif :} Burette avec eau distillée, tige de verre (ou ballon en plastique) frottée sur laine (chargée électrostatiquement).
\textbf{Protocole :} Ouvrir le robinet pour obtenir un filet d'eau fin. Approcher la tige chargée du filet sans le toucher.
\textbf{Observation :} Le filet d'eau se dévie vers la tige chargée.
\textbf{Interprétation :} Les molécules d'eau, dipolaires, s'orientent dans le champ électrique de la tige (côté \ce{\delta-} vers charge +, côté \ce{\delta+} vers charge -), prouvant leur polarité.
\end{experience}

\coursec{Acides et bases selon Brønsted}

\begin{definition}[Acide et base de Brønsted]
\begin{itemize}
    \item Un \textbf{acide} (donneur de proton) est une espèce chimique capable de \textbf{céder un ion \ce{H+}} (proton).
    \item Une \textbf{base} (accepteur de proton) est une espèce chimique capable de \textbf{capter un ion \ce{H+}}.
\end{itemize}
Une réaction acido-basique est un \textbf{transfert de proton} d'un acide vers une base.
\end{definition}

\begin{propriete}[Couples acide/base]
Un \textbf{couple acide/base} s'écrit \ce{AH}/\ce{A-} (ou \ce{BH+}/\ce{B}). L'acide \ce{AH} a perdu un proton pour donner sa base conjuguée \ce{A-}. La base \ce{A-} a gagné un proton pour donner son acide conjugué \ce{AH}.
\[ \ce{AH + H2O <=> A- + H3O+} \quad \text{ou} \quad \ce{B + H2O <=> BH+ + OH-} \]
L'eau est \textbf{ampholyte} (ou amphotère) : elle peut agir comme acide (donne \ce{H+} à une base) ou comme base (capte \ce{H+} d'un acide).
\end{propriete}

\begin{exemplebox}[Couples de l'eau]
\begin{center}
\begin{tabularx}{\linewidth}{|l|X|X|}
\hline
\rowcolor{popblueL} \textbf{Réaction} & \textbf{Acide} & \textbf{Base} \\
\hline
\ce{H2O + H2O <=> H3O+ + OH-} & \ce{H2O} (donne \ce{H+}) & \ce{H2O} (capte \ce{H+}) \\
\ce{HCl + H2O -> H3O+ + Cl-} & \ce{HCl} & \ce{H2O} \\
\ce{NH3 + H2O <=> NH4+ + OH-} & \ce{H2O} & \ce{NH3} \\
\hline
\end{tabularx}
\end{center}
Les couples sont : \ce{H3O+}/\ce{H2O}, \ce{H2O}/\ce{OH-}, \ce{HCl}/\ce{Cl-}, \ce{NH4+}/\ce{NH3}.
\end{exemplebox}

\begin{aretenir}[Force des acides/bases]
Un acide est \textbf{fort} si sa réaction avec l'eau est \textbf{totale} (flèche \ce{->}), \textbf{faible} si elle est \textbf{limitée} (flèche \ce{<=>}). La base conjuguée d'un acide fort est une base très faible (négligeable), et inversement.
\end{aretenir}

\coursec{Autoprotolyse de l'eau et produit ionique}

\begin{definition}[Autoprotolyse de l'eau]
L'eau pure réagit avec elle-même : une molécule donne un proton (acide) à une autre molécule (base).
\[ \ce{2 H2O(l) <=> H3O+(aq) + OH-(aq)} \]
C'est une transformation \textbf{très limitée} à température ambiante.
\end{definition}

\begin{propriete}[Produit ionique de l'eau \texorpdfstring{$K_e$}{Ke}]
À l'équilibre, le produit des concentrations molaires des ions formés est constant à température donnée :
\[ K_e = [\ce{H3O+}] \times [\ce{OH-}] \]
\begin{itemize}
    \item À \SI{25}{\celsius} : $K_e = \num{1,0e-14}$ (souvent noté $K_w$).
    \item $K_e$ \textbf{augmente} avec la température (la réaction est \textbf{endothermique}).
\end{itemize}
\end{propriete}

\begin{definition}[p$K_e$]
\[ \text{p}K_e = -\log_{10} K_e \]
À \SI{25}{\celsius} : $\text{p}K_e = 14,0$.
Relation fondamentale : $\text{pH} + \text{pOH} = \text{p}K_e$ avec $\text{pOH} = -\log_{10}[\ce{OH-}]$.
\end{definition}

\begin{exemplebox}[Variation de $K_e$ avec la température]
\begin{center}
\begin{tabular}{|c|c|c|}
\hline
\rowcolor{popblueL} \textbf{Température (\SI{}{\celsius})} & \textbf{$K_e$} & \textbf{p$K_e$} \\
\hline
0 & \num{0,11e-14} & 14,9 \\
25 & \num{1,0e-14} & 14,0 \\
60 & \num{1,0e-13} & 13,0 \\
100 & \num{5,5e-13} & 12,3 \\
\hline
\end{tabular}
\end{center}
\emph{Interprétation :} L'augmentation de $K_e$ avec $T$ signifie que l'équilibre se décale vers les produits (ions) : la réaction absorbe de la chaleur (endothermique).
\end{exemplebox}

\coursec{Le pH : définition, calcul et caractère des solutions}

\begin{definition}[pH d'une solution aqueuse]
Pour une solution aqueuse diluée ($[\ce{H3O+}] < \SI{1}{mol.L^{-1}}$), le pH est défini par :
\[ \text{pH} = -\log_{10} \left( \frac{[\ce{H3O+}]}{\SI{1}{mol.L^{-1}}} \right) \]
Inversement : $[\ce{H3O+}] = 10^{-\text{pH}}$ (en \SI{}{mol.L^{-1}}).
\end{definition}

\begin{propriete}[Caractère acide, neutre, basique à \SI{25}{\celsius}]
\begin{center}
\begin{tabularx}{\linewidth}{|l|X|X|X|}
\hline
\rowcolor{popblueL} \textbf{Caractère} & \textbf{Condition sur $[\ce{H3O+}]$} & \textbf{Condition sur pH} & \textbf{Relation $[\ce{H3O+}]$ vs $[\ce{OH-}]$} \\
\hline
Acide & $> \SI{1,0e-7}{mol.L^{-1}}$ & $< 7$ & $[\ce{H3O+}] > [\ce{OH-}]$ \\
\hline
Neutre & $= \SI{1,0e-7}{mol.L^{-1}}$ & $= 7$ & $[\ce{H3O+}] = [\ce{OH-}]$ \\
\hline
Basique & $< \SI{1,0e-7}{mol.L^{-1}}$ & $> 7$ & $[\ce{H3O+}] < [\ce{OH-}]$ \\
\hline
\end{tabularx}
\end{center}
\end{propriete}

\begin{attention}[Effet de la température sur la neutralité]
À une autre température, la neutralité correspond à $\text{pH} = \frac{\text{p}K_e}{2}$ (ex. pH = 6,5 à \SI{60}{\celsius}). Une solution de pH 6,5 à \SI{60}{\celsius} est \textbf{neutre}, pas acide !
\end{attention}

\begin{methode}[Calcul du pH et des concentrations]
\begin{enumerate}
    \item Identifier les espèces sources d'ions \ce{H3O+} ou \ce{OH-} (acide fort, base forte, eau).
    \item Écrire l'équation d'électroneutralité : $\sum n_i [\text{cation}_i] = \sum n_j [\text{anion}_j]$.
    \item Utiliser $K_e = [\ce{H3O+}][\ce{OH-}]$.
    \item Faire les \textbf{approximations} (espèces majoritaires $> 10^{-3}$ M, minoritaires $10^{-3}$ à $10^{-7}$ M, ultra-minoritaires $< 10^{-7}$ M) pour simplifier les calculs.
    \item Calculer pH ou concentrations.
\end{enumerate}
\end{methode}

\begin{exemplebox}[Solution d'acide chlorhydrique \SI{1,0e-3}{mol.L^{-1}} à \SI{25}{\celsius}]
\ce{HCl} fort $\rightarrow$ dissociation totale : $[\ce{H3O+}] \approx C_{\ce{HCl}} = \SI{1,0e-3}{mol.L^{-1}}$.
$[\ce{OH-}] = \frac{K_e}{[\ce{H3O+}]} = \frac{\num{1,0e-14}}{\num{1,0e-3}} = \SI{1,0e-11}{mol.L^{-1}}$.
$\text{pH} = -\log(\num{1,0e-3}) = 3,00$.
Espèces : \ce{H2O} (majoritaire), \ce{H3O+}, \ce{Cl-} (mineures), \ce{OH-} (ultra-minoritaire).
Électroneutralité : $[\ce{H3O+}] = [\ce{Cl-}] + [\ce{OH-}] \approx [\ce{Cl-}]$.
\end{exemplebox}

\coursec{Mesure du pH et indicateurs colorés}

\begin{propriete}[Méthodes de mesure]
\begin{itemize}
    \item \textbf{Papier pH (ou bandelette) :} Imprégné d'un mélange d'indicateurs. Précision \approx \pm 0,5 unité de pH. Rapide, qualitatif.
    \item \textbf{pH-mètre :} Mesure de la différence de potentiel d'une électrode de verre sensible aux \ce{H3O+}. Précision \approx \pm 0,01 à \pm 0,1 unité. Nécessite étalonnage (solutions tampons pH 4, 7, 10). Méthode de référence pour les dosages.
    \item \textbf{Indicateurs colorés :} Mélanges acide/base conjugués (\ce{HIn}/\ce{In-}) dont les formes acide et base ont des couleurs distinctes.
\end{itemize}
\end{propriete}

\begin{definition}[Zone de virage d'un indicateur]
Intervalle de pH où l'indicateur change de couleur. Généralement $\text{p}K_a \pm 1$.
\begin{itemize}
    \item pH $< \text{p}K_a - 1$ : couleur de la forme acide (\ce{HIn}).
    \item pH $> \text{p}K_a + 1$ : couleur de la forme basique (\ce{In-}).
    \item Dans la zone : mélange des deux couleurs.
\end{itemize}
\end{definition}

\begin{experience}[Préparation d'un indicateur naturel : jus de bissap (Hibiscus sabdariffa)]
\textbf{Matériel :} Fleurs de bissap séchées (marché local), eau distillée, bécher, source de chaleur, filtre, tubes à essai, solutions étalons (pH 2, 7, 12).
\textbf{Protocole :}
\begin{enumerate}
    \item Faire bouillir \SI{50}{g} de fleurs dans \SI{500}{mL} d'eau \SI{10}{min}.
    \item Filtrer. Le filtrat rouge foncé est l'indicateur (conserver au frais).
    \item Verser \SI{2}{mL} de jus dans 3 tubes. Ajouter \SI{2}{mL} de chaque solution étalon.
\end{enumerate}
\textbf{Observations :}
\begin{center}
\begin{tabular}{|c|c|c|}
\hline
\rowcolor{popblueL} \textbf{pH étalon} & \textbf{Milieu} & \textbf{Couleur observée} \\
\hline
2,0 & Acide & \textbf{Rouge} \\
7,0 & Neutre & \textbf{Violet} \\
12,0 & Basique & \textbf{Vert-bleu} \\
\hline
\end{tabular}
\end{center}
\textbf{Interprétation :} Les anthocyanes (pigments) changent de structure selon le pH. Cet indicateur permet une estimation visuelle du pH (précision \approx \pm 1).
\end{experience}

\begin{savaistu}[Indicateurs au laboratoire et dans la vie courante]
Hélianthine (rouge/jaune, pH 3,1-4,4), Bleu de bromothymol BBT (jaune/bleu, pH 6,0-7,6), Phénolphtaléine (incolore/rose, pH 8,2-10,0). Le chou rouge (anthocyanes) donne des couleurs similaires au bissap. Le vinaigre (pH $\approx 2,5$) vire le bissap au rouge ; le savon (pH $\approx 9-10$) le vire au vert.
\end{savaistu}

\coursec{Préparation de solutions par dilution}

\begin{definition}[Facteur de dilution]
Prélever un volume $V_0$ de solution mère (concentration $C_0$) et compléter au volume $V$ avec le solvant (eau).
\[ F = \frac{V}{V_0} = \frac{C_0}{C} \]
$F > 1$. La concentration diminue : $C = \frac{C_0}{F}$.
\end{definition}

\begin{methode}[Protocole de dilution (ex. fiole jaugée \SI{100}{mL} à partir de \SI{10}{mL} solution mère)]
\begin{enumerate}
    \item Rincer la pipette jaugée \SI{10}{mL} avec la solution mère (déchet).
    \item Prélever \SI{10,0}{mL} de solution mère (pipette + poire à pipeter \textbf{jamais à la bouche}).
    \item Verser dans la fiole jaugée \SI{100}{mL} (rincée à l'eau distillée).
    \item Rincer la pipette avec un peu d'eau distillée dans la fiole.
    \item Compléter à l'eau distillée jusqu'au trait de jauge (bas du ménisque au niveau du trait, œil au niveau du trait).
    \item Bouchonner et homogénéiser (retourner la fiole plusieurs fois).
\end{enumerate}
\end{methode}

\begin{attention}[Sécurité et pièges]
\begin{itemize}
    \item \textbf{Acide concentré :} Toujours verser l'acide \textbf{dans l'eau} (exothermique violent sinon). Port d'EPI (lunettes, blouse, gants).
    \item \textbf{Verreerie :} Pipette jaugée (un trait), fiole jaugée (un trait), burette (graduée). Bécher et erlenmeyer \textbf{ne sont pas jaugeables} (précision \approx 5\%).
    \item \textbf{Dilution d'un acide fort :} pH augmente de 1 unité pour une dilution 10 fois. Dilution 100 fois $\rightarrow$ pH +2.
\end{itemize}
\end{attention}

\courssub{Je vérifie mes connaissances}

\begin{exercice}[QCM : l'eau et le pH]
Pour chaque question, choisir la (ou les) bonne(s) réponse(s) en justifiant brièvement.
\begin{enumerate}
\item À \SI{25}{\celsius}, le produit ionique de l'eau vaut :
\begin{enumerate}
    \item[a.] $K_e = 10^{-7}$ \quad b. $K_e = 10^{-14}$ \quad c. $K_e = 14$ \quad d. $K_e = 10^{14}$
\end{enumerate}
\item Une solution de pH $= 4$ est :
\begin{enumerate}
    \item[a.] neutre \quad b. acide \quad c. basique \quad d. impossible à \SI{25}{\celsius}
\end{enumerate}
\item La molécule d'eau est :
\begin{enumerate}
    \item[a.] linéaire et apolaire \quad b. coudée et polaire \quad c. coudée et apolaire \quad d. linéaire et polaire
\end{enumerate}
\item Lorsqu'on dilue 10 fois une solution de pH $= 2$ (acide fort), le pH de la solution diluée vaut :
\begin{enumerate}
    \item[a.] 1 \quad b. 2 \quad c. 3 \quad d. 12
\end{enumerate}
\item La concentration en ions oxonium d'une solution de pH $= 3{,}5$ vaut :
\begin{enumerate}
    \item[a.] $\SI{3,5}{mol.L^{-1}}$ \quad b. $\SI{3,2e-4}{mol.L^{-1}}$ \quad c. $\SI{3,2e-11}{mol.L^{-1}}$ \quad d. $\SI{0,35}{mol.L^{-1}}$
\end{enumerate}
\end{enumerate}
\end{exercice}

\begin{exercice}[Vrai ou faux ?]
Répondre par \textbf{vrai} ou \textbf{faux} en justifiant chaque réponse par une phrase ou un calcul.
\begin{enumerate}
\item À \SI{25}{\celsius}, toute solution de pH égal à 7 est neutre.
\item À \SI{60}{\celsius}, une eau pure de pH égal à $6{,}5$ est acide.
\item Une solution basique ne contient aucun ion \ce{H3O+}.
\item L'autoprotolyse de l'eau est une transformation très limitée.
\item Plus le pH d'une solution acide est petit, plus sa concentration en ions \ce{H3O+} est grande.
\item Le papier pH permet une mesure du pH plus précise que le pH-mètre.
\end{enumerate}
\end{exercice}

\begin{exercice}[Texte à trous]
Compléter le texte suivant avec les mots ou expressions de la liste : \emph{oxonium}, \emph{hydroxyde}, \emph{produit ionique}, \emph{autoprotolyse}, \emph{acide}, \emph{base}, \emph{transfert de proton}, \emph{neutre}, \emph{polaire}, \emph{solvatation}.

\medskip
La molécule d'eau est coudée et \textbf{polaire}, ce qui explique son pouvoir dissolvant vis-à-vis des solides ioniques : les ions sont entourés de molécules d'eau, c'est le phénomène de \textbf{solvatation}. Selon Brønsted, un \textbf{acide} est une espèce capable de céder un proton \ce{H+}, tandis qu'une \textbf{base} est une espèce capable d'en capter un. Une réaction acido-basique est un \textbf{transfert de proton} entre deux couples. L'eau pure est le siège de la réaction d'\textbf{autoprotolyse}, qui produit des ions \textbf{oxonium} \ce{H3O+} et des ions \textbf{hydroxyde} \ce{OH-} en quantités égales : l'eau pure est donc \textbf{neutre}. Le \textbf{produit ionique} de l'eau est la constante $K_e = [\ce{H3O+}][\ce{OH-}]$.
\end{exercice}

\begin{exercice}[Définitions et équations]
\begin{enumerate}
\item Donner la définition du pH d'une solution aqueuse et préciser sa condition de validité sur la concentration.
\item Écrire l'équation-bilan de l'autoprotolyse de l'eau et donner l'expression du produit ionique $K_e$.
\item Définir le p$K_e$ et donner sa valeur à \SI{25}{\celsius}.
\item Citer deux couples acide/base de l'eau et préciser le rôle (acide ou base) joué par \ce{H2O} dans chacun.
\end{enumerate}
\end{exercice}

\courssub{Je m'entraîne}

\begin{exercice}[Solution d'acide chlorhydrique]
Une solution aqueuse d'acide chlorhydrique \ce{(H3O+ + Cl-)} a un pH égal à $2{,}7$ à \SI{25}{\celsius}.
\begin{enumerate}
\item Calculer les concentrations $[\ce{H3O+}]$ et $[\ce{OH-}]$ de cette solution.
\item En déduire la concentration $C$ de la solution d'acide chlorhydrique.
\item Dresser la liste des espèces chimiques présentes en solution et classer-les en espèces majoritaires, minoritaires et ultra-minoritaires.
\item Vérifier que l'équation d'électroneutralité de la solution est satisfaite.
\end{enumerate}
\emph{Donnée à \SI{25}{\celsius} :} $K_e = 1{,}0 \times 10^{-14}$.
\end{exercice}

\begin{exercice}[Dilution d'une solution de soude]
On prélève un volume $V_0 = \SI{20,0}{mL}$ d'une solution d'hydroxyde de sodium (soude) de concentration $C_0 = \SI{0,050}{mol.L^{-1}}$, que l'on introduit dans une fiole jaugée de $V = \SI{200,0}{mL}$ et que l'on complète avec de l'eau distillée jusqu'au trait de jauge.
\begin{enumerate}
\item Définir le facteur de dilution $F$ et calculer sa valeur.
\item Calculer la concentration $C$ de la solution diluée obtenue.
\item Calculer la concentration en ions \ce{OH-}, puis celle en ions \ce{H3O+}, et en déduire le pH de la solution diluée à \SI{25}{\celsius}.
\item Décrire en quelques lignes le protocole de cette dilution en citant la verrerie utilisée.
\end{enumerate}
\emph{Donnée à \SI{25}{\celsius} :} p$K_e = 14{,}0$.
\end{exercice}

\begin{exercice}[L'eau pure à \SI{60}{\celsius}]
À la température de \SI{60}{\celsius}, le produit ionique de l'eau est tel que p$K_e = 13{,}0$.
\begin{enumerate}
\item Rappeler la relation entre $[\ce{H3O+}]$ et $[\ce{OH-}]$ dans l'eau pure, à toute température. Justifier.
\item Calculer le pH de l'eau pure à \SI{60}{\celsius}.
\item Un élève affirme : « À \SI{60}{\celsius}, l'eau pure a un pH inférieur à 7, donc elle est devenue acide. » Cette affirmation est-elle correcte ? Expliquer pourquoi l'eau pure reste neutre.
\item Comment évolue la valeur de $K_e$ lorsque la température augmente ? Que peut-on en conclure sur le caractère endothermique ou exothermique de l'autoprotolyse de l'eau ?
\end{enumerate}
\end{exercice}

\begin{exercice}[Le pH du sang]
Lors d'une analyse médicale effectuée dans un hôpital de Yaoundé, on mesure le pH sanguin d'un patient. Chez un individu en bonne santé, le pH du sang est compris entre $7{,}35$ et $7{,}45$ à \SI{37}{\celsius}. On prendra p$K_e = 13{,}6$ à \SI{37}{\celsius}.
\begin{enumerate}
\item Calculer les concentrations en ions \ce{H3O+} correspondant aux deux valeurs limites du pH sanguin normal.
\item Calculer les concentrations en ions \ce{OH-} correspondantes.
\item Le sang est-il un milieu acide, neutre ou basique ? Justifier.
\item On parle d'acidose lorsque le pH sanguin devient inférieur à $7{,}35$. Dans ce cas, la concentration en ions \ce{H3O+} a-t-elle augmenté ou diminué ? Justifier.
\end{enumerate}
\end{exercice}

\courssub{Je me prépare au Bac}

\begin{exercice}[Type Bac — Trois solutions et indicateurs colorés]
Au laboratoire d'un lycée de Douala, on dispose de trois flacons sans étiquette contenant des solutions aqueuses incolores A, B et C, de pH respectifs $3{,}0$, $7{,}0$ et $11{,}0$ (dans le désordre), à \SI{25}{\celsius}. On dispose de trois indicateurs colorés dont les zones de virage sont rappelées ci-dessous.

\begin{center}
\begin{tabularx}{\linewidth}{|l|X|X|X|}
\hline
\rowcolor{popblueL} \textbf{Indicateur} & \textbf{Teinte acide} & \textbf{Zone de virage} & \textbf{Teinte basique} \\
\hline
Hélianthine & rouge & $3{,}1$ -- $4{,}4$ & jaune orangé \\
\hline
Bleu de bromothymol (BBT) & jaune & $6{,}0$ -- $7{,}6$ & bleu \\
\hline
Phénolphtaléine & incolore & $8{,}2$ -- $10{,}0$ & rose fuchsia \\
\hline
\end{tabularx}
\end{center}

\begin{enumerate}
\item Calculer les concentrations $[\ce{H3O+}]$ et $[\ce{OH-}]$ de chacune des trois solutions.
\item Indiquer, en le justifiant, la teinte prise par chacun des trois indicateurs dans chacune des trois solutions. On présentera les résultats sous forme d'un tableau.
\item Montrer qu'un seul indicateur ne suffit pas toujours à distinguer les trois solutions. Proposer alors une méthode expérimentale utilisant \textbf{un seul} indicateur, choisi judicieusement, puis une mesure complémentaire simple, pour identifier sans ambiguïté les trois flacons.
\item La solution de pH $= 3{,}0$ est une solution d'acide chlorhydrique. Calculer sa concentration $C$, puis le volume $V_0$ de cette solution qu'il faut prélever pour préparer $\SI{500,0}{mL}$ d'une solution de pH $= 4{,}0$ par dilution. Décrire le protocole et la verrerie.
\end{enumerate}
\emph{Donnée à \SI{25}{\celsius} :} $K_e = 1{,}0 \times 10^{-14}$.
\end{exercice}

\begin{exercice}[Type Bac — Le jus de bissap, indicateur naturel]
Les fleurs d'hibiscus (\emph{bissap}), très utilisées au Cameroun pour préparer une boisson rafraîchissante, contiennent des pigments dont la couleur dépend du pH : le jus est \textbf{rouge} en milieu acide, \textbf{violet} vers la neutralité et \textbf{vert-bleu} en milieu basique. Un groupe d'élèves de Terminale C de Bafoussam réalise les expériences suivantes à \SI{25}{\celsius}.

\textbf{Partie A : étalonnage de l'indicateur naturel.} Les élèves disposent de trois solutions étalons de pH connus : pH $= 2{,}0$ ; pH $= 7{,}0$ ; pH $= 12{,}0$.
\begin{enumerate}
\item Prévoir la couleur prise par le jus de bissap dans chaque solution étalon.
\item Calculer $[\ce{H3O+}]$ et $[\ce{OH-}]$ pour chacune des trois solutions étalons.
\end{enumerate}

\textbf{Partie B : analyse d'un vinaigre de table.} Quelques gouttes de jus de bissap ajoutées à un échantillon de vinaigre donnent une teinte rouge, identique à celle obtenue avec la solution de pH $= 2{,}0$.
\begin{enumerate}
\setcounter{enumi}{2}
\item En déduire une valeur approchée du pH du vinaigre et de sa concentration en ions \ce{H3O+}.
\item Le vinaigre contient de l'acide éthanoïque \ce{CH3COOH}, un acide faible ($\text{p}K_a(\ce{CH3COOH}/\ce{CH3COO-}) = 4{,}8$). Écrire l'équation de la réaction de l'acide éthanoïque avec l'eau.
\item La concentration en acide éthanoïque du vinaigre est-elle égale à $[\ce{H3O+}]$ ? Justifier en comparant le cas d'un acide fort et celui d'un acide faible.
\end{enumerate}

\textbf{Partie C : comparaison des méthodes de mesure.}
\begin{enumerate}
\setcounter{enumi}{5}
\item Comparer la précision des trois méthodes de mesure du pH : jus de bissap, papier pH et pH-mètre étalonné. Conclure sur la méthode à privilégier pour un dosage.
\end{enumerate}
\emph{Données à \SI{25}{\celsius} :} $K_e = 1{,}0 \times 10^{-14}$ ; $M(\ce{CH3COOH}) = \SI{60,0}{g.mol^{-1}}$.
\end{exercice}

\begin{exercice}[Type Bac — Préparation d'une solution d'acide éthanoïque]
Un laborantin d'un centre de santé de Garoua doit préparer $V = \SI{250,0}{mL}$ d'une solution S d'acide éthanoïque de pH voisin de $3{,}0$, à partir d'une solution commerciale S$_0$ de concentration $C_0 = \SI{1,0}{mol.L^{-1}}$. On travaille à \SI{25}{\celsius}.

\textbf{Partie A : la solution à préparer.}
\begin{enumerate}
\item Calculer la concentration en ions \ce{H3O+} d'une solution de pH $= 3{,}0$.
\item En supposant dans un premier temps que l'acide est fort, calculer la concentration $C$ de la solution S à préparer.
\item Calculer le volume $V_0$ de solution commerciale S$_0$ à prélever et le facteur de dilution $F$.
\item Décrire précisément le protocole de préparation de la solution S, en citant toute la verrerie utilisée et en précisant les règles de sécurité (manipulation d'un acide concentré).
\end{enumerate}

\textbf{Partie B : vérification du pH.}
\begin{enumerate}
\setcounter{enumi}{4}
\item Le laborantin mesure le pH de la solution S avec du papier pH : il lit pH $\approx 3$. Puis il ajoute quelques gouttes de jus de bissap (rouge en milieu acide, violet vers la neutralité, vert-bleu en milieu basique) dans un échantillon de S. Quelle teinte observe-t-il ?
\item En réalité, l'acide éthanoïque est un acide \textbf{faible} ($\text{p}K_a = 4{,}8$). Le pH réel de la solution S est-il exactement égal à $3{,}0$ ? Dire, en le justifiant qualitativement, s'il est légèrement supérieur ou inférieur à la valeur calculée à la question 2.
\item Expliquer comment le laborantin peut ajuster expérimentalement le pH de sa solution pour obtenir la valeur souhaitée, en utilisant le pH-mètre.
\end{enumerate}
\emph{Données à \SI{25}{\celsius} :} p$K_e = 14{,}0$ ; $M(\ce{CH3COOH}) = \SI{60,0}{g.mol^{-1}}$.
\end{exercice}



\newpage

\coursec{Corrigés des exercices}

\begin{corrige}[Exercice 1 — QCM : l'eau et le pH]
\begin{enumerate}
    \item \textbf{Réponse b.} À \SI{25}{\celsius}, dans l'eau pure : $[\ce{H3O+}] = [\ce{OH-}] = \SI{1,0e-7}{mol.L^{-1}}$, donc $K_e = [\ce{H3O+}][\ce{OH-}] = 10^{-7} \times 10^{-7} = \num{1,0e-14}$.
    \item \textbf{Réponse b.} À \SI{25}{\celsius}, la neutralité correspond à pH $= 7$. Ici pH $= 4 < 7$, donc $[\ce{H3O+}] > [\ce{OH-}]$ : la solution est \textbf{acide}.
    \item \textbf{Réponse b.} La molécule \ce{H2O} est \textbf{coudée} (angle $\approx \SI{104,5}{\degree}$, deux doublets non liants sur O) et \textbf{polaire} (moment dipolaire $\SI{1,85}{D}$ non nul).
    \item \textbf{Réponse c.} Diluer 10 fois divise $[\ce{H3O+}]$ par 10, donc le pH augmente d'une unité : pH $= 2 + 1 = 3$.
    \item \textbf{Réponse b.} $[\ce{H3O+}] = 10^{-\text{pH}} = 10^{-3,5} = \num{3,16e-4} \approx \SI{3,2e-4}{mol.L^{-1}}$.
\end{enumerate}
\end{corrige}

\begin{corrige}[Exercice 2 — Vrai ou faux ?]
\begin{enumerate}
    \item \textbf{Vrai.} À \SI{25}{\celsius}, p$K_e = 14$, donc la neutralité ($[\ce{H3O+}] = [\ce{OH-}]$) correspond exactement à pH $= \frac{\text{p}K_e}{2} = 7$.
    \item \textbf{Faux.} À \SI{60}{\celsius}, p$K_e = 13{,}0$, donc la neutralité correspond à pH $= 6{,}5$. Dans l'eau pure, $[\ce{H3O+}] = [\ce{OH-}]$ à toute température : elle reste \textbf{neutre}, même si son pH est inférieur à 7.
    \item \textbf{Faux.} Le produit ionique impose $[\ce{H3O+}] = \frac{K_e}{[\ce{OH-}]} \neq 0$. Toute solution aqueuse contient des ions \ce{H3O+} et \ce{OH-} ; dans une base, \ce{H3O+} est simplement ultra-minoritaire.
    \item \textbf{Vrai.} $K_e = \num{1,0e-14}$ est extrêmement petit : l'équilibre $\ce{2 H2O <=> H3O+ + OH-}$ est très déplacé vers l'eau moléculaire. Dans \SI{1}{L} d'eau pure, seulement $10^{-7}$ mol d'eau est ionisée.
    \item \textbf{Vrai.} pH $= -\log[\ce{H3O+}]$ est une fonction décroissante de $[\ce{H3O+}]$ : plus le pH est petit, plus $[\ce{H3O+}] = 10^{-\text{pH}}$ est grande.
    \item \textbf{Faux.} Le papier pH donne une précision d'environ $\pm 0{,}5$ unité, alors qu'un pH-mètre étalonné atteint $\pm 0{,}01$ à $\pm 0{,}1$ unité. C'est le pH-mètre la méthode la plus précise.
\end{enumerate}
\end{corrige}

\begin{corrige}[Exercice 3 — Texte à trous]
La molécule d'eau est coudée et \textbf{polaire}, ce qui explique son pouvoir dissolvant vis-à-vis des solides ioniques : les ions sont entourés de molécules d'eau, c'est le phénomène de \textbf{solvatation}. Selon Brønsted, un \textbf{acide} est une espèce capable de céder un proton \ce{H+}, tandis qu'une \textbf{base} est une espèce capable d'en capter un. Une réaction acido-basique est un \textbf{transfert de proton} entre deux couples. L'eau pure est le siège de la réaction d'\textbf{autoprotolyse}, qui produit des ions \textbf{oxonium} \ce{H3O+} et des ions \textbf{hydroxyde} \ce{OH-} en quantités égales : l'eau pure est donc \textbf{neutre}. Le \textbf{produit ionique} de l'eau est la constante $K_e = [\ce{H3O+}][\ce{OH-}]$.

\medskip
\emph{Tous les mots de la liste sont utilisés une et une seule fois.}
\end{corrige}

\begin{corrige}[Exercice 4 — Définitions et équations]
\begin{enumerate}
    \item Le pH d'une solution aqueuse diluée est défini par :
    \[ \text{pH} = -\log_{10}\left(\frac{[\ce{H3O+}]}{\SI{1}{mol.L^{-1}}}\right) \]
    Cette définition n'est valable que pour des solutions \textbf{diluées}, c'est-à-dire $[\ce{H3O+}] \leq \SI{1}{mol.L^{-1}}$ (sinon le pH deviendrait négatif et la loi n'est plus applicable).
    \item Équation-bilan de l'autoprotolyse de l'eau :
    \[ \ce{2 H2O(l) <=> H3O+(aq) + OH-(aq)} \]
    Expression du produit ionique : $K_e = [\ce{H3O+}] \times [\ce{OH-}]$ (constante à température donnée).
    \item $\text{p}K_e = -\log_{10} K_e$. À \SI{25}{\celsius}, $K_e = \num{1,0e-14}$, donc $\text{p}K_e = 14{,}0$.
    \item Couple \ce{H3O+}/\ce{H2O} : \ce{H2O} joue le rôle de \textbf{base} (elle capte un proton pour donner \ce{H3O+}). Couple \ce{H2O}/\ce{OH-} : \ce{H2O} joue le rôle d'\textbf{acide} (elle cède un proton pour donner \ce{OH-}). L'eau est donc un \textbf{ampholyte}.
\end{enumerate}
\end{corrige}

\begin{corrige}[Exercice 5 — Solution d'acide chlorhydrique]
\begin{enumerate}
    \item Concentration en ions oxonium :
    \[ [\ce{H3O+}] = 10^{-\text{pH}} = 10^{-2,7} = \SI{2,0e-3}{mol.L^{-1}} \]
    Concentration en ions hydroxyde :
    \[ [\ce{OH-}] = \frac{K_e}{[\ce{H3O+}]} = \frac{\num{1,0e-14}}{\num{2,0e-3}} = \SI{5,0e-12}{mol.L^{-1}} \]
    \item \ce{HCl} est un acide \textbf{fort} : sa réaction avec l'eau est totale :
    \[ \ce{HCl + H2O -> H3O+ + Cl-} \]
    Donc $C = [\ce{H3O+}] = \SI{2,0e-3}{mol.L^{-1}}$.
    \item Inventaire des espèces :
    \begin{center}
    \begin{tabularx}{\linewidth}{|l|X|X|}
    \hline
    \rowcolor{popblueL} \textbf{Catégorie} & \textbf{Espèces} & \textbf{Concentration (mol$\cdot$L$^{-1}$)} \\
    \hline
    Majoritaire & \ce{H2O} & $\approx 55{,}5$ \\
    \hline
    Minoritaires & \ce{H3O+}, \ce{Cl-} & $\approx \num{2,0e-3}$ \\
    \hline
    Ultra-minoritaire & \ce{OH-} & $\num{5,0e-12}$ \\
    \hline
    \end{tabularx}
    \end{center}
    \item Électroneutralité : $[\ce{H3O+}] = [\ce{Cl-}] + [\ce{OH-}]$.
    Vérification : $[\ce{Cl-}] + [\ce{OH-}] = \num{2,0e-3} + \num{5,0e-12} \approx \num{2,0e-3} = [\ce{H3O+}]$. L'équation est satisfaite : l'apport des ions \ce{OH-} (autoprotolyse) est négligeable devant celui de l'acide.
\end{enumerate}
\end{corrige}

\begin{corrige}[Exercice 6 — Dilution d'une solution de soude]
\begin{enumerate}
    \item Le facteur de dilution est le rapport du volume final au volume prélevé (ou des concentrations) :
    \[ F = \frac{V}{V_0} = \frac{C_0}{C} = \frac{200{,}0}{20{,}0} = 10 \]
    \item $C = \dfrac{C_0}{F} = \dfrac{\num{5,0e-2}}{10} = \SI{5,0e-3}{mol.L^{-1}}$.
    \item La soude est une base \textbf{forte} : $\ce{NaOH -> Na+ + OH-}$ (totale), donc $[\ce{OH-}] = C = \SI{5,0e-3}{mol.L^{-1}}$.
    \[ [\ce{H3O+}] = \frac{K_e}{[\ce{OH-}]} = \frac{\num{1,0e-14}}{\num{5,0e-3}} = \SI{2,0e-12}{mol.L^{-1}} \]
    \[ \text{pH} = -\log(\num{2,0e-12}) = 12 - \log 2 = 11{,}7 \]
    Vérification : pOH $= -\log(\num{5,0e-3}) = 2{,}30$ et pH $= 14{,}0 - 2{,}30 = 11{,}70$. La solution est bien basique.
    \item \textbf{Protocole :} rincer la pipette jaugée de \SI{20}{mL} avec la solution mère ; prélever \SI{20,0}{mL} à l'aide d'une poire à pipeter ; introduire dans une fiole jaugée de \SI{200}{mL} contenant déjà un peu d'eau distillée ; compléter à l'eau distillée jusqu'au trait de jauge (bas du ménisque, œil au niveau du trait) ; boucher et homogénéiser en retournant la fiole plusieurs fois. \textbf{Sécurité :} lunettes, blouse, gants ; toujours verser la solution concentrée dans l'eau (dilution exothermique).
\end{enumerate}
\end{corrige}

\begin{corrige}[Exercice 7 — L'eau pure à \SI{60}{\celsius}]
\begin{enumerate}
    \item Dans l'eau pure, les seuls ions proviennent de l'autoprotolyse $\ce{2 H2O <=> H3O+ + OH-}$, dont les coefficients stœchiométriques sont égaux (1 et 1). À toute température : $[\ce{H3O+}] = [\ce{OH-}]$.
    \item $K_e = [\ce{H3O+}][\ce{OH-}] = [\ce{H3O+}]^2$, donc :
    \[ [\ce{H3O+}] = \sqrt{K_e} = 10^{-\text{p}K_e/2} = 10^{-6,5} = \SI{3,2e-7}{mol.L^{-1}} \]
    \[ \text{pH} = -\log(10^{-6,5}) = 6{,}5 \]
    \item L'affirmation est \textbf{fausse}. Le caractère acide ou basique se définit par la \textbf{comparaison} de $[\ce{H3O+}]$ et $[\ce{OH-}]$, pas par la valeur 7 du pH (qui n'est le pH de neutralité qu'à \SI{25}{\celsius}). Or dans l'eau pure à \SI{60}{\celsius}, $[\ce{H3O+}] = [\ce{OH-}] = \SI{3,2e-7}{mol.L^{-1}}$ : l'eau reste \textbf{neutre}. Le pH de neutralité vaut $\frac{\text{p}K_e}{2} = 6{,}5$ à cette température.
    \item $K_e$ \textbf{augmente} quand la température augmente (p$K_e$ diminue : 14,9 à \SI{0}{\celsius}, 13,0 à \SI{60}{\celsius}). D'après la loi de Van 't Hoff (ou le principe de Le Chatelier), une élévation de température déplace l'équilibre dans le sens \textbf{endothermique} : la formation des ions est donc favorisée par la chaleur, ce qui prouve que l'autoprotolyse de l'eau est une transformation \textbf{endothermique}.
\end{enumerate}
\end{corrige}

\begin{corrige}[Exercice 8 — Le pH du sang]
\begin{enumerate}
    \item $[\ce{H3O+}] = 10^{-\text{pH}}$ :
    \begin{itemize}
        \item pH $= 7{,}35$ : $[\ce{H3O+}] = 10^{-7,35} = \SI{4,5e-8}{mol.L^{-1}}$ ;
        \item pH $= 7{,}45$ : $[\ce{H3O+}] = 10^{-7,45} = \SI{3,5e-8}{mol.L^{-1}}$.
    \end{itemize}
    \item $K_e = 10^{-13,6} = \num{2,5e-14}$ et $[\ce{OH-}] = \dfrac{K_e}{[\ce{H3O+}]}$ :
    \begin{itemize}
        \item pH $= 7{,}35$ : $[\ce{OH-}] = \dfrac{\num{2,5e-14}}{\num{4,5e-8}} = \SI{5,6e-7}{mol.L^{-1}}$ ;
        \item pH $= 7{,}45$ : $[\ce{OH-}] = \dfrac{\num{2,5e-14}}{\num{3,5e-8}} = \SI{7,1e-7}{mol.L^{-1}}$.
    \end{itemize}
    \item Dans les deux cas, $[\ce{OH-}] > [\ce{H3O+}]$ (environ 12 à 20 fois plus) : le sang est un milieu \textbf{légèrement basique}. Remarque : à \SI{37}{\celsius}, la neutralité correspond à pH $= \frac{13{,}6}{2} = 6{,}8$, et le sang (pH \approx 7,4) est bien au-dessus.
    \item Lors d'une acidose, le pH \textbf{diminue} (pH $< 7{,}35$). Comme pH $= -\log[\ce{H3O+}]$ est une fonction décroissante, la concentration en ions \ce{H3O+} a \textbf{augmenté} (elle dépasse $\SI{4,5e-8}{mol.L^{-1}}$).
\end{enumerate}
\end{corrige}

\begin{corrige}[Exercice 9 — Type Bac : trois solutions et indicateurs colorés]
\textbf{Barème indicatif (sur 10) :} question 1 : 2 pts ; question 2 : 3 pts ; question 3 : 2 pts ; question 4 : 3 pts.

\begin{enumerate}
    \item À \SI{25}{\celsius}, $[\ce{H3O+}] = 10^{-\text{pH}}$ et $[\ce{OH-}] = \dfrac{K_e}{[\ce{H3O+}]}$ :
    \begin{center}
    \begin{tabular}{|c|c|c|}
    \hline
    \rowcolor{popblueL} \textbf{pH} & \textbf{$[\ce{H3O+}]$ (mol$\cdot$L$^{-1}$)} & \textbf{$[\ce{OH-}]$ (mol$\cdot$L$^{-1}$)} \\
    \hline
    3,0 & $\num{1,0e-3}$ & $\num{1,0e-11}$ \\
    7,0 & $\num{1,0e-7}$ & $\num{1,0e-7}$ \\
    11,0 & $\num{1,0e-11}$ & $\num{1,0e-3}$ \\
    \hline
    \end{tabular}
    \end{center}
    \item On compare chaque pH à la zone de virage (pH $<$ borne basse : teinte acide ; pH $>$ borne haute : teinte basique ; dans la zone : teinte sensible) :
    \begin{center}
    \begin{tabular}{|l|c|c|c|}
    \hline
    \rowcolor{popblueL} \textbf{Indicateur} & \textbf{pH 3,0} & \textbf{pH 7,0} & \textbf{pH 11,0} \\
    \hline
    Hélianthine & rouge ($3{,}0 < 3{,}1$) & jaune orangé ($7{,}0 > 4{,}4$) & jaune orangé ($11{,}0 > 4{,}4$) \\
    \hline
    BBT & jaune ($3{,}0 < 6{,}0$) & vert (zone de virage) & bleu ($11{,}0 > 7{,}6$) \\
    \hline
    Phénolphtaléine & incolore ($3{,}0 < 8{,}2$) & incolore ($7{,}0 < 8{,}2$) & rose fuchsia ($11{,}0 > 10{,}0$) \\
    \hline
    \end{tabular}
    \end{center}
    \item L'hélianthine ne distingue pas les solutions de pH 7,0 et 11,0 (jaune orangé dans les deux cas) ; la phénolphtaléine ne distingue pas pH 3,0 et 7,0 (incolore dans les deux cas). Un seul de ces indicateurs ne suffit donc pas.
    \textbf{Méthode proposée :} utiliser l'\textbf{hélianthine} : le flacon qui vire au \textbf{rouge} est la solution de pH 3,0. Pour départager les deux autres (jaune orangé), on mesure le pH au \textbf{papier pH} (ou au pH-mètre) : celui qui indique pH $\approx 7$ est la solution neutre, l'autre (pH $\approx 11$) la solution basique. (Le BBT seul aurait aussi suffi : jaune, vert, bleu.)
    \item \ce{HCl} est un acide fort : $C = [\ce{H3O+}] = 10^{-3,0} = \SI{1,0e-3}{mol.L^{-1}}$.
    Solution fille : pH $= 4{,}0 \Rightarrow C' = \SI{1,0e-4}{mol.L^{-1}}$, $V' = \SI{500,0}{mL}$.
    Conservation de la quantité de matière : $C V_0 = C' V'$, donc :
    \[ V_0 = \frac{C' V'}{C} = \frac{\num{1,0e-4} \times 500{,}0}{\num{1,0e-3}} = \SI{50,0}{mL} \]
    (Facteur de dilution $F = 10$, cohérent avec pH $+1$.)
    \textbf{Protocole :} prélever \SI{50,0}{mL} de solution mère à la pipette jaugée (poire à pipeter), verser dans une fiole jaugée de \SI{500,0}{mL} contenant déjà de l'eau distillée, compléter jusqu'au trait de jauge, boucher et homogénéiser. \textbf{Sécurité :} lunettes, blouse, gants ; toujours verser l'acide dans l'eau.
\end{enumerate}

\begin{attention}[Points de vigilance]
Ne pas confondre $[\ce{H3O+}]$ et $[\ce{OH-}]$ dans les solutions basiques ; justifier chaque teinte par une inégalité sur le pH ; vérifier l'homogénéité des unités (mL partout) dans le calcul de $V_0$.
\end{attention}
\end{corrige}

\begin{corrige}[Exercice 10 — Type Bac : le jus de bissap, indicateur naturel]
\textbf{Barème indicatif (sur 10) :} Partie A : 3 pts ; Partie B : 4 pts ; Partie C : 3 pts.

\textbf{Partie A}
\begin{enumerate}
    \item pH $= 2{,}0$ (acide) : \textbf{rouge} ; pH $= 7{,}0$ (neutre) : \textbf{violet} ; pH $= 12{,}0$ (basique) : \textbf{vert-bleu}.
    \item Avec $[\ce{H3O+}] = 10^{-\text{pH}}$ et $[\ce{OH-}] = K_e/[\ce{H3O+}]$ :
    \begin{center}
    \begin{tabular}{|c|c|c|}
    \hline
    \rowcolor{popblueL} \textbf{pH} & \textbf{$[\ce{H3O+}]$ (mol$\cdot$L$^{-1}$)} & \textbf{$[\ce{OH-}]$ (mol$\cdot$L$^{-1}$)} \\
    \hline
    2,0 & $\num{1,0e-2}$ & $\num{1,0e-12}$ \\
    7,0 & $\num{1,0e-7}$ & $\num{1,0e-7}$ \\
    12,0 & $\num{1,0e-12}$ & $\num{1,0e-2}$ \\
    \hline
    \end{tabular}
    \end{center}
\end{enumerate}

\textbf{Partie B}
\begin{enumerate}
    \setcounter{enumi}{2}
    \item La teinte rouge est identique à celle de l'étalon pH $= 2{,}0$ : on en déduit pH(vinaigre) $\approx 2{,}0$, soit $[\ce{H3O+}] \approx \SI{1,0e-2}{mol.L^{-1}}$.
    \item Réaction de l'acide éthanoïque (acide faible) avec l'eau :
    \[ \ce{CH3COOH(aq) + H2O(l) <=> CH3COO-(aq) + H3O+(aq)} \]
    \item \textbf{Non.} Pour un acide \textbf{fort}, la dissociation est totale et $C = [\ce{H3O+}]$. Pour l'acide éthanoïque, acide \textbf{faible}, la réaction est limitée : seule une fraction des molécules cède son proton, donc $[\ce{H3O+}] < C$. La conservation de la matière s'écrit $C = [\ce{CH3COOH}] + [\ce{CH3COO-}]$, avec $[\ce{CH3COO-}] \approx [\ce{H3O+}]$, d'où $C > [\ce{H3O+}]$. Le vinaigre est donc bien plus concentré en acide que $\SI{1,0e-2}{mol.L^{-1}}$.
\end{enumerate}

\textbf{Partie C}
\begin{enumerate}
    \setcounter{enumi}{5}
    \item \begin{itemize}
        \item Jus de bissap : méthode \textbf{qualitative} (comparaison visuelle de teintes), précision de l'ordre de $\pm 1$ à $2$ unités de pH ;
        \item Papier pH : méthode \textbf{semi-quantitative}, précision $\approx \pm 0{,}5$ unité ;
        \item pH-mètre étalonné : méthode \textbf{quantitative} (potentiométrie), précision $\approx \pm 0{,}01$ à $0{,}1$ unité.
    \end{itemize}
    \textbf{Conclusion :} pour un dosage acido-basique, il faut privilégier le \textbf{pH-mètre étalonné} ; le papier pH convient à un contrôle rapide et le jus de bissap à une estimation domestique ou pédagogique.
\end{enumerate}

\begin{attention}[Points de vigilance]
Un indicateur naturel donne une \emph{estimation} et non une mesure ; ne jamais écrire $C = [\ce{H3O+}]$ pour un acide faible ; penser à justifier le caractère limité de la réaction par la double flèche \ce{<=>}.
\end{attention}
\end{corrige}

\begin{corrige}[Exercice 11 — Type Bac : préparation d'une solution d'acide éthanoïque]
\textbf{Barème indicatif (sur 10) :} Partie A : 5 pts ; Partie B : 5 pts.

\textbf{Partie A}
\begin{enumerate}
    \item $[\ce{H3O+}] = 10^{-\text{pH}} = 10^{-3,0} = \SI{1,0e-3}{mol.L^{-1}}$.
    \item Sous l'hypothèse (provisoire) d'un acide fort : $C = [\ce{H3O+}] = \SI{1,0e-3}{mol.L^{-1}}$.
    \item Facteur de dilution :
    \[ F = \frac{C_0}{C} = \frac{1{,}0}{\num{1,0e-3}} = 1000 \]
    Volume à prélever :
    \[ V_0 = \frac{V}{F} = \frac{250{,}0}{1000} = \SI{0,250}{mL} \]
    Ce volume est trop petit pour une pipette jaugée classique : on procède donc par \textbf{dilutions successives}.
    \item \textbf{Protocole (dilution en deux étapes) :}
    \begin{enumerate}
        \item \textbf{Étape 1 ($F_1 = 10$) :} prélever \SI{10,0}{mL} de S$_0$ à la pipette jaugée, verser dans une fiole jaugée de \SI{100,0}{mL}, compléter à l'eau distillée, homogénéiser : on obtient S$_1$ de concentration $C_1 = \SI{0,10}{mol.L^{-1}}$.
        \item \textbf{Étape 2 ($F_2 = 100$) :} prélever \SI{2,50}{mL} de S$_1$ à la pipette jaugée, verser dans la fiole jaugée de \SI{250,0}{mL}, compléter au trait de jauge, boucher et homogénéiser : $C = \SI{1,0e-3}{mol.L^{-1}}$ (au total $F = F_1 \times F_2 = 1000$).
    \end{enumerate}
    \textbf{Sécurité :} port de lunettes, blouse et gants ; verser \textbf{toujours l'acide concentré dans l'eau} (jamais l'inverse, risque de projections exothermiques) ; pipetage à la poire, jamais à la bouche.
\end{enumerate}

\textbf{Partie B}
\begin{enumerate}
    \setcounter{enumi}{4}
    \item La solution S est acide (pH $\approx 3$) : le jus de bissap prend une teinte \textbf{rouge}.
    \item Le pH réel n'est pas exactement 3,0 : il est \textbf{légèrement supérieur}. En effet, l'acide éthanoïque est un acide faible : sa réaction avec l'eau est limitée ($\ce{CH3COOH + H2O <=> CH3COO- + H3O+}$), donc $[\ce{H3O+}] < C = \SI{1,0e-3}{mol.L^{-1}}$, ce qui donne pH $> 3{,}0$. Estimation : $[\ce{H3O+}] \approx \sqrt{K_a \cdot C} = \sqrt{10^{-4,8} \times 10^{-3}} = 10^{-3,9}$, soit pH $\approx 3{,}9$.
    \item \textbf{Ajustement expérimental :} étalonner le pH-mètre avec des solutions tampons (pH 4 et 7) ; plonger l'électrode dans S sous agitation ; si le pH mesuré est supérieur à 3,0, ajouter \textbf{goutte à goutte} un peu de solution plus concentrée (S$_1$ ou S$_0$ diluée) en suivant le pH ; si le pH est inférieur à 3,0, ajouter de l'eau distillée goutte à goutte ; arrêter dès que l'afficheur indique pH $= 3{,}0$ de façon stable.
\end{enumerate}

\begin{attention}[Points de vigilance]
Un facteur de dilution de 1000 ne se réalise jamais en une seule étape (volume prélevé trop faible et imprécis) : penser aux dilutions en cascade. Ne pas oublier que l'hypothèse « acide fort » surestime $[\ce{H3O+}]$ pour l'acide éthanoïque : le pH réel est toujours plus élevé que le pH calculé avec cette hypothèse.
\end{attention}
\end{corrige}

\newpage

\coursec{Fiche bilan}

\courssub{Carte mentale de la leçon}

\begin{popfigure}
\begin{tikzpicture}[
    every node/.style={align=center, font=\small, inner sep=4pt},
    concept/.style={rectangle, rounded corners=8pt, draw, thick, minimum width=2.8cm, minimum height=1.2cm, font=\bfseries\small, text=white},
    subconcept/.style={rectangle, rounded corners=5pt, draw, thin, minimum width=2.2cm, minimum height=0.9cm, font=\scriptsize, text=white, opacity=0.95},
    line/.style={draw, thick, -{Stealth[length=2mm]}, shorten >=2pt, shorten <=2pt}
  ]
  % Central node
  \node[concept, fill=popblue, minimum width=3.5cm, minimum height=1.5cm] (center) {Eau, autoprotolyse\\ et pH};
  
  % Level 1 nodes (6 branches)
  \node[concept, fill=poporange] (n1) at (90:4.2) {L'eau, solvant\\ exceptionnel};
  \node[concept, fill=popgreen] (n2) at (30:4.2) {Acides et bases\\ de Brønsted};
  \node[concept, fill=poppurple] (n3) at (-30:4.2) {Autoprotolyse\\ de l'eau};
  \node[concept, fill=popblue!70!black] (n4) at (-90:4.2) {Le pH};
  \node[concept, fill=poppink] (n5) at (-150:4.2) {Mesure du pH};
  \node[concept, fill=popgold!80!black] (n6) at (150:4.2) {Dilution};
  
  % Connections center -> level 1
  \foreach \n in {n1,n2,n3,n4,n5,n6} {
    \draw[line, popmuted] (center) -- (\n);
  }
  
  % Level 2 nodes (children)
  % Branch 1: Eau
  \node[subconcept, fill=poporange] (n11) at ($(n1)+(90:2.2)$) {Molécule coudée,\\ polaire};
  \node[subconcept, fill=poporange] (n12) at ($(n1)+(150:2.2)$) {Dissout, disperse,\\ solvate, ionise};
  \draw[line, poporange!50!white] (n1) -- (n11);
  \draw[line, poporange!50!white] (n1) -- (n12);
  
  % Branch 2: Brønsted
  \node[subconcept, fill=popgreen] (n21) at ($(n2)+(30:2.2)$) {Acide : donneur\\ de \ce{H+}};
  \node[subconcept, fill=popgreen] (n22) at ($(n2)+(-30:2.2)$) {Couples \ce{AH}/\ce{A-}\\ conjugués};
  \draw[line, popgreen!50!white] (n2) -- (n21);
  \draw[line, popgreen!50!white] (n2) -- (n22);
  
  % Branch 3: Autoprotolyse
  \node[subconcept, fill=poppurple] (n31) at ($(n3)+(-30:2.2)$) {$K_e=\SI{1,0e-14}{}$\\ à \SI{25}{\celsius}};
  \node[subconcept, fill=poppurple] (n32) at ($(n3)+(-90:2.2)$) {$pK_e=14$\\ $K_e \uparrow$ si $T \uparrow$};
  \draw[line, poppurple!50!white] (n3) -- (n31);
  \draw[line, poppurple!50!white] (n3) -- (n32);
  
  % Branch 4: pH
  \node[subconcept, fill=popblue!70!black] (n41) at ($(n4)+(-90:2.2)$) {$pH=-\log[\ce{H3O+}]$};
  \node[subconcept, fill=popblue!70!black] (n42) at ($(n4)+(-150:2.2)$) {Acide / Neutre\\ / Basique};
  \draw[line, popblue!50!white] (n4) -- (n41);
  \draw[line, popblue!50!white] (n4) -- (n42);
  
  % Branch 5: Mesure
  \node[subconcept, fill=poppink] (n51) at ($(n5)+(-150:2.2)$) {pH-mètre, papier pH};
  \node[subconcept, fill=poppink] (n52) at ($(n5)+(150:2.2)$) {Indicateurs colorés\\ (zone de virage)};
  \draw[line, poppink!50!white] (n5) -- (n51);
  \draw[line, poppink!50!white] (n5) -- (n52);
  
  % Branch 6: Dilution
  \node[subconcept, fill=popgold!80!black] (n61) at ($(n6)+(150:2.2)$) {$C_iV_i=C_fV_f$};
  \node[subconcept, fill=popgold!80!black] (n62) at ($(n6)+(90:2.2)$) {Facteur $F=C_i/C_f$\\ $=V_f/V_i$};
  \draw[line, popgold!50!white] (n6) -- (n61);
  \draw[line, popgold!50!white] (n6) -- (n62);
\end{tikzpicture}
\legende{Carte mentale de synthèse : les six idées-forces de la leçon.}
\end{popfigure}

\courssub{L'essentiel à retenir}

\begin{definition}[L'eau et les couples acide/base de Brønsted]
\textbf{Structure de la molécule d'eau} : géométrie \textbf{coudée} (angle $\approx \SI{104,5}{\degree}$), liaison O--H polarisée ($\Delta\chi \approx 1,4$). Dipôle net $\neq 0$ $\rightarrow$ molécule \textbf{polaire}.
\begin{center}
\chemfig{H-[:30]O(-[:-30]H)} \quad $\rightarrow$ \quad \chemfig{H-[:30]O^{\delta-}(-[:-30]H^{\delta+})}
\end{center}
\textbf{Rôles de l'eau} : \emph{solvant} (dissout), \emph{dispersant} (colloïdes), \emph{solvantant} (entoure les ions par solvatation), \emph{ionisant} (favorise la dissociation des électrolytes).

\medskip
\textbf{Acide de Brønsted} (\cle{AH}) : espèce chimique capable de \textbf{céder} un proton \ce{H+}.
\textbf{Base de Brønsted} (\cle{B}) : espèce chimique capable de \textbf{capter} un proton \ce{H+}.
\textbf{Couple acide/base conjugué} : deux espèces liées par le gain ou la perte d'un proton, noté \ce{AH}/\ce{A-} (acide/base) ou \ce{BH+}/\ce{B} (base/acide).
\begin{itemize}
  \item Exemple : \ce{CH3COOH}/\ce{CH3COO-} (acide éthanoïque / éthanoate) ; \ce{NH4+}/\ce{NH3} (ammonium / ammoniac).
  \item L'eau est \textbf{amphotère} : acide (\ce{H2O}/\ce{HO-}) et base (\ce{H3O+}/\ce{H2O}).
\end{itemize}
\end{definition}

\begin{propriete}[Autoprotolyse de l'eau, produit ionique et pH]
\textbf{Équation-bilan de l'autoprotolyse} (réaction limitée) :
\[ \ce{2 H2O <=> H3O+ + HO-} \]
\textbf{Produit ionique de l'eau} ($K_e$) : pour \textbf{toute} solution aqueuse, à température donnée,
\[ K_e = [\ce{H3O+}] \times [\ce{HO-}] \]
À \SI{25}{\celsius} : $K_e = \SI{1,0e-14}{}$ (sans unité) et $pK_e = -\log K_e = 14,00$.
\textbf{Variation avec la température} : la réaction est endothermique $\Rightarrow$ $K_e$ \textbf{augmente} quand $T$ augmente (ex. $K_e \approx \SI{5,5e-14}{}$ à \SI{50}{\celsius}). La neutralité ($[\ce{H3O+}]=[\ce{HO-}]$) n'est à $pH=7$ qu'à \SI{25}{\celsius}.

\medskip
\textbf{pH d'une solution aqueuse} :
\[ pH = -\log_{10}\left( \frac{[\ce{H3O+}]}{\SI{1}{mol.L^{-1}}} \right) \quad \Leftrightarrow \quad [\ce{H3O+}] = 10^{-pH} \times \SI{1}{mol.L^{-1}} \]
\textbf{Concentration en ions hydroxyde} :
\[ [\ce{HO-}] = \frac{K_e}{[\ce{H3O+}]} = 10^{pH-pK_e} \times \SI{1}{mol.L^{-1}} \]
\textbf{Caractère de la solution à \SI{25}{\celsius}} :
\begin{center}
\begin{tabular}{|c|c|c|c|}\hline
\rowcolor{popblueL} \textbf{Caractère} & \textbf{Comparaison $[\ce{H3O+}]$ vs $[\ce{HO-}]$} & \textbf{pH} & \textbf{Exemple camerounais} \\\hline
Acide & $[\ce{H3O+}] > [\ce{HO-}]$ & $pH < 7$ & Jus de bissap ($pH \approx 2,5$), vin de palme frais \\\hline
Neutre & $[\ce{H3O+}] = [\ce{HO-}]$ & $pH = 7$ & Eau distillée, eau de source (environ) \\\hline
Basique & $[\ce{H3O+}] < [\ce{HO-}]$ & $pH > 7$ & Lessive artisanale (cendre), eau de mer \\\hline
\end{tabular}
\end{center}
\end{propriete}

\begin{propriete}[Formules et valeurs de référence (à connaître par cœur)]
\begin{center}
\begin{tabularx}{\linewidth}{|l|X|}\hline
\rowcolor{popblueL} \textbf{Relation} & \textbf{Utilisation / Précision} \\\hline
$K_e=[\ce{H3O+}][\ce{HO-}]$ & Valable dans \textbf{toute} solution aqueuse, à toute $T$ \\\hline
$K_e=\SI{1,0e-14}{}$ ; $pK_e=14$ à \SI{25}{\celsius} & $K_e$ \textbf{augmente} avec $T$ (réaction endothermique) \\\hline
$pH=-\log[\ce{H3O+}]$ & $[\ce{H3O+}]$ en \si{mol.L^{-1}} ; pH sans unité \\\hline
$[\ce{H3O+}]=10^{-pH}$ & Inverse : touche \texttt{10\textsuperscript{x}} ou \texttt{SHIFT log} \\\hline
$[\ce{HO-}]=10^{pH-pK_e}$ & Calcul direct depuis le pH \\\hline
$pH=7$ neutre ; $<7$ acide ; $>7$ basique & \textbf{Uniquement à \SI{25}{\celsius}} \\\hline
$C_iV_i=C_fV_f$ ; $F=\dfrac{C_i}{C_f}=\dfrac{V_f}{V_i}$ & Conservation de la quantité de matière ($n_i=n_f$) \\\hline
\end{tabularx}
\end{center}
\end{propriete}

\begin{methode}[Démarches types : caractère, approximations, mesure, dilution]
\textbf{1. Déterminer le caractère d'une solution}
\begin{enumerate}
  \item Calculer $[\ce{H3O+}]$ (depuis $C$ ou pH) et $[\ce{HO-}] = K_e/[\ce{H3O+}]$.
  \item Comparer $[\ce{H3O+}]$ et $[\ce{HO-}]$ (valable à toute $T$).
  \item \textit{Si $T=\SI{25}{\celsius}$} : comparer pH à 7.
\end{enumerate}

\textbf{2. Faire les approximations (espèces majoritaires / minoritaires)}
\begin{itemize}
  \item L'eau (solvant) est \textbf{toujours ultra-majoritaire} : $[\ce{H2O}] \approx \SI{55,5}{mol.L^{-1}}$.
  \item Solution \textbf{acide} : \ce{H3O+} majoritaire (par rapport à \ce{HO-}), \ce{HO-} \textbf{ultra-minoritaire}.
  \item Solution \textbf{basique} : \ce{HO-} majoritaire, \ce{H3O+} \textbf{ultra-minoritaire}.
  \item Ces approximations justifient l'oubli de $[\ce{HO-}]$ dans le bilan de matière d'une solution acide (et vice-versa).
\end{itemize}

\textbf{3. Mesurer le pH : choisir la méthode adaptée}
\begin{enumerate}
  \item \textbf{pH-mètre} (électrode de verre) : étalonnage 2 points (pH 4,01 ; 7,00 ; 9,21). Précision $\pm \SI{0,05}{}$ à $\pm \SI{0,1}{}$. \emph{Le plus précis.}
  \item \textbf{Papier pH universel} : comparaison visuelle. Précision $\pm \SI{0,5}{}$ à $\pm \SI{1}{}$.
  \item \textbf{Indicateur coloré} : encadrement du pH par la zone de virage. Précision qualitative.
\end{enumerate}

\textbf{4. Préparer une solution fille par dilution (protocole rigoureux)}
\begin{enumerate}
  \item Calculer $V_i = \dfrac{C_f V_f}{C_i}$ (pipette jaugée) ou $V_i = \dfrac{V_f}{F}$.
  \item Rincer la \textbf{pipette jaugée} avec la solution mère (2-3 fois).
  \item Prélever $V_i$ au-dessus du trait de jauge, ajuster au ménisque (œil au niveau), laisser s'écouler (ne pas souffler la goutte restante).
  \item Verser dans la \textbf{fiole jaugée} de volume $V_f$ (rincée à l'eau distillée).
  \item Rincer la pipette (goutte résiduelle) dans la fiole.
  \item Compléter à l'eau distillée \textbf{jusqu'au trait de jauge} (ménisque bas tangent).
  \item Bouchonner, homogénéiser (retournements lents).
\end{enumerate}
\end{methode}

\begin{experience}[Préparation d'une solution diluée : exemple chiffré]
\textbf{Objectif} : Préparer \SI{100,0}{mL} d'une solution d'acide chlorhydrique \SI{0,100}{mol.L^{-1}} à partir d'une solution mère \SI{1,00}{mol.L^{-1}}.
\textbf{Matériel} : Pipette jaugée \SI{10,0}{mL} (classe A), fiole jaugée \SI{100,0}{mL} (classe A), bécher, eau distillée, entonnoir.
\textbf{Calcul} : $V_i = \dfrac{C_f V_f}{C_i} = \dfrac{\SI{0,100}{mol.L^{-1}} \times \SI{100,0}{mL}}{\SI{1,00}{mol.L^{-1}}} = \SI{10,0}{mL}$. Facteur $F=10$.
\textbf{Observations} : La solution mère est incolore, la solution fille aussi. Pas de dégagement de chaleur notable (dilution peu exothermique pour HCl).
\textbf{Interprétation} : $n_{\ce{H3O+}}$ initial = $n_{\ce{H3O+}}$ final = \SI{1,00e-2}{mol}. La concentration est divisée par 10, le pH augmente de 1 unité (pH passe de 1,00 à 2,00).
\end{experience}

\begin{experience}[Utilisation du pH-mètre : étalonnage et mesure]
\textbf{Dispositif} : pH-mètre, électrode combinée (verre/référence), solutions tampons étalons (pH 4,01 ; 7,00 ; 9,21 à \SI{25}{\celsius}), béchers, eau distillée, papier essuie-tout.
\textbf{Protocole}
\begin{enumerate}
  \item Rincer l'électrode à l'eau distillée, tamponner (ne pas frotter le bulbe).
  \item \textbf{Étalonnage 2 points} : plonger dans tampon pH 7,00 $\rightarrow$ valider ; rincer ; plonger dans tampon pH 4,01 (ou 9,21 selon domaine) $\rightarrow$ valider. Vérifier la pente (95-105\%).
  \item Rincer, plonger dans la solution test (agitation magnétique douce), attendre stabilisation ($\Delta pH < 0,01$/min).
  \item Noter pH et température. Rincer électrode, stocker dans solution de conservation (KCl saturé).
\end{enumerate}
\textbf{Précision attendue} : $\pm \SI{0,05}{}$ pH si étalonnage récent et température stable.
\end{experience}

\begin{savaistu}[Indicateurs colorés naturels : le bissap et le chou rouge]
Au Cameroun, deux indicateurs naturels sont faciles à préparer en TP :
\begin{itemize}
  \item \textbf{Jus de bissap (fleurs d'\emph{Hibiscus sabdariffa})} : 
  \begin{itemize}
    \item \textit{Préparation} : Faire bouillir \SI{10}{g} de fleurs séchées dans \SI{200}{mL} d'eau \SI{5}{min}. Filtrer. Le filtrat rouge violacé est l'indicateur.
    \item \textit{Propriétés} : Rouge acide ($pH < 3$), Violet ($pH \approx 5-6$), Bleu/Vert basique ($pH > 8$). Zone de virage utile \textbf{4 -- 8}.
    \item \textit{Chimie} : Anthocyanes (cyanidine-3-sambubioside). Forme acide (rouge) $\ce{AH}$ / Forme basique (bleu) $\ce{A-}$.
  \end{itemize}
  \item \textbf{Jus de chou rouge} : 
  \begin{itemize}
    \item \textit{Préparation} : Mixer \SI{50}{g} de chou rouge avec \SI{100}{mL} d'eau. Filtrer.
    \item \textit{Propriétés} : Rouge ($pH < 3$), Violet ($pH 5-7$), Bleu ($pH 8$), Vert ($pH 10$), Jaune ($pH > 12$). Zone de virage large \textbf{2 -- 12}.
  \end{itemize}
\end{itemize}
\textbf{Intérêt pédagogique} : Visualiser la zone de virage, comprendre qu'un indicateur est un couple acide/base \ce{AH}/\ce{A-} aux formes colorées distinctes. Coût quasi nul, disponible localement.
\end{savaistu}

\begin{aretenir}[Checklist finale avant le Bac -- Leçon 06]
Avant l'épreuve, vérifie que tu sais \textbf{sans hésiter} :
\begin{itemize}
  \item[$\square$] Dessiner la molécule d'eau coudée, indiquer les dipôles de liaison et le dipôle moléculaire, justifier le pouvoir ionisant.
  \item[$\square$] Définir acide/base de Brønsted ; identifier les couples \ce{AH}/\ce{A-} et \ce{BH+}/\ce{B} dans une équation.
  \item[$\square$] Écrire l'équation d'autoprotolyse : \ce{2 H2O <=> H3O+ + HO-} (équilibrée atomes + charges).
  \item[$\square$] Énoncer $K_e=[\ce{H3O+}][\ce{HO-}]=\SI{1,0e-14}{}$ à \SI{25}{\celsius} et $pK_e=14$.
  \item[$\square$] Dire que $K_e$ augmente avec $T$ (neutralité $pH<7$ si $T>\SI{25}{\celsius}$).
  \item[$\square$] Passer de $[\ce{H3O+}]$ au pH et inversement (calculatrice : \texttt{log}, \texttt{10\textsuperscript{x}}).
  \item[$\square$] Calculer $[\ce{HO-}] = 10^{pH-pK_e}$ ou $K_e/[\ce{H3O+}]$.
  \item[$\square$] Conclure sur le caractère acide/neutre/basique (comparaison concentrations ou pH vs 7 à \SI{25}{\celsius}).
  \item[$\square$] Classer les méthodes de mesure par précision : pH-mètre $>$ papier pH $>$ indicateur.
  \item[$\square$] Citer la zone de virage de \textbf{deux} indicateurs usuels (ex. BBT : 6,0--7,6 ; Phénolphtaléine : 8,2--10,0 ; Bissap : 4--8).
  \item[$\square$] Décrire le protocole de préparation d'un indicateur naturel (bissap ou chou rouge).
  \item[$\square$] Calculer $V_i$ pour une dilution ($C_iV_i=C_fV_f$) et décrire les gestes justes (pipette, fiole, ménisque, homogénéisation).
\end{itemize}
\end{aretenir}

\findecours

\end{document}
